
\chapter{Database Management System (DBMS) and SQL}
\label{ch:dbms_sql}

This chapter covers the foundations of database management, transaction processing, data integrity, concurrency control, and database normalization. These topics are fundamental to understanding how enterprise systems store, retrieve, manipulate, and safeguard structured organizational information, and form a core pillar of technical recruitment examinations for IT cadres (Assistant Programmer, Database Administrator, Assistant Network Engineer, and ICT Officer).

The chapter is organized into five major areas:
\begin{enumerate}
\item \textbf{DBMS Fundamentals and Relational Architecture:} Core definitions, data models, three-schema architecture, data independence, and SQL command classification.
\item \textbf{Transaction Processing and ACID Guarantees:} Transaction lifecycles, Atomicity, Consistency, Isolation, Durability, write-ahead logging, and crash recovery.
\item \textbf{Data Integrity and Database Constraints:} Entity integrity, referential integrity, domain integrity, and declarative constraints (NOT NULL, UNIQUE, PRIMARY KEY, FOREIGN KEY, CHECK, DEFAULT).
\item \textbf{Concurrency Control and Transaction Isolation:} Lost updates, dirty reads, non-repeatable reads, phantom reads, ANSI SQL isolation levels, shared/exclusive locks, Two-Phase Locking (2PL), and deadlock resolution.
\item \textbf{Database Normalization and Dependency Theory:} Functional dependencies, normal forms from 1NF through 5NF, lossless join decomposition, dependency preservation, practical conversion workflows, and controlled denormalization.
\end{enumerate}

\newpage
\section{Database and Database Management System (DBMS) Fundamentals}

Information is the central operational asset of modern enterprises. To manage vast volumes of structured records securely and efficiently, organizations transition from primitive file systems to robust database systems.

\subsection{What is a Database?}

A \textbf{database} is an organized, structured collection of logically related data stored electronically in a computer system, designed for rapid search, retrieval, modification, and management.

\textbf{Prominent Real-World Examples:}
\begin{itemize}
\item \textbf{University Student Information System:} Maintains records of enrolled students, academic departments, semester course registrations, and grade point averages (CGPA).
\item \textbf{Core Banking System:} Maintains customer identification profiles, checking/savings accounts, ledger balances, and chronological transaction journals.
\item \textbf{Hospital Management System:} Tracks patient medical histories, clinical appointments, physician rosters, diagnostic lab reports, and pharmaceutical inventories.
\item \textbf{National Citizen Registry (NID / Civil Registry):} Stores citizen identification numbers, biometric records, parental lineages, and voter registration statuses.
\end{itemize}

\subsection{What is a Database Management System (DBMS)?}

A \textbf{Database Management System (DBMS)} is specialized system software that serves as an interface between end-users, software applications, and the physical database files. It provides centralized tools to define, construct, query, update, and administer data while enforcing access controls and data integrity.

\textbf{Industry-Standard Relational DBMS Engines:} PostgreSQL, MySQL, Oracle Database, Microsoft SQL Server, and SQLite.

\subsubsection{Eight Primary Functions of a DBMS}

\begin{enumerate}
\item \textbf{Data Definition:} Provides a Data Definition Language (DDL) allowing database administrators and developers to define schemas, tables, columns, data types, indexes, and constraints.
\item \textbf{Data Manipulation:} Provides a Data Manipulation Language (DML) enabling users to insert new records, retrieve specific subsets of data, update existing attributes, and delete obsolete entries.
\item \textbf{Data Security and Access Control:} Implements authentication, role-based authorization, and encryption mechanisms to restrict data access to authorized personnel (e.g., students may view their own semester grades but cannot modify them).
\item \textbf{Data Integrity Enforcement:} Automatically validates all data modification requests against defined structural constraints (primary keys, foreign keys, check constraints) to ensure accuracy and consistency.
\item \textbf{Transaction Management:} Groups multiple discrete SQL operations into atomic, logical units of work, guaranteeing ACID properties.
\item \textbf{Concurrency Control:} Coordinates simultaneous data access by hundreds or thousands of concurrent users, preventing data corruption and race conditions through locking protocols or Multi-Version Concurrency Control (MVCC).
\item \textbf{Data Backup and Crash Recovery:} Maintains automated transaction logs and checkpoints to restore database state following hardware failure, operating system crashes, or power disruptions.
\item \textbf{Data Independence:} Decouples abstract user queries and application code from physical disk storage structures, access paths, and file formats.
\end{enumerate}

\subsubsection{Advantages and Disadvantages of a DBMS}

\begin{table}[htbp]
\centering
\caption{Architectural Trade-offs: Advantages versus Disadvantages of a DBMS}
\label{tab:dbms_pros_cons}
\begin{tabularx}{\textwidth}{p{7.5cm} X}
\toprule
\textbf{Architectural Advantages} & \textbf{Operational Disadvantages} \\
\midrule
\textbf{Minimized Data Redundancy:} Eliminates duplicate storage across departments via centralized relational schema design. & \textbf{Substantial Licensing and Hardware Cost:} Enterprise database licenses (Oracle, SQL Server) and high-memory servers require heavy capital expenditure. \\
\textbf{Guaranteed Data Consistency:} Updates occur in a single logical location, eliminating conflicting versions of records. & \textbf{Specialized Technical Expertise Required:} Demands highly trained Database Administrators (DBAs) for query tuning, index optimization, and backup administration. \\
\textbf{Multi-User Concurrent Access:} Manages simultaneous user sessions safely without data loss. & \textbf{High Computational Overhead:} Buffer pool caching, write-ahead logging, and locking mechanisms consume significant CPU and RAM resources. \\
\textbf{Standardized Query Language:} Enables complex cross-table joins, aggregations, and subqueries via standard SQL. & \textbf{Single Point of Centralized Failure:} A failure of the primary database cluster can simultaneously halt all dependent organizational business services. \\
\bottomrule
\end{tabularx}
\end{table}

\subsection{Database versus DBMS versus Database Application}

In examination terminology, candidates must clearly differentiate between the conceptual data, the management engine, and the client application:

\begin{table}[htbp]
\centering
\caption{System Distinction: Database, DBMS, Database Application, Server, and DBA}
\label{tab:dbms_layers}
\begin{tabularx}{\textwidth}{p{3.8cm} X X}
\toprule
\textbf{Component} & \textbf{Core Role and Definition} & \textbf{Concrete University Example} \\
\midrule
\textbf{Database} & The organized collection of structured data files stored on persistent disk media. & The actual data records containing student names, enrolled courses, and exam marks. \\
\textbf{DBMS} & The underlying software engine that manages, queries, and secures the database. & \textbf{PostgreSQL} or \textbf{Oracle Database}. \\
\textbf{Database Application} & The frontend software program providing a user interface to interact with the DBMS. & University Student Web Portal / Android Student Registration App. \\
\textbf{Database Server} & The physical computer or virtualized host executing the DBMS daemon process. & Dedicated Linux server running PostgreSQL on port 5432. \\
\textbf{Database Administrator (DBA)} & The technical specialist responsible for architecture, security, performance, and backups. & The DBA managing user privileges, backup schedules, and replication health. \\
\bottomrule
\end{tabularx}
\end{table}

\subsection{DBMS versus Relational DBMS (RDBMS)}

An \textbf{RDBMS (Relational Database Management System)} is a specialized class of DBMS based on the \textbf{Relational Model} introduced by Dr. Edgar F. Codd at IBM in 1970. In an RDBMS, all data is logically represented as two-dimensional relations (tables) composed of rows (tuples) and columns (attributes), and relationships between entities are established through shared key values.

\begin{tcolorbox}[colback=blue!5!white,colframe=blue!60!black,title=\textbf{Crucial Examination Axiom: All RDBMSs are DBMSs, but not all DBMSs are RDBMSs}]
A general DBMS may organize data using hierarchical trees, network graphs, or unstructured key-value stores. An RDBMS strictly adheres to relational algebra, relational calculus, declarative constraints, and Structured Query Language (SQL).
\end{tcolorbox}

\begin{table}[htbp]
\centering
\caption{Comparison of Primary Database Management System Data Models}
\label{tab:dbms_models}
\begin{tabularx}{\textwidth}{p{3.2cm} X p{3.5cm}}
\toprule
\textbf{Database Model} & \textbf{Data Structure and Conceptual Architecture} & \textbf{Prominent Implementations} \\
\midrule
\textbf{Relational (RDBMS)} & Two-dimensional tables with rows and columns, linked via primary/foreign key relationships. Uses SQL. & PostgreSQL, MySQL, Oracle, SQL Server, SQLite. \\
\textbf{Hierarchical} & Tree-like parent-child structure; each child record has exactly one parent record. (1:N only). & IBM IMS (Information Management System). \\
\textbf{Network} & Graph-like record relationships; a child record (member) can possess multiple parent records (owners) (M:N). & IDMS, Integrated Data Store (IDS). \\
\textbf{Document (NoSQL)} & Semi-structured hierarchical documents (JSON, BSON, XML); dynamic, flexible schema. & MongoDB, CouchDB. \\
\textbf{Key-Value (NoSQL)} & Associative array mapping unique string keys to arbitrary data blobs; ultra-fast memory caching. & Redis, Amazon DynamoDB, Memcached. \\
\textbf{Graph (NoSQL)} & Nodes (entities) and edges (relationships) with associated properties; optimized for deep traversals. & Neo4j, Amazon Neptune. \\
\bottomrule
\end{tabularx}
\end{table}

\newpage
\section{Relational Model Concepts: Tables, Domains, Degree, and Cardinality}

The relational data model represents an enterprise domain using formalized mathematical constructs derived from set theory.

\subsection{Fundamental Relational Terminology}

\begin{enumerate}
\item \textbf{Table (Relation):} A two-dimensional grid of data representing a single entity set (e.g., \texttt{Students}, \texttt{Courses}). In formal relational theory, a relation is a subset of the Cartesian product of a list of attribute domains:
\[ R \subseteq \text{dom}(A_1) \times \text{dom}(A_2) \times \dots \times \text{dom}(A_n) \]
\item \textbf{Row (Tuple / Record):} A horizontal entry in a relation representing a single, distinct instance of an entity. In a \texttt{Students} table, the tuple:
\[ \langle 101, \text{'Rahim'}, \text{'CSE'}, 3.75 \rangle \]
records complete information for student Rahim.
\item \textbf{Column (Attribute / Field):} A vertical component of a relation representing a specific named property or characteristic shared by all entities in that relation (e.g., \texttt{Student\_ID}, \texttt{Department}, \texttt{CGPA}).
\item \textbf{Field Value (Atomic Cell):} The scalar data element located at the exact intersection of a specific row and column (e.g., Rahim's CGPA value \texttt{3.75}).
\item \textbf{Domain:} The formal set of all permissible, valid atomic values that can be assigned to a specific attribute. For example:
\begin{itemize}
\item $\text{dom}(\text{CGPA}) = \{ x \in \mathbb{R} \mid 0.00 \le x \le 4.00 \}$
\item $\text{dom}(\text{Age}) = \{ x \in \mathbb{Z} \mid 16 \le x \le 100 \}$
\end{itemize}
\end{enumerate}

\subsection{Degree versus Cardinality: Examination Definitions and Formulas}

Examiners frequently test candidates on the distinction between the horizontal and vertical dimensions of a relation:

\begin{itemize}
\item \textbf{Degree of a Relation:} The total number of \textbf{attributes (columns)} defining the relation schema. The degree is a structural property that changes only when a DDL command alters the table schema (\texttt{ALTER TABLE ADD COLUMN}).
\item \textbf{Cardinality of a Relation:} The total number of \textbf{tuples (rows)} currently stored in the relation instance. Cardinality is a dynamic operational metric that increases upon \texttt{INSERT} and decreases upon \texttt{DELETE}.
\end{itemize}

\begin{table}[htbp]
\centering
\caption{Worked Relational Dimensions Example: Students Table}
\label{tab:degree_cardinality_example}
\begin{tabularx}{\textwidth}{c l c c}
\toprule
\textbf{Student\_ID} & \textbf{Name} & \textbf{Department} & \textbf{CGPA} \\
\midrule
101 & Rahim & CSE & 3.75 \\
102 & Karim & EEE & 3.60 \\
103 & Nabila & CSE & 3.90 \\
104 & Tania & BBA & 3.80 \\
105 & Hasan & ME & 3.45 \\
\bottomrule
\end{tabularx}
\end{table}

For the \texttt{Students} table displayed in Table~\ref{tab:degree_cardinality_example}:
\[ \text{Degree} = \text{Total Columns} = 4 \quad (\texttt{Student\_ID, Name, Department, CGPA}) \]
\[ \text{Cardinality} = \text{Total Rows} = 5 \quad (\text{Five student records}) \]

\begin{tcolorbox}[colback=yellow!10!white,colframe=yellow!60!black,title=\textbf{Exam Memory Mnemonic: Columns Degree, Rows Cardinality}]
Remember: \textbf{D}egree $=$ \textbf{D}imensions across columns (Schema-level). \textbf{C}ardinality $=$ \textbf{C}ount of tuples/records (Instance-level).
\end{tcolorbox}

\newpage
\section{Database Schemas, Instances, and Metadata Catalogs}

A rigorous separation between database design and database contents is maintained throughout relational database theory.

\subsection{Database Schema versus Database Instance}

\begin{itemize}
\item \textbf{Database Schema (Intension):} The overall logical and physical structure, blueprint, and set of rules governing a database. It specifies table names, attribute names, data types, primary and foreign key constraints, check conditions, indexes, and views. The schema is defined during database design and changes very infrequently.
\item \textbf{Database Instance (Extension / State):} The actual snapshot collection of operational data records stored in the database at a specific, given point in time. The instance state changes dynamically with every \texttt{INSERT}, \texttt{UPDATE}, or \texttt{DELETE} transaction.
\end{itemize}

\begin{table}[htbp]
\centering
\caption{Architectural Comparison: Database Schema versus Database Instance}
\label{tab:schema_vs_instance}
\begin{tabularx}{\textwidth}{p{3.2cm} X X}
\toprule
\textbf{Evaluation Metric} & \textbf{Database Schema (Intension)} & \textbf{Database Instance (Extension)} \\
\midrule
\textbf{Conceptual Meaning} & The structural blueprint and integrity definitions. & The operational data populated inside the tables. \\
\textbf{Frequency of Change} & Very low; altered only through planned DDL migrations. & Very high; updates continuously with live transactions. \\
\textbf{Architectural Analogy} & The architectural blueprint and floor plans of a building. & The occupants, furniture, and physical items inside the building. \\
\textbf{SQL Representation} & \texttt{CREATE TABLE}, \texttt{ALTER TABLE}, constraints. & \texttt{INSERT INTO}, \texttt{UPDATE}, \texttt{DELETE}, table rows. \\
\bottomrule
\end{tabularx}
\end{table}

\subsection{Metadata and the System Catalog (Data Dictionary)}

\textbf{Metadata} is formally defined as \textit{``data about data''}. In an RDBMS, metadata is the structured catalog describing the properties, definitions, and locations of all database objects. 

The DBMS maintains this information automatically within internal system tables known as the \textbf{Data Dictionary} or \textbf{System Catalog} (e.g., \texttt{information\_schema} in PostgreSQL/MySQL, or \texttt{DBA\_TABLES} in Oracle). Metadata records include:
\begin{itemize}
\item Table and view definitions, column data types, and character encodings.
\item Constraint definitions (primary keys, foreign keys, check expressions).
\item Physical storage allocations, table spaces, data block allocations, and file extents.
\item Secondary index definitions (B-Tree, Hash, GIN) and index selectivity statistics.
\item User accounts, role assignments, passwords hashes, and object privileges (DCL permissions).
\end{itemize}

\newpage
\section{Relational Database Keys: Super, Candidate, Primary, Alternate, and Foreign Keys}

Keys are foundational relational mechanisms that uniquely identify tuples within relations and enforce referential associations across tables.

\subsection{The Hierarchy of Relational Keys}

\begin{enumerate}
\item \textbf{Super Key (SK):}
A set of one or more attributes that collectively and uniquely identifies every row in a relation. In mathematical terms, a set of attributes $K \subseteq R$ is a super key if no two distinct tuples $t_1, t_2 \in R$ have $t_1[K] = t_2[K]$.
\begin{itemize}
\item \textit{Property:} A super key may contain extraneous, redundant attributes that are unnecessary for unique identification. For example, if \texttt{Student\_ID} is unique, then $\{\texttt{Student\_ID}\}$, $\{\texttt{Student\_ID, Name}\}$, and $\{\texttt{Student\_ID, Name, CGPA}\}$ are all valid super keys.
\end{itemize}

\item \textbf{Candidate Key (CK):}
A \textbf{minimal super key}—a super key from which no attribute can be removed without destroying the property of unique identification.
\begin{itemize}
\item \textit{Property:} If $K$ is a candidate key, then no proper subset $K' \subset K$ is a super key. A relation can have multiple candidate keys.
\item \textit{Example:} If an employee table has unique \texttt{Employee\_ID} and unique national \texttt{NID}, then $\{\texttt{Employee\_ID}\}$ and $\{\texttt{NID}\}$ are both candidate keys.
\end{itemize}

\item \textbf{Primary Key (PK):}
The specific candidate key officially selected by the database designer as the primary, default mechanism for uniquely identifying tuples in a table.
\begin{itemize}
\item \textit{Mandatory Invariants:} (1) Must be strictly unique across all rows; (2) Can \textbf{never contain NULL} (enforces Entity Integrity); (3) Each table can possess \textbf{exactly one} primary key constraint.
\end{itemize}

\item \textbf{Alternate Key (Secondary Key):}
Any candidate key that was \textbf{not} chosen as the primary key. In the employee example, if \texttt{Employee\_ID} is designated as the primary key, then \texttt{NID} becomes the alternate key (typically enforced via a \texttt{UNIQUE NOT NULL} constraint).

\item \textbf{Foreign Key (FK):}
An attribute or set of attributes in a referencing (child) relation that references a candidate key (almost always the primary key) in a referenced (parent) relation.
\begin{itemize}
\item \textit{Property:} Enforces \textbf{Referential Integrity}, ensuring that child rows cannot point to nonexistent parent records. Foreign keys may contain \texttt{NULL} values unless explicitly restricted by a \texttt{NOT NULL} constraint.
\end{itemize}

\item \textbf{Composite Key:}
A key (candidate, primary, or foreign) that is composed of \textbf{two or more attributes} rather than a single attribute.
\begin{itemize}
\item \textit{Example:} In a university \texttt{Enrollment} table, neither \texttt{Student\_ID} nor \texttt{Course\_ID} is unique alone, but the combined pair $(\texttt{Student\_ID, Course\_ID})$ uniquely identifies an enrollment.
\end{itemize}

\item \textbf{Natural Key versus Surrogate Key:}
\begin{itemize}
\item \textbf{Natural Key:} An attribute with inherent business meaning that already exists in the real world and possesses uniqueness (e.g., National ID number, Vehicle VIN, ISBN number).
\item \textbf{Surrogate Key:} An artificial, system-generated numerical identifier (e.g., auto-incrementing integer \texttt{BIGSERIAL}, \texttt{IDENTITY}, or \texttt{UUIDv4}) created solely for database management without real-world business meaning.
\end{itemize}
\end{enumerate}

\begin{figure}[htbp]
\centering
\begin{tcolorbox}[colback=white,colframe=black,title=\textbf{Mathematical Relationship Among Relational Keys}]
\[ \text{Primary Key} \in \text{Candidate Keys} \subseteq \text{Super Keys} \]
Every Primary Key is a Candidate Key. Every Candidate Key is a minimal Super Key. Every Super Key guarantees unique row identification.
\end{tcolorbox}
\end{figure}

\begin{table}[htbp]
\centering
\caption{Comprehensive Comparison Matrix of Relational Database Keys}
\label{tab:relational_keys_comparison}
\begin{tabularx}{\textwidth}{p{2.8cm} p{3.2cm} c c X}
\toprule
\textbf{Key Type} & \textbf{Minimal Uniqueness?} & \textbf{Allows NULL?} & \textbf{Per Table Limit} & \textbf{Primary Operational Purpose} \\
\midrule
\textbf{Super Key} & No (may contain redundant fields) & No (on identifying set) & Unlimited & Theoretical superset of identifying attributes. \\
\textbf{Candidate Key} & \textbf{Yes} (strictly minimal) & No & Multiple permissible & Set of all eligible candidate identifiers. \\
\textbf{Primary Key} & \textbf{Yes} (minimal candidate) & \textbf{Strictly No} & \textbf{Exactly One} & Default table identifier; target of foreign keys. \\
\textbf{Alternate Key} & \textbf{Yes} (minimal candidate) & Depends on config & Multiple permissible & Enforces uniqueness on non-primary candidate keys. \\
\textbf{Foreign Key} & Not necessarily & \textbf{Yes} (unless NOT NULL) & Multiple permissible & Enforces referential relationships between tables. \\
\textbf{Composite Key} & Yes (jointly minimal) & Jointly No & Up to one as PK & Uniquely identifies records across multiple attributes. \\
\bottomrule
\end{tabularx}
\end{table}

\newpage
\section{Entity Relationships and Cardinality Constraints (1:1, 1:N, M:N)}

Relationships define the logical associations connecting entities in a relational database.

\subsection{The Three Cardinality Ratios}

\begin{enumerate}
\item \textbf{One-to-One (1:1) Relationship:}
Each record in Table A is associated with at most one record in Table B, and each record in Table B is associated with at most one record in Table A.
\begin{itemize}
\item \textit{Real-World Example:} A citizen and their official passport document.
\item \textit{Relational Implementation:} Place a Foreign Key in either table referencing the other, and enforce a \textbf{UNIQUE} constraint on that Foreign Key column (or merge into a single shared table).
\end{itemize}

\item \textbf{One-to-Many (1:N) Relationship:}
A single record in Table A can be associated with multiple records in Table B, but each record in Table B is associated with at most one record in Table A.
\begin{itemize}
\item \textit{Real-World Example:} An academic Department employs many Faculty Instructors, but each Instructor belongs to exactly one Department.
\item \textit{Relational Implementation:} Place the Foreign Key on the \textbf{``Many'' side} referencing the Primary Key of the ``One'' side (e.g., place \texttt{Department\_ID} inside the \texttt{Instructors} table).
\end{itemize}

\item \textbf{Many-to-Many (M:N) Relationship:}
Multiple records in Table A can be associated with multiple records in Table B, and vice versa.
\begin{itemize}
\item \textit{Real-World Example:} A Student enrolls in multiple Courses, and a Course contains multiple enrolled Students.
\item \textit{Relational Implementation:} An M:N relationship \textbf{cannot be implemented directly with a single foreign key}. It mandates the creation of a third intermediate table, known as a \textbf{Junction Table (Associative Entity / Bridge Table)}.
\end{itemize}
\end{enumerate}

\begin{table}[htbp]
\centering
\caption{Implementation Strategy for Relational Cardinality Mappings}
\label{tab:relationship_implementation}
\begin{tabularx}{\textwidth}{l p{4.2cm} X}
\toprule
\textbf{Relationship Ratio} & \textbf{Representative Domain Example} & \textbf{Relational Table Schema Implementation} \\
\midrule
\textbf{One-to-One (1:1)} & Person $\longleftrightarrow$ Passport & Place \texttt{Passport\_ID} in \texttt{Persons} with a \texttt{UNIQUE FOREIGN KEY}. \\
\textbf{One-to-Many (1:N)} & Department $\longleftrightarrow$ Students & Place \texttt{Department\_ID} as a \texttt{FOREIGN KEY} inside the \texttt{Students} table. \\
\textbf{Many-to-Many (M:N)} & Students $\longleftrightarrow$ Courses & Create junction table \texttt{Enrollment(Student\_ID, Course\_ID, Grade)} with composite PK. \\
\bottomrule
\end{tabularx}
\end{table}

\newpage
\section{Three-Schema Architecture and Data Independence}

Proposed by the ANSI/SPARC study group in 1975, the \textbf{Three-Schema Architecture} provides a framework separating user perspectives from physical database storage to achieve data independence.

\subsection{The Three Levels of ANSI/SPARC Architecture}

\begin{enumerate}
\item \textbf{External Level (View Level):}
Describes how individual users or distinct application modules perceive the data. It consists of multiple \textbf{External Schemas (User Views)}, tailoring data presentation to specific security clearances and user roles while hiding irrelevant or confidential records (e.g., student view vs. registrar view vs. accountant view).
\item \textbf{Conceptual Level (Logical Level):}
Describes the entire logical structure of the whole database for the entire community of users. It specifies all entity types, attributes, relationships, user constraints, and data types without exposing low-level physical disk layouts.
\item \textbf{Internal Level (Physical Level):}
Describes the low-level physical storage structures, data access paths, file organizations, B-Tree indexes, hashing algorithms, record page allocations, and disk block pointer arrangements.
\end{enumerate}

\subsection{Data Independence: Physical versus Logical}

\textbf{Data Independence} is the capacity to modify the schema at one level of the database architecture without necessitating changes to the schemas at the next higher level.

\begin{table}[htbp]
\centering
\caption{Comparison: Physical Data Independence versus Logical Data Independence}
\label{tab:data_independence_comparison}
\begin{tabularx}{\textwidth}{p{3.2cm} X X}
\toprule
\textbf{Evaluation Dimension} & \textbf{Physical Data Independence} & \textbf{Logical Data Independence} \\
\midrule
\textbf{Level Modified} & \textbf{Internal (Physical) Schema}. & \textbf{Conceptual (Logical) Schema}. \\
\textbf{Level Protected} & \textbf{Conceptual Schema} (and External Views). & \textbf{External Schemas} (User Views \& Applications). \\
\textbf{Concrete Examples} & Reorganizing disk data blocks, switching from B-Tree to Hash index, moving files to an SSD array. & Adding a new column (\texttt{Phone\_Number}) to a table, introducing a new table (\texttt{Scholarships}). \\
\textbf{Implementation Ease} & Relatively straightforward; handled natively by RDBMS storage managers. & Difficult; requires creating and maintaining views to preserve legacy queries. \\
\bottomrule
\end{tabularx}
\end{table}

\newpage
\section{Structured Query Language (SQL) Taxonomy: DDL, DML, DQL, DCL, and TCL}

SQL is the ANSI/ISO standard declarative language for interacting with relational databases. SQL statements are categorized into five distinct functional sub-languages.

\begin{table}[htbp]
\centering
\caption{The Five Functional Command Categories of SQL}
\label{tab:sql_command_categories}
\begin{tabularx}{\textwidth}{l p{3.2cm} X p{4.2cm}}
\toprule
\textbf{Category} & \textbf{Full Language Name} & \textbf{Core Purpose and Function} & \textbf{Primary SQL Commands} \\
\midrule
\textbf{DDL} & Data Definition Language & Defines, modifies, and alters database structures and schemas. & \texttt{CREATE, ALTER, DROP, TRUNCATE, RENAME} \\
\textbf{DML} & Data Manipulation Language & Inserts, updates, and deletes data records stored within tables. & \texttt{INSERT, UPDATE, DELETE} \\
\textbf{DQL} & Data Query Language & Queries and retrieves data subsets from one or more tables. & \texttt{SELECT} \\
\textbf{DCL} & Data Control Language & Administers user permissions, privileges, and database access rights. & \texttt{GRANT, REVOKE} \\
\textbf{TCL} & Transaction Control Language & Manages transactions, commits changes, and establishes savepoints. & \texttt{COMMIT, ROLLBACK, SAVEPOINT} \\
\bottomrule
\end{tabularx}
\end{table}

\begin{tcolorbox}[colback=teal!5!white,colframe=teal!60!black,title=\textbf{Crucial Examination Invariant: DROP versus TRUNCATE versus DELETE}]
A favorite examination topic across all IT recruitment tests is distinguishing these three commands:
\begin{itemize}
\item \textbf{\texttt{DELETE} (DML):} Removes specific rows based on a \texttt{WHERE} clause (or all rows if omitted). Each row deletion is individually written to the transaction log; triggers fire; \textbf{can be rolled back} inside a transaction. Retains table structure.
\item \textbf{\texttt{TRUNCATE} (DDL):} Deallocates the data pages containing the table rows in a single operation. Extremely fast; does not fire row delete triggers; resets auto-increment identity counters. Retains table structure.
\item \textbf{\texttt{DROP} (DDL):} Completely removes the table structure, metadata, indexes, triggers, and all data from the database catalog permanently.
\end{itemize}
\end{tcolorbox}



\newpage
\section{ACID Properties of Database Transactions}

A \textbf{database transaction} is a logical unit of work comprising one or more SQL operations (reads, inserts, updates, deletes) executed against a database such that either all operations succeed and become permanent, or the entire transaction is rolled back with zero side effects.

\subsection{The Banking Transfer Paradigm}

Consider a fund transfer of \textbf{BDT 3,000} from Account A to Account B:
\begin{enumerate}
\item \textbf{Operation 1:} Read balance of Account A and verify $A \ge 3,000$.
\item \textbf{Operation 2:} Deduct BDT 3,000 from Account A ($A \leftarrow A - 3,000$).
\item \textbf{Operation 3:} Read balance of Account B.
\item \textbf{Operation 4:} Add BDT 3,000 to Account B ($B \leftarrow B + 3,000$).
\end{enumerate}
If a power blackout occurs immediately after Operation 2, Account A loses BDT 3,000 while Account B receives nothing—violating total money conservation. To guarantee reliable transaction processing, a DBMS enforces the four \textbf{ACID Properties}.

\subsection{The Four ACID Properties}

\begin{enumerate}
\item \textbf{Atomicity (``All or Nothing''):}
Guarantees that a transaction is treated as a single indivisible unit. Either all statements in the transaction execute and commit successfully, or the entire transaction is aborted and any partial changes are completely undone (rolled back) to the pre-transaction state.
\begin{itemize}
\item \textit{Enforcement Mechanism:} Write-Ahead Logging (WAL) and Undo Logs (Rollback Segments).
\end{itemize}

\item \textbf{Consistency (``Preserving Invariant Rules''):}
Ensures that a transaction transforms the database from one valid state satisfying all defined integrity constraints and business rules to another valid state.
\begin{itemize}
\item \textit{Example:} In the banking transfer, the global invariant is conservation of total balance: $A_{\text{final}} + B_{\text{final}} = A_{\text{initial}} + B_{\text{initial}}$. Furthermore, no account balance may drop below zero ($Balance \ge 0$). If an operation attempts to withdraw BDT 7,000 from an account holding only BDT 5,000, the constraint violation aborts the transaction.
\item \textit{Enforcement Mechanism:} Declarative constraints (Primary Key, Foreign Key, CHECK, NOT NULL) and application validation logic.
\end{itemize}

\item \textbf{Isolation (``Concurrent Independence''):}
Dictates that concurrently executing transactions must not interfere with one another. Each transaction must execute as if it were the sole transaction operating on the system.
\begin{itemize}
\item \textit{Example (The Last Laptop Problem):} An inventory system has \textbf{1 laptop remaining} in stock. Two customers, T1 and T2, attempt to purchase it simultaneously. Both read $\text{Stock} = 1$. Without isolation, both purchase requests succeed, leaving the database with an invalid inventory of $\text{Stock} = -1$. With isolation, the DBMS serializes or locks the record so only one customer obtains the unit while the other is notified that the product is sold out.
\item \textit{Enforcement Mechanism:} Concurrency control protocols (Two-Phase Locking, Timestamp Ordering) and Multi-Version Concurrency Control (MVCC).
\end{itemize}

\item \textbf{Durability (``Permanent Survival''):}
Guarantees that once a transaction successfully commits, its recorded modifications become permanent and will survive any subsequent system failure, power outage, or operating system crash.
\begin{itemize}
\item \textit{Enforcement Mechanism:} Flushing transaction commit log records to non-volatile stable storage (Write-Ahead Logging / Redo Logs) and database checkpoints.
\end{itemize}
\end{enumerate}

\begin{table}[htbp]
\centering
\caption{Master Architectural Summary: The ACID Properties of Database Transactions}
\label{tab:acid_properties_master}
\begin{tabularx}{\textwidth}{l p{4.2cm} X X}
\toprule
\textbf{Property} & \textbf{Core Transactional Rule} & \textbf{Failure Mode Prevented} & \textbf{Primary DBMS Mechanism} \\
\midrule
\textbf{Atomicity} & \textbf{All or nothing}; entire unit commits or rolls back. & Partial transaction execution and split states. & Undo logging, write-ahead logging (WAL), rollback. \\
\textbf{Consistency} & Moves database from one valid state to another valid state. & Violation of database constraints and invariants. & Schema constraints, foreign keys, triggers, business rules. \\
\textbf{Isolation} & Concurrent transactions execute without mutual interference. & Concurrency anomalies (lost updates, dirty reads). & Locking protocols (2PL), MVCC snapshot isolation. \\
\textbf{Durability} & Committed changes persist permanently across crashes. & Loss of committed updates upon server reboot. & Redo logs, WAL flushed to disk, battery-backed write buffers. \\
\bottomrule
\end{tabularx}
\end{table}

\newpage
\section{Database Constraints: NOT NULL, UNIQUE, PRIMARY KEY, FOREIGN KEY, CHECK, and DEFAULT}

\textbf{Database Constraints} are declarative rules enforced by the database engine on table columns to restrict the values that can be inserted, updated, or maintained, preserving data validity.

\subsection{The Six Fundamental Declarative Constraints}

\begin{enumerate}
\item \textbf{\texttt{NOT NULL} Constraint:} Enforces that an attribute cannot contain a \texttt{NULL} (missing/unknown) value. Mandatory attributes (e.g., student name) must be defined with \texttt{NOT NULL}.
\item \textbf{\texttt{UNIQUE} Constraint:} Guarantees that all non-null values in a column (or combination of columns) are strictly distinct across all rows. Unlike primary keys, SQL standard permits multiple columns to possess unique constraints, and most engines allow \texttt{NULL} values in unique columns.
\item \textbf{\texttt{PRIMARY KEY} Constraint:} Uniquely identifies each row in a relation. Implicitly combines \textbf{\texttt{UNIQUE}} and \textbf{\texttt{NOT NULL}}. Each table can possess exactly one primary key constraint, which may consist of one column or a composite set of columns.
\item \textbf{\texttt{FOREIGN KEY} Constraint:} Enforces referential integrity by establishing a link between a column in the child table and a candidate/primary key in the parent table. Ensures that child records cannot reference nonexistent parent entities.
\item \textbf{\texttt{CHECK} Constraint:} Enforces an explicit logical expression that must evaluate to \texttt{TRUE} (or \texttt{UNKNOWN} for nulls) for every inserted or updated row:
\begin{lstlisting}[language=SQL]
CGPA DECIMAL(3,2) CHECK (CGPA >= 0.00 AND CGPA <= 4.00)
\end{lstlisting}
\item \textbf{\texttt{DEFAULT} Constraint:} Supplies a predefined fallback literal value or function result when an \texttt{INSERT} statement omits a value for that column (e.g., \texttt{Status VARCHAR(20) DEFAULT 'Active'}).
\end{enumerate}

\subsection{Foreign Key Referential Integrity Actions}

When a referenced parent row is updated or deleted, the foreign key definition dictates the cascading action applied to referencing child records:

\begin{table}[htbp]
\centering
\caption{Foreign Key Referential Actions on Parent Delete / Update}
\label{tab:foreign_key_actions}
\begin{tabularx}{\textwidth}{l X X}
\toprule
\textbf{Action} & \textbf{Operational Behavior on Parent Row Deletion} & \textbf{Typical Application Scenario} \\
\midrule
\textbf{\texttt{CASCADE}} & Automatically deletes (or updates) all matching child rows in the referencing table. & Deleting an \texttt{Order} automatically deletes all corresponding \texttt{Order\_Items}. \\
\textbf{\texttt{SET NULL}} & Sets the foreign key column in all matching child rows to \texttt{NULL} (requires child FK to be nullable). & Deleting a \texttt{Department} sets the \texttt{Department\_ID} of affected employees to \texttt{NULL}. \\
\textbf{\texttt{SET DEFAULT}} & Sets the foreign key column in matching child rows to its defined default value. & Deleting a manager assigns employees to a designated default general manager ID. \\
\textbf{\texttt{RESTRICT}} & Immediately aborts and rejects the parent deletion if any dependent child rows exist. & Prevents deleting a \texttt{Course} if students are actively enrolled in it. \\
\textbf{\texttt{NO ACTION}} & Similar to \texttt{RESTRICT}; deferred constraint validation checks integrity at transaction end. & Standard SQL default; raises an error if child rows remain orphaned at commit. \\
\bottomrule
\end{tabularx}
\end{table}

\subsection{Unified SQL DDL Example Implementing Complete Constraint Suite}

\begin{lstlisting}[language=SQL,caption={Production SQL DDL Incorporating All Six Constraint Classes}]
-- Parent Relation
CREATE TABLE Departments (
    Department_ID INT PRIMARY KEY,
    Department_Name VARCHAR(100) NOT NULL UNIQUE
);

-- Child Relation
CREATE TABLE Students (
    Student_ID INT PRIMARY KEY,
    Student_Name VARCHAR(100) NOT NULL,
    Email VARCHAR(150) UNIQUE,
    Age INT CHECK (Age >= 16 AND Age <= 65),
    CGPA DECIMAL(3,2) CHECK (CGPA BETWEEN 0.00 AND 4.00),
    Enrollment_Status VARCHAR(20) DEFAULT 'Active',
    Department_ID INT,
    
    CONSTRAINT fk_student_department
        FOREIGN KEY (Department_ID)
        REFERENCES Departments(Department_ID)
        ON DELETE RESTRICT
        ON UPDATE CASCADE
);
\end{lstlisting}

\newpage
\section{Database Integrity Rules: Entity, Referential, Domain, and User-Defined Integrity}

Database integrity guarantees that information remains accurate, valid, and consistent throughout its operational lifecycle.

\begin{table}[htbp]
\centering
\caption{The Four Types of Database Integrity}
\label{tab:database_integrity_types}
\begin{tabularx}{\textwidth}{p{3.2cm} X X}
\toprule
\textbf{Integrity Type} & \textbf{Core Definition and Purpose} & \textbf{Primary Enforcement Mechanism} \\
\midrule
\textbf{Entity Integrity} & Ensures that every row in a table is uniquely identifiable. Mandates that the primary key is unique and \textbf{can never contain a NULL value}. & \texttt{PRIMARY KEY} constraint. \\
\textbf{Referential Integrity} & Ensures that relationships between tables remain consistent. A foreign key value must either match an existing primary key value in the parent table or be \texttt{NULL}. & \texttt{FOREIGN KEY} constraint with \texttt{REFERENCES}. \\
\textbf{Domain Integrity} & Ensures that all data values recorded in an attribute column belong strictly to the permitted, valid data domain. & Data types (\texttt{INT}, \texttt{DATE}), \texttt{CHECK} constraints, \texttt{NOT NULL}. \\
\textbf{User-Defined Integrity} & Custom enterprise-specific business rules that cannot be captured by standard relational constraints. & Database Triggers, Stored Procedures, and application business logic. \\
\bottomrule
\end{tabularx}
\end{table}

\newpage
\section{Transaction Lifecycle States and Transaction Control Language (TCL)}

During execution, a database transaction transitions through several distinct states:

\subsection{The Six Transaction States}

\begin{enumerate}
\item \textbf{Active:} The initial state where the transaction begins execution and performs read/write operations.
\item \textbf{Partially Committed:} The state entered after the final statement of the transaction has executed, but before the commit changes are durably flushed to disk storage.
\item \textbf{Committed:} The state entered after all transaction operations have been durably recorded in transaction logs; changes are now permanent.
\item \textbf{Failed:} The state entered when normal execution can no longer proceed due to hardware failure, software error, or constraint violation.
\item \textbf{Aborted:} The state entered after the transaction has been rolled back and the database has been restored to its pre-transaction state.
\item \textbf{Terminated:} The final state when the transaction leaves the system entirely (either after commit or rollback).
\end{enumerate}

\subsection{Transaction Control Language (TCL) Commands}

\begin{itemize}
\item \textbf{\texttt{BEGIN} / \texttt{START TRANSACTION}:} Marks the beginning of an explicit transaction block.
\item \textbf{\texttt{COMMIT}:} Durably writes all transaction changes to disk, releasing locks and terminating the transaction successfully.
\item \textbf{\texttt{ROLLBACK}:} Aborts the transaction and reverts all uncommitted modifications back to the initial state.
\item \textbf{\texttt{SAVEPOINT}:} Establishes an intermediate bookmark within a long transaction, permitting partial rollback (\texttt{ROLLBACK TO savepoint\_name}) without aborting preceding statements.
\end{itemize}

\newpage
\section{Concurrency Anomalies: Lost Updates, Dirty Reads, Non-Repeatable Reads, and Phantom Reads}

When multiple transactions execute concurrently without adequate isolation, four classical data anomalies manifest:

\begin{enumerate}
\item \textbf{Lost Update Anomaly:}
Occurs when two transactions read the identical initial record, calculate updates concurrently, and one transaction writes its update, completely overwriting the other transaction's update without accounting for it.
\begin{itemize}
\item \textit{Numerical Example:} Balance $= 1,000$. T1 reads $1,000$ and calculates $1,000 + 200 = 1,200$. T2 reads $1,000$ and calculates $1,000 - 100 = 900$. T1 writes $1,200$. T2 subsequently writes $900$. T1's deposit is completely lost! (Correct balance must be \textbf{1,100}).
\end{itemize}

\item \textbf{Dirty Read Anomaly (Reading Uncommitted Data):}
Occurs when Transaction T2 reads a data item modified by Transaction T1 before T1 commits. If T1 subsequently aborts and rolls back, T2 has processed data that never existed in the database.

\item \textbf{Non-Repeatable Read (Fuzzy Read) Anomaly:}
Occurs when Transaction T1 reads a specific row twice during its execution, but Transaction T2 modifies or deletes that row and commits in the interval between the two reads. T1 observes different attribute values for the identical row.

\item \textbf{Phantom Read Anomaly:}
Occurs when Transaction T1 executes a range query (e.g., \texttt{WHERE Salary > 50000}), and Transaction T2 inserts or deletes a new row satisfying the predicate and commits. When T1 re-executes the identical range query, a new ``phantom'' record appears.
\end{enumerate}

\newpage
\section{Transaction Isolation Levels: Read Uncommitted, Read Committed, Repeatable Read, and Serializable}

The ANSI/ISO SQL standard defines four transaction isolation levels representing differing trade-offs between concurrency throughput and anomaly protection.

\begin{table}[htbp]
\centering
\caption{ANSI SQL Transaction Isolation Levels and Concurrency Anomalies}
\label{tab:isolation_levels_matrix}
\begin{tabularx}{\textwidth}{l c c c c}
\toprule
\textbf{Isolation Level} & \textbf{Dirty Read} & \textbf{Non-Repeatable Read} & \textbf{Phantom Read} & \textbf{Lost Update} \\
\midrule
\textbf{Read Uncommitted} & \textcolor{red}{Possible} & \textcolor{red}{Possible} & \textcolor{red}{Possible} & \textcolor{red}{Possible} \\
\textbf{Read Committed} & \textbf{Prevented} & \textcolor{red}{Possible} & \textcolor{red}{Possible} & \textbf{Prevented} \\
\textbf{Repeatable Read} & \textbf{Prevented} & \textbf{Prevented} & \textcolor{red}{Possible}* & \textbf{Prevented} \\
\textbf{Serializable} & \textbf{Prevented} & \textbf{Prevented} & \textbf{Prevented} & \textbf{Prevented} \\
\bottomrule
\end{tabularx}
\end{table}

\textit{*Note on Modern RDBMSs:} In PostgreSQL and MySQL InnoDB, the \texttt{Repeatable Read} level is implemented via Multi-Version Concurrency Control (MVCC) snapshot isolation, which also prevents phantom reads in standard queries.

\newpage
\section{Concurrency Control Protocols: Shared/Exclusive Locks, Two-Phase Locking (2PL), and Deadlocks}

\subsection{Lock Types: Shared versus Exclusive}

\begin{itemize}
\item \textbf{Shared Lock (S-Lock / Read Lock):} Acquired by a transaction before reading an item. Multiple transactions can hold compatible shared locks simultaneously on the same item.
\item \textbf{Exclusive Lock (X-Lock / Write Lock):} Acquired before updating or deleting an item. Only one transaction can hold an exclusive lock, blocking all other read and write requests until released.
\end{itemize}

\subsection{Two-Phase Locking (2PL) Protocol}

The \textbf{Two-Phase Locking (2PL)} protocol guarantees \textbf{conflict serializability} by dividing lock handling into two distinct phases:
\begin{enumerate}
\item \textbf{Growing Phase:} The transaction may acquire locks, but cannot release any lock.
\item \textbf{Shrinking Phase:} The transaction may release locks, but cannot acquire any new lock.
\end{enumerate}
\textbf{Strict 2PL:} Retains all exclusive (X) locks until the transaction commits or aborts, eliminating cascading rollbacks.

\subsection{Deadlocks in Databases}

A \textbf{Deadlock} is a circular waiting condition where two or more transactions are permanently blocked because each holds a lock on a resource that the other transaction requires.
\begin{itemize}
\item \textit{Example:} T1 holds lock on A and requests lock on B. Simultaneously, T2 holds lock on B and requests lock on A.
\item \textit{Resolution:} The DBMS deadlock detector constructs a \textbf{Wait-For Graph (WFG)}. If a cycle is detected, the DBMS aborts and rolls back a designated victim transaction to break the cycle.
\end{itemize}

\newpage
\section{Database Recovery Architecture: Transaction Logs, Write-Ahead Logging (WAL), Checkpoints, and Undo/Redo}

\subsection{Write-Ahead Logging (WAL)}

The fundamental principle of database crash recovery is \textbf{Write-Ahead Logging (WAL)}:
\begin{tcolorbox}[colback=white,colframe=black,title=\textbf{The WAL Invariant Rule}]
\textbf{Before any modified data page is written from memory to permanent database disk files, the corresponding log record describing the change must already be flushed to non-volatile stable storage.}
\end{tcolorbox}
This guarantees that if the server crashes mid-operation, the DBMS can inspect the log to determine whether to \textbf{REDO} (reapply committed updates) or \textbf{UNDO} (revert uncommitted updates).

\subsection{Checkpoints and Recovery Operations}

\begin{itemize}
\item \textbf{Checkpoint:} A periodic synchronization operation where the DBMS flushes all modified buffer pool pages to disk and records a checkpoint marker in the transaction log. During recovery, the DBMS only needs to process log records written after the checkpoint, drastically shortening recovery time.
\item \textbf{UNDO Operation:} Reverses changes made by transactions that were active or uncommitted at the moment of failure.
\item \textbf{REDO Operation:} Reapplies changes made by transactions that committed prior to failure but whose data pages were still sitting dirty in volatile RAM.
\end{itemize}



\newpage
\section{Database Normalization Foundations and Anomalies (Insertion, Update, Deletion)}

\textbf{Database Normalization} is the formal theoretical process of decomposing poorly structured, monolithic tables into smaller, well-structured, cohesive relations to eliminate unnecessary data redundancy and prevent operational anomalies.

\subsection{The Monolithic Design Fallacy}

Consider a poorly designed, un-normalized table storing student enrollment, departmental details, and course instruction within a single table:

\begin{table}[htbp]
\centering
\caption{The Un-Normalized Monolithic Table: Redundancy and Flaws}
\label{tab:unnormalized_table}
\begin{tabularx}{\textwidth}{c l c l c l c}
\toprule
\textbf{Student\_ID} & \textbf{Name} & \textbf{Dept\_ID} & \textbf{Dept\_Name} & \textbf{Course\_ID} & \textbf{Course\_Name} & \textbf{Grade} \\
\midrule
101 & Rahim & 10 & CSE & CSE101 & Programming & A \\
101 & Rahim & 10 & CSE & CSE102 & Database & A- \\
102 & Karim & 20 & EEE & CSE101 & Programming & B+ \\
103 & Nabila & 10 & CSE & CSE101 & Programming & A \\
\bottomrule
\end{tabularx}
\end{table}

This design creates severe data redundancy: student Rahim's name and department are repeated across multiple rows, and course details are duplicated for every enrolled student. This redundancy directly causes three major database anomalies:

\subsection{The Three Classic Database Anomalies}

\begin{enumerate}
\item \textbf{Insertion Anomaly:}
Occurs when certain information cannot be recorded into the database without simultaneously inserting unrelated data, or without introducing artificial \texttt{NULL} values.
\begin{itemize}
\item \textit{Concrete Example:} The university introduces a brand-new course, \texttt{CSE103 (Computer Networks)}. However, because the table's composite primary key is $(\texttt{Student\_ID, Course\_ID})$, the course \textbf{cannot be added to the database} until at least one student enrolls, because primary keys prohibit \texttt{NULL} student IDs!
\end{itemize}

\item \textbf{Update (Modification) Anomaly:}
Occurs when a single real-world fact is stored redundantly in multiple rows, requiring every occurrence to be updated simultaneously.
\begin{itemize}
\item \textit{Concrete Example:} If \texttt{CSE101} is renamed from \textit{``Programming''} to \textit{``Structured Programming''}, all 500 rows where students take this course must be updated. If the update fails mid-way, the database enters an inconsistent state where the identical course possesses two conflicting names.
\end{itemize}

\item \textbf{Deletion Anomaly:}
Occurs when deleting one specific fact unintentionally destroys unrelated information that should have been preserved.
\begin{itemize}
\item \textit{Concrete Example:} Student 102 (Karim) is the sole student enrolled in the \texttt{EEE101} course. If Karim drops out and his record is deleted, the university \textbf{unintentionally deletes all record of the existence of EEE101}!
\end{itemize}
\end{enumerate}

\newpage
\section{Functional Dependencies: Determinants, Trivial, Full, Partial, and Transitive Dependencies}

\textbf{Functional Dependency (FD)} is the core mathematical concept underpinning relational normalization theory.

\subsection{Mathematical Definition and Determinants}

Let $R$ be a relation schema, and let $X$ and $Y$ be arbitrary subsets of attributes of $R$. A \textbf{Functional Dependency}, denoted:
\[ X \rightarrow Y \]
(read as \textit{``$X$ functionally determines $Y$''} or \textit{``$Y$ is functionally dependent on $X$''}) means that for any legal instance $r(R)$, whenever two tuples $t_1, t_2 \in r$ have identical values for $X$, they must have identical values for $Y$:
\[ t_1[X] = t_2[X] \implies t_1[Y] = t_2[Y] \]
The left-hand side ($X$) is termed the \textbf{Determinant}, and the right-hand side ($Y$) is the \textbf{Dependent}.

\subsection{Taxonomy of Functional Dependencies}

\begin{enumerate}
\item \textbf{Trivial Functional Dependency:}
A dependency $X \rightarrow Y$ is trivial if $Y \subseteq X$ (the right-hand side is already a subset of the left-hand side).
\begin{itemize}
\item \textit{Example:} $\{\texttt{Student\_ID, Course\_ID}\} \rightarrow \texttt{Student\_ID}$. Trivial dependencies provide no constraint power.
\end{itemize}

\item \textbf{Non-Trivial Functional Dependency:}
A dependency $X \rightarrow Y$ is non-trivial if $Y \not\subseteq X$ (at least one attribute on the right is not present on the left).
\begin{itemize}
\item \textit{Example:} $\texttt{Student\_ID} \rightarrow \texttt{Student\_Name}$.
\end{itemize}

\item \textbf{Full Functional Dependency:}
In a relation with a composite determinant $X$, an attribute $Y$ is \textbf{fully functionally dependent} on $X$ if $X \rightarrow Y$ holds, but for no proper subset $X' \subset X$ does $X' \rightarrow Y$ hold.
\begin{itemize}
\item \textit{Example:} In \texttt{Enrollment}, $\{\texttt{Student\_ID, Course\_ID}\} \rightarrow \texttt{Grade}$. Grade depends on both student and course; neither attribute alone determines grade.
\end{itemize}

\item \textbf{Partial Functional Dependency:}
Occurs when a non-prime attribute depends on only a \textbf{proper subset} of a composite candidate key.
\begin{itemize}
\item \textit{Example:} In \texttt{Enrollment}, $\{\texttt{Student\_ID, Course\_ID}\} \rightarrow \texttt{Student\_Name}$. Here, $\texttt{Student\_ID} \rightarrow \texttt{Student\_Name}$ alone. Because \texttt{Student\_ID} is a proper subset of the composite primary key, this is a \textbf{partial dependency}.
\end{itemize}

\item \textbf{Transitive Functional Dependency:}
Occurs when an attribute depends on a candidate key indirectly through another non-prime attribute. Formally, if $X \rightarrow Y$ and $Y \rightarrow Z$ hold (where $Y \not\rightarrow X$ and $Z \not\subseteq Y$), then $X \rightarrow Z$ is a \textbf{transitive dependency}:
\[ \texttt{Student\_ID} \rightarrow \texttt{Dept\_ID} \quad \text{and} \quad \texttt{Dept\_ID} \rightarrow \texttt{Dept\_Name} \implies \texttt{Student\_ID} \rightarrow \texttt{Dept\_Name} \]
\end{enumerate}

\begin{table}[htbp]
\centering
\caption{Functional Dependency Taxonomy and Alignment with Normal Forms}
\label{tab:fd_taxonomy_comparison}
\begin{tabularx}{\textwidth}{p{2.8cm} X X}
\toprule
\textbf{Dependency Type} & \textbf{Formal Characteristic} & \textbf{Normal Form Addressing It} \\
\midrule
\textbf{Trivial} & $Y \subseteq X$; holds unconditionally by definition. & Foundational relational property. \\
\textbf{Full Dependency} & Non-prime attribute depends on the entire candidate key. & \textbf{Second Normal Form (2NF)}. \\
\textbf{Partial Dependency} & Non-prime attribute depends on a proper subset of a composite key. & \textbf{Eliminated by 2NF}. \\
\textbf{Transitive Dependency} & Non-prime attribute depends on a key through an intermediate non-key. & \textbf{Eliminated by 3NF}. \\
\textbf{Multivalued (MVD)} & Attribute independently determines multiple values of another ($X \twoheadrightarrow Y$). & \textbf{Eliminated by 4NF}. \\
\textbf{Join Dependency} & Relation is losslessly reconstructed from multiple project-joins. & \textbf{Eliminated by 5NF (PJNF)}. \\
\bottomrule
\end{tabularx}
\end{table}

\newpage
\section{Relational Normal Forms: 1NF, 2NF, 3NF, BCNF, 4NF, and 5NF}

Normalization proceeds in a sequence of progressively stricter levels termed \textbf{Normal Forms}.

\subsection{First Normal Form (1NF)}

\begin{tcolorbox}[colback=white,colframe=black,title=\textbf{1NF Definition}]
A relation is in \textbf{First Normal Form (1NF)} if and only if:
\begin{enumerate}
\item Every attribute contains only \textbf{atomic (indivisible) scalar values} within its domain.
\item The relation contains \textbf{no repeating groups} or multi-valued attributes within a single cell.
\item Each row can be uniquely identified by a primary key.
\end{enumerate}
\end{tcolorbox}
\textbf{Violation Example:} Storing comma-separated values (e.g., \texttt{Courses: "CSE101, CSE102"}).\\
\textbf{Resolution:} Unnest each course into a distinct row with a composite key $(\texttt{Student\_ID, Course\_ID})$.

\subsection{Second Normal Form (2NF)}

\begin{tcolorbox}[colback=white,colframe=black,title=\textbf{2NF Definition}]
A relation is in \textbf{Second Normal Form (2NF)} if and only if:
\begin{enumerate}
\item It satisfies \textbf{First Normal Form (1NF)}.
\item It contains \textbf{no partial functional dependencies}—every non-prime attribute must be fully functionally dependent on every candidate key.
\end{enumerate}
\end{tcolorbox}
\textbf{Crucial Invariant:} If a relation in 1NF has \textbf{only single-attribute candidate keys}, it is \textbf{automatically in 2NF} because partial dependency can only exist when candidate keys are composite!

\subsection{Third Normal Form (3NF)}

\begin{tcolorbox}[colback=white,colframe=black,title=\textbf{3NF Definition (Formal and Practical)}]
A relation is in \textbf{Third Normal Form (3NF)} if and only if:
\begin{enumerate}
\item It satisfies \textbf{Second Normal Form (2NF)}.
\item For every non-trivial functional dependency $X \rightarrow A$, at least one of the following holds:
\begin{itemize}
\item $X$ is a \textbf{Super Key} of the relation; \textbf{OR}
\item $A$ is a \textbf{Prime Attribute} (meaning $A$ is a member of at least one candidate key).
\end{itemize}
\end{enumerate}
\end{tcolorbox}
\textbf{Practical Rule of Thumb:} No non-prime attribute should depend transitively on a candidate key (\textit{``Attributes must depend on the key, the whole key, and nothing but the key''}).

\subsection{Boyce-Codd Normal Form (BCNF)}

Introduced by Raymond F. Boyce and Edgar F. Codd in 1974, \textbf{BCNF} is a stricter version of 3NF that eliminates anomalies occurring when relations possess overlapping composite candidate keys.

\begin{tcolorbox}[colback=white,colframe=black,title=\textbf{BCNF Definition}]
A relation is in \textbf{Boyce-Codd Normal Form (BCNF)} if and only if for every non-trivial functional dependency:
\[ X \rightarrow Y \]
\textbf{$X$ must be a Super Key.}
\end{tcolorbox}

\subsubsection{3NF versus BCNF: The Classic Teaching Example}

Consider a university relation: \texttt{Teaching(Student, Course, Instructor)} under the business rules:
\begin{enumerate}
\item For each course, each student has exactly one instructor: $(\texttt{Student, Course}) \rightarrow \texttt{Instructor}$.
\item Each instructor teaches exactly one course: $\texttt{Instructor} \rightarrow \texttt{Course}$.
\item A course can be taught by multiple instructors.
\end{enumerate}
Candidate keys are: $CK_1 = (\texttt{Student, Course})$ and $CK_2 = (\texttt{Student, Instructor})$.
\begin{itemize}
\item Examine $\texttt{Instructor} \rightarrow \texttt{Course}$: \texttt{Instructor} is \textbf{NOT a super key} (violating BCNF).
\item However, \texttt{Course} is a \textbf{prime attribute} (part of $CK_1$). Thus, this relation satisfies \textbf{3NF}, but \textbf{violates BCNF}!
\end{itemize}

\subsection{Fourth Normal Form (4NF) and Multivalued Dependencies}

A \textbf{Multivalued Dependency (MVD)}, denoted $X \twoheadrightarrow Y$, occurs when $X$ independently determines a set of values for $Y$, regardless of the values of other attributes.
\begin{tcolorbox}[colback=white,colframe=black,title=\textbf{4NF Definition}]
A relation is in \textbf{Fourth Normal Form (4NF)} if and only if it is in BCNF and for every non-trivial multivalued dependency $X \twoheadrightarrow Y$, \textbf{$X$ is a super key}.
\end{tcolorbox}
\textbf{Example:} \texttt{Student\_Activities(Student\_ID, Skill, Hobby)}. A student's skills (Python, Java) and hobbies (Cricket, Reading) are completely independent. Cross-product multiplication creates 4 rows. 4NF decomposes this into \texttt{Student\_Skills(Student\_ID, Skill)} and \texttt{Student\_Hobbies(Student\_ID, Hobby)}.

\subsection{Fifth Normal Form (5NF / Project-Join Normal Form)}

A relation is in \textbf{Fifth Normal Form (5NF)} if and only if every non-trivial \textbf{Join Dependency} in the relation is implied by its candidate keys. 5NF applies when a table can be decomposed into three or more smaller tables, but not into two tables without introducing spurious rows upon reconstruction.

\newpage
\section{Comprehensive Normalization Case Study: 1NF to 3NF/BCNF Step-by-Step with SQL Implementation}

Recruitment examinations routinely require candidates to take an un-normalized schema, identify functional dependencies, and decompose it through 1NF, 2NF, and 3NF.

\subsection{Initial Problem Schema and Assumptions}

\begin{lstlisting}[language=SQL]
Enrollment_Master (
    Student_ID, Student_Name, Department_ID, Department_Name,
    Course_ID, Course_Name, Instructor_ID, Instructor_Name, Grade
)
\end{lstlisting}

\textbf{Business Rules / Invariants:}
\begin{enumerate}
\item Each student has a unique \texttt{Student\_ID} and belongs to exactly one department.
\item Each department has a unique \texttt{Department\_ID} and a unique \texttt{Department\_Name}.
\item Each course has a unique \texttt{Course\_ID} and has one designated primary instructor.
\item Each instructor has a unique \texttt{Instructor\_ID} and name.
\item A student receives a single letter \texttt{Grade} for each completed course enrollment.
\end{enumerate}

\subsection{Step 1: Identifying Functional Dependencies and Primary Key}

Based on the domain business rules, the set of functional dependencies $F$ is:
\begin{align*}
FD_1: &\quad \texttt{Student\_ID} \rightarrow \texttt{Student\_Name, Department\_ID} \\
FD_2: &\quad \texttt{Department\_ID} \rightarrow \texttt{Dept\_Name} \\
FD_3: &\quad \texttt{Course\_ID} \rightarrow \texttt{Course\_Name, Instructor\_ID} \\
FD_4: &\quad \texttt{Instructor\_ID} \rightarrow \texttt{Instructor\_Name} \\
FD_5: &\quad \{\texttt{Student\_ID, Course\_ID}\} \rightarrow \texttt{Grade}
\end{align*}
The minimal Candidate Key for the un-decomposed table is:
\[ \boxed{\text{Candidate Key} = \{\texttt{Student\_ID, Course\_ID}\}} \]

\subsection{Step 2: Conversion to 1NF}
Every cell contains an atomic scalar value, and repeating groups are unnested. The relation is in 1NF.

\subsection{Step 3: Conversion to 2NF (Eliminating Partial Dependencies)}
In the 1NF relation, non-prime attributes \texttt{Student\_Name} and \texttt{Department\_ID} depend only on \texttt{Student\_ID} ($FD_1$), while \texttt{Course\_Name} and \texttt{Instructor\_ID} depend only on \texttt{Course\_ID} ($FD_3$). These are \textbf{partial dependencies}.
Decompose into three relations:
\begin{enumerate}
\item \texttt{Students\_2NF(Student\_ID, Student\_Name, Department\_ID, Department\_Name)} (PK: \texttt{Student\_ID}).
\item \texttt{Courses\_2NF(Course\_ID, Course\_Name, Instructor\_ID, Instructor\_Name)} (PK: \texttt{Course\_ID}).
\item \texttt{Enrollment\_2NF(Student\_ID, Course\_ID, Grade)} (PK: $\{\texttt{Student\_ID, Course\_ID}\}$).
\end{enumerate}

\subsection{Step 4: Conversion to 3NF (Eliminating Transitive Dependencies)}
In \texttt{Students\_2NF}, $\texttt{Department\_ID} \rightarrow \texttt{Department\_Name}$ is a transitive dependency through non-key attribute \texttt{Department\_ID}.
In \texttt{Courses\_2NF}, $\texttt{Instructor\_ID} \rightarrow \texttt{Instructor\_Name}$ is a transitive dependency through non-key attribute \texttt{Instructor\_ID}.
Decompose into five final relations:
\begin{enumerate}
\item \textbf{\texttt{Departments:}} $(\underline{\texttt{Department\_ID}}, \texttt{Department\_Name})$
\item \textbf{\texttt{Students:}} $(\underline{\texttt{Student\_ID}}, \texttt{Student\_Name}, \texttt{Department\_ID})$
\item \textbf{\texttt{Instructors:}} $(\underline{\texttt{Instructor\_ID}}, \texttt{Instructor\_Name})$
\item \textbf{\texttt{Courses:}} $(\underline{\texttt{Course\_ID}}, \texttt{Course\_Name}, \texttt{Instructor\_ID})$
\item \textbf{\texttt{Enrollment:}} $(\underline{\texttt{Student\_ID}}, \underline{\texttt{Course\_ID}}, \texttt{Grade})$
\end{enumerate}

\subsection{Final Production SQL DDL Implementation}

\begin{lstlisting}[language=SQL,caption={Normalized 3NF/BCNF Schema with Relational Integrity}]
-- 1. Departments Table
CREATE TABLE Departments (
    Department_ID INT PRIMARY KEY,
    Department_Name VARCHAR(100) NOT NULL UNIQUE
);

-- 2. Students Table
CREATE TABLE Students (
    Student_ID INT PRIMARY KEY,
    Student_Name VARCHAR(100) NOT NULL,
    Department_ID INT NOT NULL,
    FOREIGN KEY (Department_ID) REFERENCES Departments(Department_ID)
);

-- 3. Instructors Table
CREATE TABLE Instructors (
    Instructor_ID VARCHAR(10) PRIMARY KEY,
    Instructor_Name VARCHAR(100) NOT NULL
);

-- 4. Courses Table
CREATE TABLE Courses (
    Course_ID VARCHAR(20) PRIMARY KEY,
    Course_Name VARCHAR(100) NOT NULL,
    Instructor_ID VARCHAR(10) NOT NULL,
    FOREIGN KEY (Instructor_ID) REFERENCES Instructors(Instructor_ID)
);

-- 5. Enrollment Junction Table
CREATE TABLE Enrollment (
    Student_ID INT,
    Course_ID VARCHAR(20),
    Grade VARCHAR(3),
    PRIMARY KEY (Student_ID, Course_ID),
    FOREIGN KEY (Student_ID) REFERENCES Students(Student_ID) ON DELETE CASCADE,
    FOREIGN KEY (Course_ID) REFERENCES Courses(Course_ID) ON DELETE RESTRICT
);
\end{lstlisting}

\newpage
\section{Lossless Join Decomposition, Dependency Preservation, and Controlled Denormalization}

\subsection{Lossless Join Decomposition}

A decomposition of relation $R$ into $R_1$ and $R_2$ is \textbf{lossless-join} if and only if computing the natural join of the decomposed relations exactly reproduces the original relation without losing information or creating false, spurious tuples:
\[ R = R_1 \bowtie R_2 \]
\textbf{The Binary Lossless Test:} A binary decomposition of $R$ into $R_1$ and $R_2$ is lossless with respect to a set of functional dependencies $F$ if and only if the common attributes $(R_1 \cap R_2)$ functionally determine at least one of the decomposed relations:
\[ (R_1 \cap R_2) \rightarrow R_1 \quad \text{OR} \quad (R_1 \cap R_2) \rightarrow R_2 \]

\subsection{Dependency Preservation}

A decomposition is \textbf{dependency-preserving} if all functional dependencies in the original set $F$ can be verified by inspecting individual decomposed relations independently, without executing expensive multi-table joins.
\begin{itemize}
\item \textit{Important Theorem:} 3NF synthesis algorithms always guarantee both \textbf{lossless join} and \textbf{dependency preservation}. BCNF decomposition guarantees lossless join, but may \textbf{not} always preserve all functional dependencies.
\end{itemize}

\subsection{Controlled Denormalization}

\textbf{Denormalization} is the deliberate, strategic introduction of controlled data redundancy into a normalized database schema to reduce costly SQL joins and accelerate high-frequency read queries in read-intensive applications (e.g., data warehouses, reporting dashboards).

\begin{table}[htbp]
\centering
\caption{Architectural Comparison: Normalization versus Controlled Denormalization}
\label{tab:norm_vs_denorm}
\begin{tabularx}{\textwidth}{p{3.2cm} X X}
\toprule
\textbf{Evaluation Dimension} & \textbf{Normalized Database (3NF / BCNF)} & \textbf{Denormalized Database} \\
\midrule
\textbf{Primary Objective} & Minimize redundancy and eliminate write anomalies. & Optimize read query throughput and minimize joins. \\
\textbf{Data Redundancy} & Strictly minimized. & Intentionally reintroduced in controlled areas. \\
\textbf{Write Performance} & Fast \texttt{INSERT}, \texttt{UPDATE}, and \texttt{DELETE}. & Slower writes; requires updating multiple duplicate rows. \\
\textbf{Read Performance} & Requires complex multi-table \texttt{JOIN} queries. & Ultra-fast single-table scans. \\
\textbf{Storage Footprint} & Minimal disk and buffer pool storage. & Larger storage requirement due to repeated fields. \\
\textbf{Typical Use Case} & Online Transaction Processing (\textbf{OLTP}) systems. & Online Analytical Processing (\textbf{OLAP}) and reporting. \\
\bottomrule
\end{tabularx}
\end{table}



% =============================================================================
% SECTION 21: SQL DATA TYPES, STORAGE ARCHITECTURE & NUMERIC PRECISION
% =============================================================================
\section{SQL Data Types, Storage Architecture, and Numeric Precision}
\label{sec:dbms_sql_data_types}

\subsection{The Role of Data Types in Relational Database Management}
An \textbf{SQL data type} is an attribute-level metadata declaration that specifies the domain of permissible values, the internal physical storage allocation, the collation order, and the mathematical or logical operations applicable to that column. Defining rigorous, semantically precise data types is a primary engineering responsibility in relational database design, directly impacting:
\begin{itemize}
    \item \textbf{Domain Integrity and Validity:} Prevents illegal states (e.g., storing alphabetic characters in a currency field or calendar dates with day 32).
    \item \textbf{Physical Storage Efficiency:} Minimizes disk and RAM footprint per page, increasing buffer cache hit ratios and reducing I/O bandwidth.
    \item \textbf{Query Execution Performance:} Aligns data fields to native CPU word boundaries, enabling SIMD vectorization and hardware-accelerated comparisons.
    \item \textbf{Arithmetic and Financial Accuracy:} Eliminates rounding, truncation, and precision drift inherent in approximate floating-point representations.
\end{itemize}

\subsection{Numeric Data Types: Fixed-Point versus Floating-Point}
SQL divides numeric data types into two distinct categories: \textbf{exact numeric types} and \textbf{approximate numeric types}.

\begin{table}[htbp]
\centering
\small
\caption{SQL Numeric Data Types and Storage Profiles}
\label{tab:sql_numeric_types}
\begin{tabularx}{\textwidth}{l p{3.2cm} c X}
\toprule
\textbf{Data Type} & \textbf{Internal Representation} & \textbf{Storage} & \textbf{Typical Engineering Use Case} \\
\midrule
\texttt{SMALLINT} & Signed 16-bit integer & 2 bytes & Low-cardinality counters, calendar years ($-32,768$ to $32,767$). \\
\texttt{INT} / \texttt{INTEGER} & Signed 32-bit integer & 4 bytes & Standard primary keys, student IDs ($-2.14 \times 10^9$ to $+2.14 \times 10^9$). \\
\texttt{BIGINT} & Signed 64-bit integer & 8 bytes & High-frequency transaction IDs, core banking ledger serials. \\
\texttt{DECIMAL(p, s)} & Exact packed decimal & Variable & Monetary balances, accounting, tax rates, currency exchange. \\
\texttt{NUMERIC(p, s)} & Exact packed decimal & Variable & Fully synonymous with \texttt{DECIMAL} in ANSI SQL. \\
\texttt{FLOAT(p)} & Approximate binary float & 4 or 8 bytes & Scientific simulations, geospatial coordinate readings. \\
\texttt{REAL} & Single-precision float & 4 bytes & Graphic rendering coords, statistical sensor telemetry. \\
\texttt{DOUBLE PRECISION} & Double-precision float & 8 bytes & High-range scientific computations (IEEE 754 floating point). \\
\bottomrule
\end{tabularx}
\end{table}

\begin{quote}
\textbf{Critical Examination Distinction: DECIMAL versus FLOAT in Banking Systems}
\begin{itemize}
    \item \textbf{\texttt{DECIMAL(p, s)} / \texttt{NUMERIC(p, s)}}: Stores exact decimal fractions by representing the significand as a base-10 packed integer. The parameter $p$ specifies the \textbf{precision} (total count of significant decimal digits, both left and right of the decimal point), while $s$ specifies the \textbf{scale} (number of digits to the right of the decimal point). For example, \texttt{DECIMAL(12, 2)} represents numbers up to $9,999,999,999.99$. In banking, financial balance calculations require exact decimal fractions; rounding errors of even $\$0.000001$ compound over millions of transactions into irreconcilable ledger discrepancies.
    \item \textbf{\texttt{FLOAT} / \texttt{DOUBLE PRECISION}}: Implements IEEE 754 binary floating-point representation ($V = (-1)^s \times 1.m \times 2^{e - \text{bias}}$). Because fractions such as $0.1$ or $0.01$ cannot be represented as finite binary fractions (they result in repeating binary sequences: $0.0001100110011\dots_2$), standard binary arithmetic incurs rounding errors. \textbf{Rule: Never store monetary amounts, financial ledgers, or account balances in \texttt{FLOAT} or \texttt{DOUBLE}.}
\end{itemize}
\end{quote}

\subsection{Character and String Data Types: CHAR versus VARCHAR}
\begin{table}[htbp]
\centering
\small
\caption{Character and String Data Types}
\label{tab:sql_string_types}
\begin{tabularx}{\textwidth}{l p{3.5cm} X}
\toprule
\textbf{Data Type} & \textbf{Storage Characteristics} & \textbf{Recommended Application} \\
\midrule
\texttt{CHAR(n)} & Fixed length of $n$ characters; padded with trailing spaces to $n$ bytes. & Uniform-length codes: ISO-3166 alpha-2 country codes (\texttt{'BD'}, \texttt{'US'}), MD5/SHA-256 hashes, currency codes (\texttt{'BDT'}, \texttt{'USD'}). \\
\texttt{VARCHAR(n)} & Variable length up to $n$ characters; stored with a 1- or 2-byte length prefix. & Variable-length strings: Customer names, email addresses, street addresses, remarks. \\
\texttt{TEXT} & Variable length with unbounded or multi-gigabyte limits (engine dependent). & Long-form narrative content, audit logs, JSON payloads, article bodies. \\
\texttt{NCHAR(n)} & Fixed-length Unicode UTF-16/UCS-2 national character string. & Fixed-length multilingual text in DBMSs requiring explicit national types (e.g., MSSQL). \\
\texttt{NVARCHAR(n)} & Variable-length Unicode UTF-16 national character string. & Variable multilingual strings containing international scripts. \\
\bottomrule
\end{tabularx}
\end{table}

\textbf{Storage Trade-Off Analysis:}
\begin{itemize}
    \item \texttt{CHAR(n)} allocates exactly $n$ character units regardless of actual payload length. If \texttt{'CSE'} is stored in a \texttt{CHAR(10)} column, the DBMS pads it with 7 trailing spaces. It offers faster random memory access in fixed-length records because tuple offsets are predictable constants.
    \item \texttt{VARCHAR(n)} allocates only the actual characters plus length overhead (typically 1 byte for lengths $< 256$, 2 bytes for lengths $\ge 256$). When updated to a longer string, however, \texttt{VARCHAR} can cause \textbf{row migration} and page fragmentation if the expanded row no longer fits within the original data page.
\end{itemize}

\subsection{Temporal Data Types: Date, Time, and Timestamps}
\begin{itemize}
    \item \textbf{\texttt{DATE}:} Encodes year, month, and day (\texttt{YYYY-MM-DD}). Requires 3 or 4 bytes. Example: \texttt{'2026-11-19'}.
    \item \textbf{\texttt{TIME}:} Encodes hour, minute, second, and fractional seconds (\texttt{HH:MM:SS.ffffff}). Example: \texttt{'14:30:00'}.
    \item \textbf{\texttt{TIMESTAMP}:} Combined date and time without time zone offset (\texttt{YYYY-MM-DD HH:MM:SS}).
    \item \textbf{\texttt{TIMESTAMP WITH TIME ZONE} (\texttt{TIMESTAMPTZ}):} Records timestamp along with UTC offset or normalizes stored time to UTC. Essential for global financial networks to preserve chronologically absolute transaction sequences across geographic timezones.
    \item \textbf{\texttt{INTERVAL}:} Represents a duration of time (e.g., \texttt{INTERVAL '3 DAYS'}, \texttt{INTERVAL '2 HOURS'}).
\end{itemize}

\subsection{Boolean, Binary, and Semi-Structured Types}
\begin{itemize}
    \item \textbf{\texttt{BOOLEAN}:} Logical values \texttt{TRUE}, \texttt{FALSE}, and \texttt{NULL} (Three-Valued Logic). Implemented natively in PostgreSQL; mapped to \texttt{TINYINT(1)} in MySQL and \texttt{NUMBER(1)} in Oracle.
    \item \textbf{\texttt{BLOB} / \texttt{BYTEA} / \texttt{VARBINARY}:} Raw uninterpreted byte streams used for cryptographic keys, signatures, or embedded certificates.
    \item \textbf{Object Storage Architecture vs BLOBs:} Best practice in modern software engineering dictates storing large media files (images, PDF statements, audio) in external distributed object stores (e.g., Amazon S3, MinIO) while holding the immutable URL/URI pointer and metadata hash in the relational database.
\end{itemize}

% =============================================================================
% SECTION 22: SQL DATA QUERY LANGUAGE (DQL) & QUERY EXECUTION ORDER
% =============================================================================
\section{SQL Data Query Language (DQL): Syntax, Filtering, and Execution Order}
\label{sec:dbms_sql_dql}

\subsection{The Declarative Nature of SQL}
Structured Query Language is fundamentally \textbf{declarative}: the engineer specifies \textit{what} relational dataset must be retrieved, not \textit{how} the physical execution engine must navigate disk blocks, scan B-trees, or construct hash tables. The DBMS Cost-Based Optimizer (CBO) compiles the declarative SQL query into an algebraic operator tree, estimates cardinalities using statistical histograms, and executes the optimal physical plan.

\subsection{Standard DQL Syntax and Query Clauses}
The canonical syntax of an SQL Data Query statement spans six primary clauses:
\begin{lstlisting}[language=SQL]
SELECT [DISTINCT] column_expressions
FROM source_table_or_join
[WHERE row_level_predicates]
[GROUP BY grouping_columns]
[HAVING group_level_predicates]
[ORDER BY sort_expressions [ASC | DESC]]
[LIMIT row_count OFFSET skip_count];
\end{lstlisting}

\subsection{Relational Projection and DISTINCT Tuple Elimination}
\begin{itemize}
    \item \textbf{Explicit Projection versus \texttt{SELECT *}:} In production applications, always explicitly name required columns. Using \texttt{SELECT *} introduces network transmission overhead, prevents the query planner from executing \textbf{Index-Only Scans} (covering indexes), and creates runtime fragility when table schemas evolve.
    \item \textbf{\texttt{DISTINCT}:} Eliminates duplicate tuples from the output relation. The deduplication operator evaluates the \textbf{entire combination} of projected columns:
\begin{lstlisting}[language=SQL]
-- Evaluates the distinct composite tuple (Department, CGPA)
SELECT DISTINCT Department, CGPA FROM Students;
\end{lstlisting}
\end{itemize}

\subsection{Predicate Evaluation in the WHERE Clause}
The \texttt{WHERE} clause restricts rows retrieved by evaluating a boolean predicate on each individual tuple before grouping.
\begin{table}[htbp]
\centering
\small
\caption{SQL Relational and Logical Predicate Operators}
\label{tab:sql_where_operators}
\begin{tabularx}{\textwidth}{l p{4.5cm} X}
\toprule
\textbf{Operator} & \textbf{Relational Semantics} & \textbf{Example Code} \\
\midrule
\texttt{=}, \texttt{<>} or \texttt{!=} & Equality and inequality tests & \texttt{Department <> 'CSE'} \\
\texttt{>}, \texttt{<}, \texttt{>=}, \texttt{<=} & Numerical/lexicographical order & \texttt{Salary >= 50000.00} \\
\texttt{BETWEEN a AND b} & Inclusive closed interval $[a, b]$ & \texttt{CGPA BETWEEN 3.50 AND 4.00} \\
\texttt{IN (v1, v2, ...)} & Set membership test ($\exists v \in S$) & \texttt{Department IN ('CSE', 'EEE', 'BBA')} \\
\texttt{LIKE 'pattern'} & Wildcard string matching & \texttt{Name LIKE 'R\%'} or \texttt{Code LIKE 'CS\_\_'} \\
\texttt{IS NULL} / \texttt{IS NOT NULL} & Three-valued logic null tests & \texttt{Email IS NULL} \\
\texttt{AND}, \texttt{OR}, \texttt{NOT} & Boolean logical combinators & \texttt{Age < 25 AND CGPA > 3.75} \\
\bottomrule
\end{tabularx}
\end{table}

\begin{quote}
\textbf{The Three-Valued Logic (3VL) Null Trap}
In SQL, \texttt{NULL} does not denote the number zero or an empty string; it denotes the \textbf{absence of a value} (unknown or inapplicable). Consequently, evaluating any direct equality comparison with \texttt{NULL} (e.g., \texttt{WHERE Email = NULL}) does not return \texttt{TRUE} or \texttt{FALSE}---it evaluates to \textbf{\texttt{UNKNOWN}}. Since the \texttt{WHERE} filter discards tuples that evaluate to anything other than \texttt{TRUE}, \texttt{WHERE Email = NULL} silently drops all records. Always use \texttt{IS NULL} or \texttt{IS NOT NULL}.
\end{quote}

\subsection{Pattern Matching with Wildcard Mechanics}
The \texttt{LIKE} operator utilizes two standard ANSI wildcards:
\begin{itemize}
    \item \texttt{\%} matches zero, one, or multiple arbitrary characters.
    \item \texttt{\_} matches exactly one single arbitrary character.
\end{itemize}
\begin{lstlisting}[language=SQL]
SELECT * FROM Students WHERE Name LIKE 'R%';     -- Begins with 'R'
SELECT * FROM Students WHERE Name LIKE '%a';     -- Ends with 'a'
SELECT * FROM Students WHERE Name LIKE '_____';  -- Exactly five characters
SELECT * FROM Students WHERE Code LIKE '_S%';    -- 'S' as second character
\end{lstlisting}

\subsection{Logical Query Processing Order}
Although written beginning with the \texttt{SELECT} keyword, relational engines execute clauses in a strictly defined logical pipeline. Understanding this sequence explains standard compiler constraints:

\begin{figure}[htbp]
\centering
\begin{tcolorbox}[colback=codebg, colframe=primary, title=\textbf{The 8-Stage Logical Query Processing Sequence}]
\begin{enumerate}
    \item \textbf{\texttt{FROM} and \texttt{JOIN}:} Cross-join source relations and apply \texttt{ON} join filters.
    \item \textbf{\texttt{WHERE}:} Evaluate row-level filtering predicates on candidate tuples.
    \item \textbf{\texttt{GROUP BY}:} Partition remaining tuples into equivalence classes (groups).
    \item \textbf{\texttt{HAVING}:} Evaluate group-level predicates on aggregated groups.
    \item \textbf{\texttt{SELECT}:} Evaluate projection expressions, mathematical formulas, and aliases.
    \item \textbf{\texttt{DISTINCT}:} Deduplicate projected output rows.
    \item \textbf{\texttt{ORDER BY}:} Sort the final relational result set.
    \item \textbf{\texttt{LIMIT} / \texttt{OFFSET}:} Slice and paginate the final sorted tuple stream.
\end{enumerate}
\end{tcolorbox}
\end{figure}

\textbf{Core Examination Takeaway:} Because \texttt{WHERE} is executed at Step 2 while \texttt{SELECT} column aliases are computed at Step 5, you \textbf{cannot} reference a \texttt{SELECT} alias in the \texttt{WHERE} clause:
\begin{lstlisting}[language=SQL]
-- INVALID: Salary_Bonus alias is evaluated AFTER the WHERE clause
SELECT Salary * 0.10 AS Salary_Bonus FROM Employees WHERE Salary_Bonus > 5000;

-- VALID: Must repeat the expression or wrap inside a Subquery / CTE
SELECT Salary * 0.10 AS Salary_Bonus FROM Employees WHERE (Salary * 0.10) > 5000;
\end{lstlisting}

% =============================================================================
% SECTION 23: ADVANCED SQL SUBQUERIES & PREDICATES
% =============================================================================
\section{Advanced SQL Subqueries, Correlated Execution, and Set Predicates}
\label{sec:dbms_sql_subqueries}

\subsection{Subquery Taxonomy}
A \textbf{subquery} (nested query) is an SQL query contained inside an outer query. Subqueries are categorized according to their return cardinality and structural correlation:

\begin{table}[htbp]
\centering
\small
\caption{Taxonomy of SQL Subqueries}
\label{tab:sql_subquery_taxonomy}
\begin{tabularx}{\textwidth}{l p{3.2cm} p{4.2cm} X}
\toprule
\textbf{Subquery Type} & \textbf{Return Dimension} & \textbf{Allowed Operators} & \textbf{Typical Practical Application} \\
\midrule
\textbf{Scalar} & $1 \times 1$ (1 row, 1 col) & \texttt{=}, \texttt{<>}, \texttt{>}, \texttt{<}, \texttt{>=}, \texttt{<=} & Comparing an attribute against an aggregated metric (e.g., \texttt{> AVG(Salary)}). \\
\textbf{Single-Row} & $1 \times k$ (1 row, $k$ cols) & Relational tuple comparison & Finding records matching an exact multi-attribute tuple. \\
\textbf{Multi-Row} & $m \times 1$ ($m$ rows, 1 col) & \texttt{IN}, \texttt{NOT IN}, \texttt{ANY}, \texttt{ALL} & Filtering by set membership against dynamically generated keys. \\
\textbf{Correlated} & Evaluated per outer row & \texttt{EXISTS}, \texttt{NOT EXISTS}, Relational & Cross-referencing outer tuple attributes (e.g., department-level maximums). \\
\bottomrule
\end{tabularx}
\end{table}

\subsection{Correlated Subqueries versus Uncorrelated Subqueries}
\begin{itemize}
    \item \textbf{Uncorrelated Subquery:} The nested query executes exactly \textbf{once} independently of the outer query. Its result is cached and consumed by the outer query:
\begin{lstlisting}[language=SQL]
-- Subquery executes once to compute company-wide average
SELECT Name, Salary FROM Employees 
WHERE Salary > (SELECT AVG(Salary) FROM Employees);
\end{lstlisting}
    \item \textbf{Correlated Subquery:} The inner query references an attribute from the current candidate tuple of the outer query (\texttt{e1.Department\_ID}). Logically, the inner query must re-evaluate for \textbf{every candidate row} in the outer table:
\begin{lstlisting}[language=SQL]
-- Find employees who earn more than the average salary of THEIR OWN department
SELECT e1.Employee_ID, e1.Name, e1.Salary, e1.Department_ID
FROM Employees e1
WHERE e1.Salary > (
    SELECT AVG(e2.Salary) 
    FROM Employees e2 
    WHERE e2.Department_ID = e1.Department_ID
);
\end{lstlisting}
\end{itemize}

\subsection{Set Existence Predicates: EXISTS versus IN}
\begin{itemize}
    \item \textbf{\texttt{EXISTS}:} Returns \texttt{TRUE} as soon as the inner query returns at least one tuple (Boolean short-circuit evaluation). The projection list in \texttt{SELECT 1} is ignored by the engine; only tuple existence is verified.
\begin{lstlisting}[language=SQL]
-- Retrieve students enrolled in at least one course
SELECT s.Student_ID, s.Name FROM Students s
WHERE EXISTS (
    SELECT 1 FROM Enrollment e WHERE e.Student_ID = s.Student_ID
);
\end{lstlisting}
    \item \textbf{\texttt{NOT IN} versus \texttt{NOT EXISTS} (The Critical NULL Hazard):} If the subquery in a \texttt{NOT IN} predicate returns even a \textbf{single \texttt{NULL}} value, the entire \texttt{NOT IN} expression evaluates to \texttt{UNKNOWN}, resulting in \textbf{zero rows returned}. In contrast, \texttt{NOT EXISTS} safely evaluates boolean existence without being corrupted by nulls.
\end{itemize}

\subsection{Quantified Set Predicates: ANY and ALL}
\begin{itemize}
    \item \textbf{\texttt{v > ANY (Subquery)}:} True if $v$ is greater than \textbf{at least one} element in the subquery result. Mathematically equivalent to:
    \[
    v > \min_{x \in S}(x)
    \]
    \item \textbf{\texttt{v > ALL (Subquery)}:} True if $v$ is greater than \textbf{every} element in the subquery result. Mathematically equivalent to:
    \[
    v > \max_{x \in S}(x)
    \]
\end{itemize}

% =============================================================================
% SECTION 24: SQL AGGREGATE FUNCTIONS & NULL SEMANTICS
% =============================================================================
\section{SQL Aggregate Functions, Statistical Calculations, and NULL Handling}
\label{sec:dbms_sql_aggregates}

\subsection{Mathematical Taxonomy of Aggregate Functions}
An \textbf{aggregate function} maps a multiset of domain values from multiple rows to a single scalar summary value:
\begin{equation}
f: \mathcal{M}(D) \to D_{\text{scalar}}
\end{equation}

\begin{table}[htbp]
\centering
\small
\caption{Standard ANSI SQL Aggregate Functions}
\label{tab:sql_aggregates}
\begin{tabularx}{\textwidth}{l p{2.8cm} X l}
\toprule
\textbf{Function} & \textbf{Mathematical Definition} & \textbf{Relational Semantics} & \textbf{NULL Handling} \\
\midrule
\texttt{COUNT(*)} & $|\mathcal{R}|$ & Counts total candidate rows in relation. & Preserves NULLs. \\
\texttt{COUNT(col)} & $|\{t \in \mathcal{R} \mid t[\text{col}] \neq \text{NULL}\}|$ & Counts non-NULL values in specified attribute. & Ignores NULLs. \\
\texttt{COUNT(DISTINCT col)} & $|\text{set}(t[\text{col}]) \setminus \{\text{NULL}\}|$ & Counts unique non-NULL values. & Ignores NULLs. \\
\texttt{SUM(col)} & $\sum_{i=1}^{k} x_i$ & Arithmetic sum over non-NULL numeric multiset. & Ignores NULLs. \\
\texttt{AVG(col)} & $\frac{\sum x_i}{k}$ & Arithmetic mean ($\text{SUM(col)} / \text{COUNT(col)}$). & Ignores NULLs. \\
\texttt{MIN(col)} & $\min(x_1, \dots, x_k)$ & Minimum value according to collation order. & Ignores NULLs. \\
\texttt{MAX(col)} & $\max(x_1, \dots, x_k)$ & Maximum value according to collation order. & Ignores NULLs. \\
\bottomrule
\end{tabularx}
\end{table}

\begin{quote}
\textbf{Fundamental Examination Property: Aggregation over NULLs}
\begin{enumerate}
    \item Except for \texttt{COUNT(*)}, \textbf{all aggregate functions strictly ignore NULL values}.
    \item If an aggregate function (\texttt{SUM}, \texttt{AVG}, \texttt{MIN}, \texttt{MAX}) is executed over a column containing \textbf{only NULL values} (or an empty relation), the result is \textbf{\texttt{NULL}}, whereas \texttt{COUNT} over an empty set returns \textbf{\texttt{0}}.
    \item In \texttt{AVG(col)}, the denominator is $\text{COUNT}(col)$, \textbf{not} $\text{COUNT}(*)$. Rows with \texttt{NULL} do not deflate the computed average.
\end{enumerate}
\end{quote}

% =============================================================================
% SECTION 25: GROUP BY & THE WHERE VS. HAVING PARADIGM
% =============================================================================
\section{Data Aggregation and Group Filtering: GROUP BY and HAVING}
\label{sec:dbms_sql_group_by}

\subsection{The Equivalence Partitioning of GROUP BY}
The \texttt{GROUP BY} clause collapses all tuples sharing identical values in the specified grouping attributes into a single summary row.
\begin{lstlisting}[language=SQL]
SELECT Department, COUNT(*) AS Total_Employees, AVG(Salary) AS Avg_Salary
FROM Employees
GROUP BY Department;
\end{lstlisting}

\textbf{The ANSI SQL Single-Value Rule:} Every non-aggregated attribute appearing in the \texttt{SELECT} projection list \textbf{must strictly appear in the \texttt{GROUP BY} clause}. Failure to comply causes relational indeterminacy (a single output row would have to arbitrarily select between conflicting values).

\subsection{The WHERE versus HAVING Paradigm}
The distinction between \texttt{WHERE} and \texttt{HAVING} is one of the most frequently tested concepts in technical recruitment examinations:

\begin{table}[htbp]
\centering
\small
\caption{Comprehensive Comparison: WHERE versus HAVING}
\label{tab:where_vs_having}
\begin{tabularx}{\textwidth}{p{3.2cm} X X}
\toprule
\textbf{Feature Dimension} & \textbf{WHERE Clause} & \textbf{HAVING Clause} \\
\midrule
\textbf{Filtering Target} & Individual candidate tuples (rows). & Aggregated groups of tuples. \\
\textbf{Pipeline Position} & Executed \textbf{before} \texttt{GROUP BY} aggregation. & Executed \textbf{after} \texttt{GROUP BY} aggregation. \\
\textbf{Aggregate Functions} & Strictly disallowed (e.g., \texttt{WHERE AVG(s) > 50} is illegal). & Allowed and standard (e.g., \texttt{HAVING AVG(s) > 50}). \\
\textbf{Index Utilization} & Can utilize B-tree indexes for fast index range scans. & Operates on temporary in-memory group buffers; cannot use raw table indexes. \\
\textbf{Primary Objective} & Minimize the volume of rows fed into the grouping engine. & Discard computed groups that fail summary criteria. \\
\bottomrule
\end{tabularx}
\end{table}

\subsection{Multi-Stage Execution Pipeline Example}
\begin{lstlisting}[language=SQL]
SELECT Department_ID, COUNT(*) AS Headcount, AVG(Salary) AS Avg_Salary
FROM Employees
WHERE Salary >= 40000.00                   -- Step 1: Filter raw employee rows
GROUP BY Department_ID                     -- Step 2: Form groups by department
HAVING COUNT(*) > 1                        -- Step 3: Filter groups with > 1 employee
ORDER BY Avg_Salary DESC;                  -- Step 4: Sort resulting group records
\end{lstlisting}

% =============================================================================
% SECTION 26: RELATIONAL JOIN ARCHITECTURES & SET OPERATIONS
% =============================================================================
\section{Relational JOIN Architectures and Set Operations}
\label{sec:dbms_sql_joins}

\subsection{Relational Join Mechanics}
A \textbf{JOIN} is an algebraic binary operator that combines attributes from two relations ($R$ and $S$) based on a join condition (predicate).

\begin{figure}[htbp]
\centering
\small
\begin{tabularx}{\textwidth}{l p{4.2cm} X}
\toprule
\textbf{Join Type} & \textbf{Mathematical Definition} & \textbf{Relational Behavior} \\
\midrule
\textbf{INNER JOIN} & $R \bowtie_{\theta} S = \sigma_{\theta}(R \times S)$ & Retains only tuples that satisfy the join predicate in both relations. \\
\textbf{LEFT OUTER JOIN} & $R \text{ } {}_{{}^{\text{L}}}\!\bowtie_{\theta} S$ & Retains all tuples from $R$; un-matched attributes from $S$ are populated with \texttt{NULL}. \\
\textbf{RIGHT OUTER JOIN} & $R \bowtie_{\theta}^{\text{R}}\text{ } S$ & Retains all tuples from $S$; un-matched attributes from $R$ are populated with \texttt{NULL}. \\
\textbf{FULL OUTER JOIN} & $R \text{ } {}_{{}^{\text{L}}}\!\bowtie_{\theta}^{\text{R}}\text{ } S$ & Retains all tuples from both $R$ and $S$; pads non-matching sides with \texttt{NULL}. \\
\textbf{CROSS JOIN} & $R \times S$ & Unconstrained Cartesian product ($|R| \times |S|$ total tuples). \\
\textbf{SELF JOIN} & $R \bowtie_{\theta} R$ (aliased as $R_1, R_2$) & Joins a table to itself to model recursive hierarchies. \\
\bottomrule
\end{tabularx}
\end{figure}

\subsection{SELF JOIN for Hierarchical Traversal}
The self-join is standard for representing organizational reporting lines, where a tuple's foreign key points to the primary key of the same relation:
\begin{lstlisting}[language=SQL]
-- Retrieve employee names along with their direct supervisor's name
SELECT e.Name AS Employee_Name, m.Name AS Manager_Name
FROM Employees e
LEFT JOIN Employees m ON e.Manager_ID = m.Employee_ID;
\end{lstlisting}
Using a \texttt{LEFT JOIN} ensures that top-level executives (whose \texttt{Manager\_ID} is \texttt{NULL}) are preserved in the result.

\subsection{Horizontal versus Vertical Operations: JOIN versus UNION}
\begin{itemize}
    \item \textbf{\texttt{JOIN}:} Combines relations \textbf{horizontally} by appending columns from related tables based on common keys.
    \item \textbf{\texttt{UNION}:} Combines relations \textbf{vertically} by concatenating rows from two union-compatible queries.
\end{itemize}

\begin{quote}
\textbf{UNION versus UNION ALL}
\begin{itemize}
    \item \textbf{\texttt{UNION}:} Merges row sets and \textbf{eliminates duplicate tuples}. Requires an implicit sorting or hashing pass over all returned records, introducing significant CPU and memory overhead on large datasets.
    \item \textbf{\texttt{UNION ALL}:} Concatenates row sets directly \textbf{without checking for duplicates}. Incurs virtually zero overhead. Rule: Always use \texttt{UNION ALL} unless deduplication is an explicit business requirement.
\end{itemize}
\end{quote}

% =============================================================================
% SECTION 27: REFERENTIAL ACTIONS & FOREIGN KEY CASCADE DYNAMICS
% =============================================================================
\section{Referential Actions and Foreign Key Cascade Dynamics}
\label{sec:dbms_sql_referential_actions}

\subsection{Foreign Key Constraints and Referential Integrity}
Referential integrity dictates that every non-null foreign key value in a referencing (child) relation must match an existing primary/candidate key value in the referenced (parent) relation. When a referenced parent tuple is modified or deleted, the DBMS enforces one of four standardized \textbf{referential actions}:

\begin{table}[htbp]
\centering
\small
\caption{Foreign Key Referential Action Matrix}
\label{tab:referential_actions}
\begin{tabularx}{\textwidth}{l X X}
\toprule
\textbf{Action} & \textbf{Behavior on Parent Tuple Deletion} & \textbf{Behavior on Parent Key Update} \\
\midrule
\texttt{CASCADE} & Automatically deletes all child rows referencing the deleted parent. & Automatically updates referencing child foreign keys to the new key value. \\
\texttt{SET NULL} & Sets the referencing child foreign key columns to \texttt{NULL} (requires nullable FK). & Sets referencing child foreign key columns to \texttt{NULL}. \\
\texttt{SET DEFAULT} & Sets the referencing child foreign key columns to their declared default values. & Sets referencing child foreign key columns to their declared default values. \\
\texttt{RESTRICT} & Immediately rejects the parent deletion if any referencing child exists. & Immediately rejects the parent key update if any referencing child exists. \\
\texttt{NO ACTION} & Rejects the operation if the constraint is violated when checked (allows deferred checking). & Rejects the operation if the constraint is violated when checked. \\
\bottomrule
\end{tabularx}
\end{table}

\subsection{Production Banking Architecture Rule}
\begin{quote}
\textbf{Production Principle: Never Use CASCADE on Financial Ledgers}
In financial core banking systems, transactional history must remain immutable for legal, compliance, and auditing requirements. If an account or customer profile is deactivated or closed, applying \texttt{ON DELETE CASCADE} to the \texttt{Bank\_Transactions} table would delete historical deposit and withdrawal ledgers, destroying financial auditability. Real-world banking engines enforce \textbf{\texttt{RESTRICT}} on ledger tables and employ \textbf{soft deletes} (\texttt{Status = 'Closed'}) rather than physical tuple destruction.
\end{quote}



% =============================================================================
% SECTION: FORMAL RELATIONAL QUERY LANGUAGES (ALGEBRA & CALCULUS)
% =============================================================================
\section{Formal Relational Query Languages: Relational Algebra and Relational Calculus}
\label{sec:dbms_relational_algebra}

\subsection{Relational Algebra: The Procedural Foundation}
\textbf{Relational Algebra} is a formal procedural query language that operates on one or more relations to produce another relation as its output without mutating the underlying operands. It is termed \textit{procedural} because an algebraic expression prescribes the step-by-step mathematical sequence of operational transforms required to compute the resultant relation. Relational algebra provides the theoretical foundation for relational database query optimization and execution trees.

\begin{table}[htbp]
\centering
\small
\caption{Fundamental Operators of Relational Algebra and SQL Equivalents}
\label{tab:relational_algebra_ops}
\begin{tabularx}{\textwidth}{l c p{5.5cm} X}
\toprule
\textbf{Operator} & \textbf{Symbol} & \textbf{Relational Algebraic Definition} & \textbf{SQL DQL Equivalent} \\
\midrule
\textbf{Selection} & $\sigma_{\theta}(R)$ & Filters tuples satisfying predicate $\theta$: $\{t \in R \mid \theta(t) = \text{true}\}$ & \texttt{WHERE} clause filter \\
\textbf{Projection} & $\pi_{A_1, \dots, A_k}(R)$ & Extracts attributes $A_1, \dots, A_k$ and eliminates duplicate tuples & \texttt{SELECT DISTINCT col1, col2} \\
\textbf{Union} & $R \cup S$ & Set union of two union-compatible relations & \texttt{UNION} operator \\
\textbf{Set Difference} & $R - S$ & Tuples present in relation $R$ but absent in relation $S$ & \texttt{EXCEPT} (PostgreSQL) / \texttt{MINUS} (Oracle) \\
\textbf{Cartesian Product} & $R \times S$ & Pairs every tuple of $R$ with every tuple of $S$ & \texttt{CROSS JOIN} \\
\textbf{Rename} & $\rho_{S}(R)$ & Renames relation $R$ to $S$ or attributes to $B_1, \dots, B_n$ & Column and Table Aliases (\texttt{AS}) \\
\textbf{Theta Join} & $R \bowtie_{\theta} S$ & Equi/theta combination: $\sigma_{\theta}(R \times S)$ & \texttt{INNER JOIN ... ON} \\
\textbf{Natural Join} & $R \bowtie S$ & Equijoin on all shared attribute names with deduplication & \texttt{NATURAL JOIN} \\
\bottomrule
\end{tabularx}
\end{table}

\subsection{Relational Calculus: The Declarative Foundation}
In contrast to relational algebra, \textbf{Relational Calculus} is a non-procedural, declarative query language based on first-order predicate logic. Rather than specifying operational steps, an expression in relational calculus states \textit{what} relational criteria the resultant tuples must satisfy.
\begin{itemize}
    \item \textbf{Tuple Relational Calculus (TRC):} Formulates queries using tuple variables that range over relations:
    \begin{equation}
    \{t \mid P(t)\}
    \end{equation}
    where $t$ is a tuple variable and $P(t)$ is a first-order logic formula. For example, retrieving all CSE students:
    \begin{equation}
    \{t \mid t \in \text{STUDENT} \land t[\text{Department}] = \text{'CSE'}\}
    \end{equation}
    \item \textbf{Domain Relational Calculus (DRC):} Formulates queries using domain variables that range over individual attribute domains rather than entire tuples:
    \begin{equation}
    \{\langle x_1, x_2, \dots, x_n \rangle \mid P(x_1, x_2, \dots, x_n)\}
    \end{equation}
    For example:
    \begin{equation}
    \{\langle i, n, d, c \rangle \mid \langle i, n, d, c \rangle \in \text{STUDENT} \land d = \text{'CSE'}\}
    \end{equation}
\end{itemize}

\begin{quote}
\textbf{Core Theoretical Principle: Codd's Equivalence Theorem}
E. F. Codd proved that the expressive power of basic Relational Algebra is mathematically equivalent to the safe subsets of Tuple Relational Calculus and Domain Relational Calculus. A query language that possesses this baseline expressive power is defined as \textbf{Relational Complete}. SQL is a declarative realization of relational calculus with algebraic execution foundations.
\end{quote}

% =============================================================================
% SECTION: DATABASE VIEWS & MATERIALIZED VIEWS
% =============================================================================
\section{Database Views, Updatability, and Materialized Views}
\label{sec:dbms_views}

\subsection{Virtual Database Views}
A \textbf{Database View} is a virtual table defined as a stored SQL query expression in the system catalog. A standard view contains no physically materialized row records on disk; when a client executes a query against a view, the RDBMS query rewriter merges the view definition with the user's outer query into a unified abstract syntax tree (AST).
\begin{lstlisting}[language=SQL]
CREATE VIEW Employee_Public_Info AS
SELECT Employee_ID, Name, Department
FROM Employees;
\end{lstlisting}

\subsection{Updatable Views and the WITH CHECK OPTION}
Under specific constraints, DML operations (\texttt{INSERT}, \texttt{UPDATE}, \texttt{DELETE}) performed against a view propagate to the underlying base table. A view is updatable if:
\begin{enumerate}
    \item It references exactly one base table in its \texttt{FROM} clause.
    \item It does not contain \texttt{DISTINCT}, \texttt{GROUP BY}, \texttt{HAVING}, or aggregate functions.
    \item The base table primary key and all non-nullable columns without defaults are present.
\end{enumerate}

\textbf{\texttt{WITH CHECK OPTION}:} Guarantees that any row inserted or updated through the view strictly conforms to the view's defining \texttt{WHERE} predicate, preventing data corruption:
\begin{lstlisting}[language=SQL]
CREATE VIEW IT_Employees AS
SELECT Employee_ID, Name, Department
FROM Employees
WHERE Department = 'IT'
WITH CHECK OPTION;
\end{lstlisting}
If a client attempts to execute \texttt{UPDATE IT\_Employees SET Department = 'HR' WHERE Employee\_ID = 101}, the RDBMS engine immediately throws a constraint violation error.

\subsection{Materialized Views}
A \textbf{Materialized View} physically persists the calculated query result on disk, acting as a concrete physical table with associated B-tree indexes. It is engineered for expensive analytical reporting over massive tables where re-executing complex aggregations on every request would degrade OLTP throughput.
\begin{lstlisting}[language=SQL]
-- PostgreSQL Materialized View Definition
CREATE MATERIALIZED VIEW Department_Summary AS
SELECT Department, COUNT(*) AS Total_Employees, AVG(Salary) AS Average_Salary
FROM Employees
GROUP BY Department;

-- Query the precomputed physical result
SELECT * FROM Department_Summary WHERE Average_Salary > 50000;

-- Refresh the materialized table after OLTP bulk updates
REFRESH MATERIALIZED VIEW Department_Summary;
\end{lstlisting}

\begin{table}[htbp]
\centering
\small
\caption{Architectural Comparison: Standard Virtual View versus Materialized View}
\label{tab:view_vs_materialized_view}
\begin{tabularx}{\textwidth}{p{3.2cm} X X}
\toprule
\textbf{Feature Dimension} & \textbf{Standard Virtual View} & \textbf{Materialized View} \\
\midrule
\textbf{Physical Storage} & Zero data disk storage (stores only SQL query text). & Allocates physical disk pages for precomputed tuples and indexes. \\
\textbf{Query Processing} & Dynamic query expansion: executes underlying SQL at query time. & Direct sequential/index read from pre-aggregated physical pages. \\
\textbf{Data Freshness} & Always 100\% real-time consistent with base table mutations. & Asynchronous / Stale until explicitly or incrementally refreshed. \\
\textbf{Write Overhead} & Zero base-table write overhead. & May require refresh scheduling or write-amplification triggers. \\
\textbf{Optimal Application} & Access control, row/column security, query abstraction. & OLAP dashboards, data warehousing, expensive multi-table reporting. \\
\bottomrule
\end{tabularx}
\end{table}

% =============================================================================
% SECTION: COMMON TABLE EXPRESSIONS (CTE) & RECURSIVE QUERIES
% =============================================================================
\section{Common Table Expressions (CTE) and Recursive Hierarchical Queries}
\label{sec:dbms_cte}

\subsection{Common Table Expressions (The WITH Clause)}
A \textbf{Common Table Expression (CTE)} is a temporary named result set defined within the execution scope of a single SQL statement (\texttt{SELECT}, \texttt{INSERT}, \texttt{UPDATE}, or \texttt{DELETE}) using the \texttt{WITH} clause. CTEs decompose unwieldy nested subqueries into modular, readable blocks:
\begin{lstlisting}[language=SQL]
WITH Dept_Avg AS (
    SELECT Department_ID, AVG(Salary) AS Avg_Salary
    FROM Employees
    GROUP BY Department_ID
),
High_Paying_Dept AS (
    SELECT Department_ID, Avg_Salary
    FROM Dept_Avg
    WHERE Avg_Salary > 50000.00
)
SELECT e.Name, e.Salary, h.Avg_Salary
FROM Employees e
INNER JOIN High_Paying_Dept h ON e.Department_ID = h.Department_ID;
\end{lstlisting}

\subsection{Recursive CTEs: Hierarchical Graph Traversal}
A \textbf{Recursive CTE} references its own intermediate result set to traverse acyclic hierarchical trees (e.g., organizational reporting lines, bill-of-materials, folder structures). A recursive CTE contains two distinct computational members joined by \texttt{UNION ALL}:
\begin{enumerate}
    \item \textbf{Anchor Member:} Computes the initial starting base relation (the root of the tree).
    \item \textbf{Recursive Member:} Iteratively joins the recursive relation with the base table until the join yields an empty result set (termination condition).
\end{enumerate}

\begin{lstlisting}[language=SQL]
-- Traversal of complete organizational reporting hierarchy starting at CEO (ID=101)
WITH RECURSIVE Employee_Hierarchy AS (
    -- 1. Anchor Member: Select top-level manager
    SELECT Employee_ID, Name, Manager_ID, 1 AS Hierarchy_Level
    FROM Employees
    WHERE Employee_ID = 101

    UNION ALL

    -- 2. Recursive Member: Join subordinate employees with hierarchy result
    SELECT e.Employee_ID, e.Name, e.Manager_ID, h.Hierarchy_Level + 1
    FROM Employees e
    INNER JOIN Employee_Hierarchy h ON e.Manager_ID = h.Employee_ID
)
SELECT * FROM Employee_Hierarchy ORDER BY Hierarchy_Level, Employee_ID;
\end{lstlisting}

% =============================================================================
% SECTION: SQL WINDOW FUNCTIONS & ANALYTIC EVALUATION
% =============================================================================
\section{SQL Window Functions, Ranking, and Analytic Frames}
\label{sec:dbms_window_functions}

\subsection{The Analytic Window Evaluation Paradigm}
Unlike aggregate functions with \texttt{GROUP BY} (which collapse distinct tuples into single grouped summary rows), a \textbf{Window Function} performs calculations across a set of related rows while \textbf{strictly preserving individual row identities and cardinalities} in the final output. The calculation scope is governed by the \texttt{OVER} clause:
\begin{equation}
\text{Function}(A) \quad \text{OVER} \quad \left( \text{PARTITION BY } P_1, \dots, P_k \quad \text{ORDER BY } O_1, \dots, O_m \quad \text{Frame} \right)
\end{equation}

\subsection{Ranking Functions: ROW\_NUMBER versus RANK versus DENSE\_RANK}
The distinctions between ranking functions under tied score values are a major focus in technical examinations:

\begin{table}[htbp]
\centering
\small
\caption{Comparative Matrix: SQL Ranking Functions}
\label{tab:ranking_functions}
\begin{tabularx}{\textwidth}{l p{4.2cm} c X}
\toprule
\textbf{Function} & \textbf{Tie-Breaking Behavior} & \textbf{Rank Sequence} & \textbf{Use Case} \\
\midrule
\texttt{ROW\_NUMBER()} & Assigns strictly unique sequential integers; ties resolved arbitrarily or via secondary sort keys. & 1, 2, 3, 4, 5 & Pagination offsets, identifying duplicate records for deletion. \\
\texttt{RANK()} & Assigns identical rank to ties; \textbf{skips subsequent numbers} based on duplicate count. & 1, 2, 2, 4, 5 & Competitive sports rankings, academic grade percentiles. \\
\texttt{DENSE\_RANK()} & Assigns identical rank to ties; \textbf{never skips subsequent integers} (no gaps). & 1, 2, 2, 3, 4 & Top-N salary extraction, multi-tier pricing brackets. \\
\bottomrule
\end{tabularx}
\end{table}

\subsection{Finding the Top Two Salaries per Department via DENSE\_RANK}
\begin{lstlisting}[language=SQL]
WITH Department_Ranks AS (
    SELECT Employee_ID, Name, Department_ID, Salary,
           DENSE_RANK() OVER (
               PARTITION BY Department_ID 
               ORDER BY Salary DESC
           ) AS Salary_Tier
    FROM Employees
)
SELECT * FROM Department_Ranks WHERE Salary_Tier <= 2;
\end{lstlisting}

\subsection{Navigation Functions and Cumulative Running Totals}
\begin{itemize}
    \item \textbf{\texttt{LAG(col, offset)} and \texttt{LEAD(col, offset)}:} Accesses attribute values from preceding or succeeding tuples within the partition without executing an expensive self-join. Essential for computing period-over-period financial growth.
    \item \textbf{Cumulative Running Totals with Explicit Window Framing:}
\begin{lstlisting}[language=SQL]
SELECT Transaction_ID, Created_At, Amount,
       SUM(Amount) OVER (
           ORDER BY Created_At, Transaction_ID
           ROWS BETWEEN UNBOUNDED PRECEDING AND CURRENT ROW
       ) AS Cumulative_Ledger_Balance
FROM Bank_Transactions;
\end{lstlisting}
\end{itemize}

% =============================================================================
% SECTION: PROCEDURAL SQL: PROCEDURES, FUNCTIONS, TRIGGERS & CURSORS
% =============================================================================
\section{Procedural SQL: Stored Procedures, Functions, Triggers, and Cursors}
\label{sec:dbms_procedural_sql}

\subsection{Stored Procedures versus User-Defined Functions}
\begin{table}[htbp]
\centering
\small
\caption{Comparative Analysis: Stored Procedures versus Stored Functions}
\label{tab:procedures_vs_functions}
\begin{tabularx}{\textwidth}{p{3.2cm} X X}
\toprule
\textbf{Architectural Metric} & \textbf{Stored Procedure} & \textbf{Stored Function (UDF)} \\
\midrule
\textbf{Primary Objective} & Executes multi-step business transactions and administrative batch workloads. & Computes and returns an encapsulated transformation value or table. \\
\textbf{Return Specification} & Optional; returns zero, one, or multiple parameters via \texttt{OUT} / \texttt{INOUT} modes. & Mandatory; must return a scalar value or a defined relation (\texttt{TABLE}). \\
\textbf{Invocation Method} & Explicit standalone invocation via \texttt{CALL proc\_name()} or \texttt{EXEC}. & Directly embedded within SQL expressions (\texttt{SELECT}, \texttt{WHERE}, \texttt{HAVING}). \\
\textbf{Transaction Control} & Permitted to commit or roll back transactions internally (\texttt{COMMIT}/\texttt{ROLLBACK}). & Strictly prohibited from executing transaction control statements in standard SQL. \\
\bottomrule
\end{tabularx}
\end{table}

\subsection{Database Triggers: Event-Driven Integrity and Auditing}
A \textbf{Trigger} is an active database rule that fires automatically upon the occurrence of a specified data manipulation event (\texttt{INSERT}, \texttt{UPDATE}, \texttt{DELETE}) or DDL mutation.
\begin{itemize}
    \item \textbf{Timing Constraints:} \texttt{BEFORE} (pre-operation validation/transformation), \texttt{AFTER} (post-commit audit logging), and \texttt{INSTEAD OF} (diverts DML on complex non-updatable views).
    \item \textbf{Granularity Constraints:} \texttt{FOR EACH ROW} (fires iteratively for every mutated tuple) vs. \texttt{FOR EACH STATEMENT} (fires exactly once per SQL statement regardless of affected rows).
    \item \textbf{Pseudorecords:} \texttt{OLD} (captures pre-mutation tuple state in \texttt{UPDATE}/\texttt{DELETE}) and \texttt{NEW} (captures candidate tuple state in \texttt{INSERT}/\texttt{UPDATE}).
\end{itemize}

\begin{lstlisting}[language=SQL]
-- PostgreSQL PL/pgSQL Comprehensive Salary Audit Trigger
CREATE TABLE Salary_Audit (
    Audit_ID    BIGINT GENERATED ALWAYS AS IDENTITY PRIMARY KEY,
    Employee_ID INT NOT NULL,
    Old_Salary  DECIMAL(12, 2),
    New_Salary  DECIMAL(12, 2),
    Mutated_At  TIMESTAMP DEFAULT CURRENT_TIMESTAMP
);

CREATE OR REPLACE FUNCTION audit_employee_salary()
RETURNS TRIGGER AS $$
BEGIN
    IF OLD.Salary IS DISTINCT FROM NEW.Salary THEN
        INSERT INTO Salary_Audit (Employee_ID, Old_Salary, New_Salary)
        VALUES (OLD.Employee_ID, OLD.Salary, NEW.Salary);
    END IF;
    RETURN NEW;
END;
$$ LANGUAGE plpgsql;

CREATE TRIGGER trg_salary_mutation
AFTER UPDATE OF Salary ON Employees
FOR EACH ROW
EXECUTE FUNCTION audit_employee_salary();
\end{lstlisting}

\subsection{Database Cursors: Row-by-Row Procedural Processing}
While standard SQL processes sets of tuples simultaneously in declarative batches, a \textbf{Cursor} provides a procedural mechanism to iterate through query results one tuple at a time:
\begin{equation}
\text{DECLARE} \implies \text{OPEN} \implies \text{FETCH} \implies \text{PROCESS} \implies \text{CLOSE} \implies \text{DEALLOCATE}
\end{equation}
\textbf{Performance Warning:} Cursors bypass the relational set optimizer, incur substantial context-switching overhead between the SQL and procedural execution engines, and prolong transactional lock duration. Rule: Exhaust set-based SQL joins and window functions before resorting to procedural cursors.

% =============================================================================
% SECTION: NULL VALUES & SQL'S THREE-VALUED LOGIC (3VL)
% =============================================================================
\section{NULL Values, Three-Valued Logic (3VL), and Defensive Functions}
\label{sec:dbms_null_3vl}

\subsection{The Semantics of NULL in Relational Theory}
In relational database systems, \texttt{NULL} is not a data value: it is a syntactic marker representing the \textbf{absence of a value} (unknown, missing, unrecorded, or inapplicable). It is categorically distinct from numeric zero ($0$), a blank space (\texttt{' '}), an empty string (\texttt{''}), or boolean \texttt{FALSE}.

\subsection{Three-Valued Logic (3VL) Truth Tables}
Because any comparison against an unknown value yields an unknown result, SQL extends classical two-valued Boolean logic ($\{\text{TRUE}, \text{FALSE}\}$) to include \textbf{\texttt{UNKNOWN}}:

\begin{table}[htbp]
\centering
\small
\caption{SQL Three-Valued Logic Truth Tables}
\label{tab:3vl_truth_tables}
\begin{minipage}{0.32\textwidth}
\centering
\begin{tabular}{l l l}
\toprule
$P$ & $Q$ & $P \land Q$ \\
\midrule
T & T & \textbf{T} \\
T & F & \textbf{F} \\
T & U & \textbf{U} \\
F & U & \textbf{F} \\
U & U & \textbf{U} \\
\bottomrule
\end{tabular}
\end{minipage}
\hfill
\begin{minipage}{0.32\textwidth}
\centering
\begin{tabular}{l l l}
\toprule
$P$ & $Q$ & $P \lor Q$ \\
\midrule
T & T & \textbf{T} \\
T & F & \textbf{T} \\
T & U & \textbf{T} \\
F & U & \textbf{U} \\
U & U & \textbf{U} \\
\bottomrule
\end{tabular}
\end{minipage}
\hfill
\begin{minipage}{0.32\textwidth}
\centering
\begin{tabular}{l l}
\toprule
$P$ & $\neg P$ \\
\midrule
T & \textbf{F} \\
F & \textbf{T} \\
U & \textbf{U} \\
\bottomrule
\end{tabular}
\end{minipage}
\end{table}

\textbf{The WHERE Acceptance Filter:} In SQL, the \texttt{WHERE} and \texttt{HAVING} clauses evaluate boolean predicates on tuples and retain \textbf{only those rows that evaluate strictly to \texttt{TRUE}}. Tuples evaluating to \texttt{FALSE} or \texttt{UNKNOWN} are discarded.

\subsection{Defensive Null-Handling Functions}
\begin{itemize}
    \item \textbf{\texttt{COALESCE(e1, e2, ..., en)}:} Evaluates arguments in strict left-to-right order and returns the first non-null expression:
\begin{lstlisting}[language=SQL]
SELECT Name, COALESCE(Phone, Mobile, 'Unreachable') AS Primary_Contact FROM Users;
\end{lstlisting}
    \item \textbf{\texttt{NULLIF(a, b)}:} Returns \texttt{NULL} if $a = b$; otherwise returns $a$. Highly utilized as a defensive mathematical guard against fatal runtime divide-by-zero errors:
\begin{lstlisting}[language=SQL]
-- Safe division: if Total_Orders is 0, NULLIF returns NULL, preventing division error
SELECT Total_Revenue / NULLIF(Total_Orders, 0) AS Average_Order_Value FROM Sales;
\end{lstlisting}
\end{itemize}



% =============================================================================
% SECTION 28: DATABASE INDEXING ARCHITECTURES & PERFORMANCE OPTIMIZATION
% =============================================================================
\section{Database Indexing Architectures and Query Optimization}
\label{sec:dbms_sql_indexes}

\subsection{The Engineering Motivation for Indexing}
An \textbf{index} is an auxiliary data structure maintained by the RDBMS to accelerate the search and retrieval of records without performing an exhaustive \textbf{Full Table Scan} (scanning every data page on disk from beginning to end, incurring $O(N)$ I/O operations). Just as the index at the back of a textbook maps keywords directly to page numbers, a database index maps search key values to their corresponding physical storage addresses (\textbf{Tuple Identifiers / RowIDs}).

\begin{figure}[htbp]
\centering
\small
\begin{tabularx}{\textwidth}{p{3.5cm} X}
\toprule
\textbf{Database Operation} & \textbf{Index Performance Acceleration Mechanism} \\
\midrule
\textbf{Point Lookups} & Replaces $O(N)$ sequential page scans with $O(\log N)$ B-tree traversals. \\
\textbf{Range Scans} & Navigates to the range lower bound, then performs a sequential scan along the leaf node linked list. \\
\textbf{Relational JOINs} & Accelerates inner table probe phase in Index Nested Loop Joins. \\
\textbf{Sorting (ORDER BY)} & Retrieves tuples in pre-sorted B-tree leaf order, eliminating in-memory / tempdb disk sorts. \\
\textbf{Unique Enforcement} & Immediately verifies key uniqueness during \texttt{INSERT} and \texttt{UPDATE} via B-tree node searches. \\
\bottomrule
\end{tabularx}
\end{figure}

\subsection{Index Types: Clustered versus Non-Clustered}
\begin{itemize}
    \item \textbf{Clustered Index:} Dictates the \textbf{physical ordering} of actual table rows on disk data pages. Because table rows can only be physically sorted in one sequence on disk, \textbf{there can be at most one clustered index per table}. In modern engines (e.g., MySQL InnoDB), the primary key automatically acts as the clustered index; the leaf nodes of the B+ tree contain the complete row data pages.
    \item \textbf{Non-Clustered (Secondary) Index:} A standalone B+ tree structure maintained completely separately from the physical table data pages. The leaf nodes contain the indexed column values paired with a \textbf{row locator} (a physical RowID in heap-organized tables, or the clustered primary key value in clustered tables). Traversing a secondary index to fetch non-indexed attributes requires an extra I/O lookup called a \textbf{Bookmark / Key Lookup}.
    \item \textbf{Composite (Multi-Column) Index:} An index defined over multiple attributes:
\begin{lstlisting}[language=SQL]
CREATE INDEX idx_emp_dept_salary ON Employees(Department_ID, Salary);
\end{lstlisting}
    \textbf{The Leftmost Prefix Rule:} A composite index on $(A, B)$ can accelerate queries filtering on $A$ alone, or $A$ and $B$ together, but \textbf{cannot} accelerate queries filtering on $B$ alone without scanning the entire index tree.
\end{itemize}

\subsection{The Index Performance Trade-Off Matrix}
\begin{table}[htbp]
\centering
\small
\caption{The Dual Nature of Database Indexing}
\label{tab:index_tradeoffs}
\begin{tabularx}{\textwidth}{X X}
\toprule
\textbf{Architectural Advantages} & \textbf{Engineering Costs and Overhead} \\
\midrule
Accelerates \texttt{SELECT} point lookups from $O(N)$ to $O(\log N)$. & Incurs additional storage overhead (indexes can exceed base table size). \\
Reduces disk I/O, server CPU cycles, and RAM buffer pool pressure. & Drastically degrades write throughput: every \texttt{INSERT}, \texttt{DELETE}, and indexed \texttt{UPDATE} must update all B-tree leaf pages. \\
Enables \textbf{Covering Index} scans (retrieving data entirely from index leaf nodes without touching heap tables). & Causes B-tree page splits and storage fragmentation under random primary key insertions. \\
\bottomrule
\end{tabularx}
\end{table}

% =============================================================================
% SECTION 29: JAVA DATABASE CONNECTIVITY (JDBC) ARCHITECTURE
% =============================================================================
\section{Java Database Connectivity (JDBC) Architecture and Security}
\label{sec:dbms_jdbc}

\subsection{Definition and Purpose of JDBC}
\textbf{Java Database Connectivity (JDBC)} is a standard Java API (part of the Java Standard Edition specification under \texttt{java.sql} and \texttt{javax.sql}) that defines a vendor-neutral interface enabling Java applications to establish connections with relational database engines, execute SQL queries and updates, and navigate structured tabular results.

\subsection{The Four-Tier JDBC Architectural Hierarchy}
\begin{figure}[htbp]
\centering
\begin{tcolorbox}[colback=codebg, colframe=primary, title=\textbf{The Canonical JDBC Component Hierarchy}]
\begin{enumerate}
    \item \textbf{Java Enterprise Application:} The business logic client that initiates transactional persistence calls.
    \item \textbf{JDBC API (\texttt{java.sql.*}):} Abstract interfaces defining standard database interactions (\texttt{Connection}, \texttt{Statement}, \texttt{PreparedStatement}, \texttt{ResultSet}).
    \item \textbf{JDBC Driver Manager / DataSource:} The registry that loads vendor drivers and manages connection pooling.
    \item \textbf{Database-Specific JDBC Driver:} Vendor-supplied implementation of the JDBC SPI that translates Java method calls into the proprietary wire protocol of the target RDBMS.
    \item \textbf{RDBMS Server:} The physical database engine (e.g., Oracle, PostgreSQL, MySQL).
\end{enumerate}
\end{tcolorbox}
\end{figure}

\subsection{The Four Types of JDBC Drivers}
Recruitment examinations frequently test the classification, architecture, and network topology of JDBC drivers:

\begin{table}[htbp]
\centering
\small
\caption{Comparative Architectural Analysis of the Four JDBC Driver Types}
\label{tab:jdbc_driver_types}
\begin{tabularx}{\textwidth}{p{2.6cm} p{3.2cm} p{3.8cm} X}
\toprule
\textbf{Driver Type} & \textbf{Common Name} & \textbf{Architectural Communication Path} & \textbf{Engineering Evaluation} \\
\midrule
\textbf{Type 1} & JDBC-ODBC Bridge & Java App $\to$ JDBC $\to$ ODBC Bridge $\to$ ODBC Driver $\to$ DB & Obsolete; platform-dependent C libraries required. Removed in Java 8. \\
\textbf{Type 2} & Native-API Driver & Java App $\to$ JDBC $\to$ Vendor Native C/C++ Client $\to$ DB & Requires vendor client libraries installed on every client machine; platform-dependent. \\
\textbf{Type 3} & Network-Protocol Driver & Java App $\to$ JDBC $\to$ Middleware Gateway $\to$ Native Wire $\to$ DB & Pure Java client; routes through central middleware proxy. Complex infrastructure. \\
\textbf{Type 4} & Thin / Direct Native-Wire Driver & Java App $\to$ 100\% Pure Java Driver $\to$ DB Wire Protocol & \textbf{Modern Industry Standard.} Completely platform-independent; connects directly via TCP sockets. \\
\bottomrule
\end{tabularx}
\end{table}

\subsection{Core JDBC Interfaces and Statement Hierarchy}
\begin{itemize}
    \item \textbf{\texttt{Connection}:} Represents an active physical session with the database. Controls transaction demarcations via \texttt{setAutoCommit(false)}, \texttt{commit()}, and \texttt{rollback()}.
    \item \textbf{\texttt{Statement}:} Executes static SQL queries without parameters. Highly vulnerable to SQL injection if constructed using string concatenation.
    \item \textbf{\texttt{PreparedStatement}:} Represents a \textbf{pre-compiled} parameterized SQL statement. The query structure is compiled and optimized by the RDBMS engine \textit{before} runtime parameters are bound via positional placeholders (\texttt{?}).
\end{itemize}

\begin{quote}
\textbf{Defense-in-Depth: SQL Injection Prevention via PreparedStatement}
String concatenation dynamically alters the query's abstract syntax tree (AST):
\begin{lstlisting}[language=Java]
// VULNERABLE: Direct concatenation permits SQL injection
String query = "SELECT * FROM Users WHERE User = '" + userInput + "'";
\end{lstlisting}
If \texttt{userInput} is \texttt{' OR '1'='1}, the logical condition collapses to tautology, granting unauthorized access. 

With \texttt{PreparedStatement}, the database engine compiles the SQL command structure \textbf{prior} to receiving parameters:
\begin{lstlisting}[language=Java]
String sql = "SELECT * FROM Users WHERE Username = ? AND Status = ?";
PreparedStatement ps = connection.prepareStatement(sql);
ps.setString(1, username);
ps.setString(2, "Active");
\end{lstlisting}
The engine treats all supplied parameter inputs strictly as \textbf{literal domain data}, never as executable syntactic tokens.
\end{quote}

\subsection{Production JDBC Transaction Pattern: Banking Fund Transfer}
The following standard enterprise pattern demonstrates atomic transaction management with exception handling, rollback guarantees, and deterministic resource deallocation using Java's \texttt{try-with-resources}:

\begin{lstlisting}[language=Java]
public void transferFunds(Connection conn, long fromAcc, long toAcc, BigDecimal amount) 
        throws SQLException {
    String debitSql  = "UPDATE Accounts SET Balance = Balance - ? " +
                       "WHERE Account_ID = ? AND Status = 'Active' AND Balance >= ?";
    String creditSql = "UPDATE Accounts SET Balance = Balance + ? " +
                       "WHERE Account_ID = ? AND Status = 'Active'";

    boolean originalAutoCommit = conn.getAutoCommit();
    try {
        conn.setAutoCommit(false); // Demarcate transaction boundary

        try (PreparedStatement debitStmt = conn.prepareStatement(debitSql);
             PreparedStatement creditStmt = conn.prepareStatement(creditSql)) {

            // Step 1: Execute atomic debit on source account
            debitStmt.setBigDecimal(1, amount);
            debitStmt.setLong(2, fromAcc);
            debitStmt.setBigDecimal(3, amount);
            int debited = debitStmt.executeUpdate();
            if (debited != 1) {
                throw new SQLException("Debit failed: Insufficient balance or inactive account.");
            }

            // Step 2: Execute atomic credit on destination account
            creditStmt.setBigDecimal(1, amount);
            creditStmt.setLong(2, toAcc);
            int credited = creditStmt.executeUpdate();
            if (credited != 1) {
                throw new SQLException("Credit failed: Destination account inactive or absent.");
            }

            conn.commit(); // Persist both operations atomically (ACID Atomicity)
        }
    } catch (SQLException ex) {
        conn.rollback(); // Undo partial updates on any operational failure
        throw ex;
    } finally {
        conn.setAutoCommit(originalAutoCommit); // Restore original connection state
    }
}
\end{lstlisting}

% =============================================================================
% SECTION 30: DOCUMENT-ORIENTED NOSQL & APACHE MONGODB
% =============================================================================
\section{Document-Oriented NoSQL Systems: Apache MongoDB}
\label{sec:dbms_mongodb}

\subsection{The Paradigm Shift to Document NoSQL}
Relational databases enforce strict normalization, static schemas, and relational join constraints. In high-velocity web architectures, content management, and real-time event analytics, systems encounter semi-structured payloads, hierarchical entities, and massive horizontal write loads. \textbf{Apache MongoDB} is a leading open-source, document-oriented NoSQL database that stores data as self-describing, binary-encoded JSON structures termed \textbf{BSON (Binary JSON)}.

\begin{table}[htbp]
\centering
\small
\caption{Architectural Structural Mapping: RDBMS versus MongoDB}
\label{tab:rdbms_vs_mongodb_mapping}
\begin{tabularx}{\textwidth}{l l X}
\toprule
\textbf{Relational Concept (RDBMS)} & \textbf{MongoDB Concept} & \textbf{Architectural Description} \\
\midrule
Database & Database & Logical container grouping related data collections. \\
Table & Collection & Schema-flexible grouping of documents (no fixed column layout). \\
Row / Tuple & Document & Individual record encoded as a self-describing BSON object. \\
Column / Attribute & Field & Key-value pair within a document. \\
Primary Key & \texttt{\_id} Field & Unique identifier mandatory in every document (12-byte \texttt{ObjectId}). \\
Foreign Key / Join & Embedded Document / \texttt{\$lookup} & Denormalized hierarchical subdocument or pipeline reference lookup. \\
\bottomrule
\end{tabularx}
\end{table}

\subsection{BSON Architecture and Document Representation}
\begin{lstlisting}[language=Java]
{
  "_id": 101,
  "name": "Rahim",
  "department": "CSE",
  "cgpa": 3.75,
  "contact": {
    "email": "rahim@example.com",
    "phone": "01700000000"
  },
  "courses": ["CSE101", "CSE102", "CSE205"]
}
\end{lstlisting}
Unlike flat relational tuples that require join tables for multivalued attributes, a single MongoDB document naturally encapsulates nested key-value objects and variable-length arrays.

\subsection{MongoDB CRUD Operations}
\begin{itemize}
    \item \textbf{Create (Insert):}
\begin{lstlisting}[language=Java]
db.students.insertOne({ _id: 101, name: "Rahim", department: "CSE", cgpa: 3.75 });
db.students.insertMany([
  { _id: 102, name: "Karim", department: "EEE", cgpa: 3.60 },
  { _id: 103, name: "Nabila", department: "CSE", cgpa: 3.90 }
]);
\end{lstlisting}
    \item \textbf{Read (Query) and Comparison Operators:}
\begin{lstlisting}[language=Java]
db.students.find({ department: "CSE", cgpa: { $gt: 3.50 } });
db.students.find({ department: { $in: ["CSE", "EEE"] } });
-- Projection: 1 includes attribute, 0 suppresses it
db.students.find({ department: "CSE" }, { name: 1, cgpa: 1, _id: 0 });
\end{lstlisting}
    \item \textbf{Update (Atomic Modifiers):}
\begin{lstlisting}[language=Java]
db.students.updateOne({ _id: 101 }, { $set: { cgpa: 3.85 } });
db.students.updateMany({ department: "CSE" }, { $set: { status: "Active" } });
db.students.updateOne({ _id: 101 }, { $inc: { cgpa: 0.05 } }); -- Increment
db.students.updateOne({ _id: 101 }, { $push: { courses: "CSE301" } }); -- Array push
\end{lstlisting}
    \item \textbf{Delete:}
\begin{lstlisting}[language=Java]
db.students.deleteOne({ _id: 101 });
db.students.deleteMany({ department: "EEE" });
\end{lstlisting}
\end{itemize}

\subsection{The Aggregation Pipeline Framework}
The MongoDB Aggregation Framework processes streams of documents through multi-stage computational pipelines:
\begin{lstlisting}[language=Java]
// Compute average CGPA and total student count per department for active students
db.students.aggregate([
  { $match: { status: "Active" } },
  { $group: {
      _id: "$department",
      averageCGPA: { $avg: "$cgpa" },
      totalStudents: { $sum: 1 }
  }},
  { $sort: { averageCGPA: -1 } },
  { $limit: 5 }
]);
\end{lstlisting}

% =============================================================================
% SECTION 31: MONGODB DATA MODELING, REPLICATION & SHARDING
% =============================================================================
\section{MongoDB Data Modeling, Replica Sets, and Horizontal Sharding}
\label{sec:dbms_mongodb_scaling}

\subsection{Data Modeling: Embedding versus Referencing}
A fundamental design decision in MongoDB is choosing between \textbf{denormalized embedded subdocuments} and \textbf{normalized document references}:

\begin{table}[htbp]
\centering
\small
\caption{Architectural Comparison: Embedding versus Referencing}
\label{tab:embedding_vs_referencing}
\begin{tabularx}{\textwidth}{p{3.0cm} X X}
\toprule
\textbf{Design Criterion} & \textbf{Embedded Subdocuments} & \textbf{Document References (\texttt{DBRef} / ID)} \\
\midrule
\textbf{Structural Layout} & Related entities are stored inside the parent document. & Entities are stored in separate collections and linked via IDs. \\
\textbf{Read Performance} & Maximum: complete hierarchical record retrieved in a single I/O read. & Lower: requires application-level joins or server-side \texttt{\$lookup} pipelines. \\
\textbf{Write Atomicity} & Atomic writes guaranteed natively within single document boundaries. & Multi-document operations require distributed ACID transaction sessions. \\
\textbf{Data Redundancy} & High risk of redundancy if child entities are shared across parents. & Highly normalized; minimal redundancy. \\
\textbf{Storage Limit} & Strictly bounded by the \textbf{16 MiB BSON maximum document limit}. & Unbounded: collections scale independently. \\
\textbf{Recommended For} & $1:1$ or bounded $1:\text{Few}$ relationships (addresses, line items). & $1:N$ (unbounded) and $M:N$ relationships (users to followers). \\
\bottomrule
\end{tabularx}
\end{table}

\subsection{High Availability via Replica Sets}
A \textbf{Replica Set} is a cluster of MongoDB nodes maintaining identical datasets to provide fault tolerance and automatic failover.
\begin{itemize}
    \item \textbf{Primary Node:} The single designated leader that accepts all client write operations and records state mutations into an append-only replication log termed the \textbf{oplog (Operations Log)}.
    \item \textbf{Secondary Nodes:} Replicate and apply the primary's oplog asynchronously to maintain synchronization. Can serve read-only queries when configured with appropriate \textbf{Read Preferences} (e.g., \texttt{secondaryPreferred}).
    \item \textbf{Automatic Raft/Paxos-like Failover:} If the primary heartbeat ceases for $> 10$ seconds, eligible secondaries initiate an automated election. The node with the most up-to-date oplog is elected the new primary without manual administrative intervention.
\end{itemize}

\subsection{Horizontal Scalability via Sharding}
When datasets exceed the physical disk capacity or write throughput of a single server, MongoDB employs \textbf{Sharding} to partition collections across a distributed cluster:

\begin{figure}[htbp]
\centering
\begin{tcolorbox}[colback=codebg, colframe=primary, title=\textbf{The Three Core Components of a Sharded Cluster}]
\begin{enumerate}
    \item \textbf{Shards:} Dedicated replica sets that store physical contiguous subsets (\textbf{chunks}) of partitioned collection data.
    \item \textbf{Query Routers (\texttt{mongos}):} Stateless interface processes that intercept client requests, consult cluster metadata, and route queries directly to the relevant shard(s).
    \item \textbf{Config Servers:} A dedicated replica set that persists cluster configuration metadata and the exact mappings of chunk boundaries across shards.
\end{enumerate}
\end{tcolorbox}
\end{figure}

\begin{quote}
\textbf{Golden Examination Distinction: Replication versus Sharding}
\begin{itemize}
    \item \textbf{Replication = High Availability and Redundancy.} Every replica set node holds a \textbf{duplicate copy} of the exact same data.
    \item \textbf{Sharding = Horizontal Scalability and Capacity.} Data is \textbf{partitioned} across machines; each shard holds a unique, disjoint fraction of the total dataset.
\end{itemize}
\end{quote}



% =============================================================================
% SECTION: ADVANCED INDEXING & QUERY OPTIMIZATION
% =============================================================================
\section{Advanced Indexing, Query Optimization Engines, and SARGability}
\label{sec:dbms_advanced_indexing}

\subsection{B+ Tree Architecture versus Hash Indexing}
Most modern relational database engines deploy \textbf{B+ Trees} as their default indexing data structure:
\begin{itemize}
    \item \textbf{B+ Tree Mechanics:} A balanced search tree where all data pointers or table rows reside strictly in the \textbf{leaf nodes}. Leaf nodes are interconnected via bidirectional linked lists, enabling $O(\log N)$ point lookups as well as extremely efficient sequential $O(K)$ \textbf{range scans} without navigating back up the tree hierarchy.
    \item \textbf{Hash Indexes:} Uses a mathematical hash function $h(k)$ to map search keys directly into bucket addresses. Hash indexes deliver $O(1)$ point lookups for exact equality matches (\texttt{WHERE id = 101}), but are \textbf{fundamentally incapable of supporting range scans} (\texttt{WHERE age BETWEEN 20 AND 30}) or prefix lookups because hashing destroys natural order.
\end{itemize}

\subsection{The Cost-Based Optimizer (CBO) Pipeline}
When a declarative SQL query enters the database, the Cost-Based Optimizer determines the physical execution plan by estimating CPU cycles, memory buffers, and I/O page reads:

\begin{figure}[htbp]
\centering
\begin{tcolorbox}[colback=codebg, colframe=primary, title=\textbf{The 5-Stage Query Processing Pipeline}]
\begin{enumerate}
    \item \textbf{Parser and Validator:} Validates SQL syntax, checks table and column existence against the system catalog, and generates an internal parse tree.
    \item \textbf{Query Rewriter / Preprocessor:} Flattens subqueries, eliminates redundant conditions, unfolds views, and applies relational algebra equivalence rewrites.
    \item \textbf{Cost-Based Optimizer (CBO):} Evaluates alternate access paths (Sequential Scan vs. Index Scan), join orders, and join algorithms (Nested Loop, Hash Join, Merge Join) using statistical data distribution histograms.
    \item \textbf{Execution Engine:} Interprets the generated physical execution plan tree via the Volcano iterator model (\texttt{open()}, \texttt{next()}, \texttt{close()}).
    \item \textbf{Storage Engine and Buffer Manager:} Fetches required physical 8 KB pages into memory, reads tuples, and returns the result stream.
\end{enumerate}
\end{tcolorbox}
\end{figure}

\subsection{Execution Plan Inspection: EXPLAIN and EXPLAIN ANALYZE}
\begin{itemize}
    \item \textbf{\texttt{EXPLAIN}:} Displays the query execution plan generated by the optimizer without running the query, showing estimated startup costs, total costs, planned row counts, and projected byte widths.
    \item \textbf{\texttt{EXPLAIN ANALYZE}:} Actually executes the query in the database engine and displays \textbf{actual runtime elapsed times and true row counts} alongside the optimizer's original estimates. Essential for identifying outdated table statistics or skewed cost estimates.
\end{itemize}

\subsection{SARGable Queries (Search Argument Able)}
A predicate is defined as \textbf{SARGable} if the query optimizer can utilize an index range scan to satisfy the condition rather than resorting to an expensive full table scan.
\begin{quote}
\textbf{The Primary Rule of SARGability:} Never wrap an indexed column in an arithmetic expression or scalar function on the left-hand side of a comparison operator!
\end{quote}

\begin{table}[htbp]
\centering
\small
\caption{Non-SARGable versus SARGable Refactoring Matrix}
\label{tab:sargable_queries}
\begin{tabularx}{\textwidth}{X X}
\toprule
\textbf{Non-SARGable (Forces Full Table Scan)} & \textbf{SARGable Refactored (Permits B-Tree Range Scan)} \\
\midrule
\texttt{WHERE YEAR(Join\_Date) = 2025} & \texttt{WHERE Join\_Date >= '2025-01-01' AND Join\_Date < '2026-01-01'} \\
\texttt{WHERE Salary * 1.10 > 55000.00} & \texttt{WHERE Salary > 55000.00 / 1.10} \\
\texttt{WHERE SUBSTRING(Name, 1, 1) = 'R'} & \texttt{WHERE Name LIKE 'R\%'} \\
\texttt{WHERE UPPER(Email) = 'USER@EXAMPLE.COM'} & Functional index \texttt{ON Users(UPPER(Email))} or normalized storage. \\
\bottomrule
\end{tabularx}
\end{table}

% =============================================================================
% SECTION: DATABASE STORAGE ARCHITECTURE & BUFFER MANAGEMENT
% =============================================================================
\section{Database Storage Architecture, Slotted Pages, and Buffer Management}
\label{sec:dbms_storage_architecture}

\subsection{Physical Storage Hierarchy: Pages and Blocks}
A database engine does not read or write individual table rows directly to raw disk storage. The foundational unit of hardware disk I/O and RAM caching is the \textbf{Database Page} (typically \textbf{8 KiB} in PostgreSQL and Microsoft SQL Server; 16 KiB in MySQL InnoDB).
\begin{itemize}
    \item \textbf{Slotted-Page Architecture:} To support variable-length records, a database page divides its physical structure into:
    \begin{enumerate}
        \item \textbf{Page Header:} Contains LSN (Log Sequence Number) for WAL recovery, transaction visibility metadata, and free space offsets.
        \item \textbf{Slot Directory (Line Pointers):} Grows downwards from the top of the page, storing pointer offsets to each physical record.
        \item \textbf{Free Space:} Unallocated page bytes between the slot directory and the record heap.
        \item \textbf{Tuple Data Storage:} Grows upwards from the bottom of the page, storing actual attribute byte arrays.
    \end{enumerate}
    This slotted architecture guarantees that a row's physical address (\textbf{Tuple ID / TID / RowID}) remains permanent even if data within the page is compacted or rearranged.
\end{itemize}

\subsection{The Buffer Pool and Dirty Page Lifecycle}
To bridge the nanosecond latency of CPU registers and RAM with the millisecond latency of magnetic disks or solid-state drives, the RDBMS manages a dedicated region of system memory termed the \textbf{Buffer Pool (Buffer Cache)}:
\begin{figure}[htbp]
\centering
\small
\begin{tabularx}{\textwidth}{p{3.2cm} X}
\toprule
\textbf{Buffer Pool Component} & \textbf{Architectural Function} \\
\midrule
\textbf{Page Frame Buffer} & Array of in-memory page slots (each matching the disk page size, e.g., 8 KB). \\
\textbf{Buffer Hash Table} & Maps disk page identifiers (\texttt{File\_ID, Block\_Number}) to memory frame addresses in $O(1)$. \\
\textbf{Dirty Page} & A page modified in the buffer pool by an \texttt{UPDATE} or \texttt{INSERT} but not yet written to disk. \\
\textbf{Eviction Policy (LRU/Clock)} & Determines which clean or flushed page frame is evicted when free memory is exhausted. \\
\textbf{Checkpointer Daemon} & Background worker that periodically flushes dirty buffer pages to persistent disk. \\
\bottomrule
\end{tabularx}
\end{figure}

\begin{quote}
\textbf{Write-Ahead Logging Invariant for Buffer Eviction (WAL Protocol)}
A dirty page resident in the buffer pool \textbf{can never be written to persistent disk storage} until the corresponding log records describing all uncommitted mutations to that page have first been synchronously flushed to the non-volatile transaction log (\texttt{WAL} / \texttt{redo log}). This guarantees that if power fails mid-write, the engine can execute crash recovery via \textbf{Redo} and \textbf{Undo} logs.
\end{quote}

% =============================================================================
% SECTION: ADVANCED CONCURRENCY: MVCC & 2PC
% =============================================================================
\section{Advanced Concurrency: Multi-Version Concurrency Control and Two-Phase Commit}
\label{sec:dbms_mvcc_2pc}

\subsection{Multi-Version Concurrency Control (MVCC)}
Traditional lock-based concurrency forces read operations to acquire shared ($S$) locks, causing readers to block writers and writers to block readers. \textbf{Multi-Version Concurrency Control (MVCC)} resolves this bottleneck under a central operational principle:
\begin{center}
\textbf{Readers never block writers, and writers never block readers.}
\end{center}

\textbf{How MVCC Works in PostgreSQL and InnoDB:}
\begin{itemize}
    \item When an \texttt{UPDATE} occurs, the engine does not physically overwrite the old row in-place. Instead, it marks the existing tuple version as deleted and inserts a \textbf{new version} of the row with an updated transaction creation identifier (\texttt{xmin}/\texttt{xmax} in PostgreSQL; Undo segment roll-pointers in InnoDB).
    \item When a query begins, the engine establishes a \textbf{Snapshot} consisting of the set of all transaction IDs committed prior to the query's start time. The query reads only those tuple versions visible within its snapshot, ignoring uncommitted or newer mutations.
    \item \textbf{Vacuuming / Garbage Collection:} Old, dead row versions obsolete to all active transaction snapshots must be periodically pruned by background maintenance daemons (e.g., PostgreSQL \texttt{VACUUM}) to reclaim page storage.
\end{itemize}

\subsection{Pessimistic versus Optimistic Concurrency Control}
\begin{itemize}
    \item \textbf{Pessimistic Concurrency Control:} Assumes high data collision frequency. Acquires exclusive locks ($X$) upfront before mutating data (e.g., \texttt{SELECT ... FOR UPDATE}). Prevents conflicts at the cost of reduced throughput and potential deadlocks.
    \item \textbf{Optimistic Concurrency Control (OCC):} Assumes collisions are rare. Transactions execute reads and modifications in private memory without locking. During the commit phase, the transaction verifies whether another transaction has altered the record (commonly checked via a numeric \texttt{Version} column). If an inconsistent state is detected, the transaction aborts and retries:
\begin{lstlisting}[language=SQL]
UPDATE Products SET Stock = 8, Version = Version + 1
WHERE Product_ID = 101 AND Version = 5;
\end{lstlisting}
\end{itemize}

\subsection{Two-Phase Commit (2PC) in Distributed Systems}
When a logical transaction spans multiple autonomous, network-separated database nodes (e.g., debiting a core banking server in Dhaka and crediting a remittance server in Singapore), standard local ACID transactions cannot ensure global atomicity. The \textbf{Two-Phase Commit (2PC)} protocol coordinates participants through a central \textbf{Coordinator}:

\begin{figure}[htbp]
\centering
\begin{tcolorbox}[colback=codebg, colframe=primary, title=\textbf{The Two Phases of Distributed 2PC Coordination}]
\begin{enumerate}
    \item \textbf{Phase 1: Prepare Phase}
    \begin{itemize}
        \item The Coordinator sends a \texttt{PREPARE} message to all participant nodes.
        \item Each participant executes local operations, flushes transaction mutations to its local redo log, and locks resources.
        \item If successful, the node votes \texttt{VOTE\_COMMIT}; if a local failure occurs, it votes \texttt{VOTE\_ABORT}.
    \end{itemize}
    \item \textbf{Phase 2: Commit or Abort Phase}
    \begin{itemize}
        \item If \textbf{all participants vote \texttt{VOTE\_COMMIT}}, the Coordinator writes a \texttt{GLOBAL\_COMMIT} log record and broadcasts \texttt{COMMIT} to all nodes. Each node commits locally and releases locks.
        \item If \textbf{any participant votes \texttt{VOTE\_ABORT}} (or a communication timeout occurs), the Coordinator broadcasts \texttt{GLOBAL\_ABORT}, causing all nodes to roll back.
    \end{itemize}
\end{enumerate}
\end{tcolorbox}
\end{figure}

\begin{quote}
\textbf{Vital Recruitment Distinction: 2PC versus 2PL}
\begin{itemize}
    \item \textbf{2PC (Two-Phase Commit):} A \textbf{distributed consensus and recovery protocol} ensuring global atomicity across multiple distinct database servers via \textit{Prepare} and \textit{Commit} phases.
    \item \textbf{2PL (Two-Phase Locking):} A \textbf{local concurrency-control locking protocol} ensuring serializability within an individual RDBMS via \textit{Growing} and \textit{Shrinking} phases.
\end{itemize}
\end{quote}

% =============================================================================
% SECTION: DATABASE PARTITIONING, REPLICATION & THE CAP THEOREM
% =============================================================================
\section{Database Partitioning, High-Availability Replication, and the CAP Theorem}
\label{sec:dbms_partitioning_cap}

\subsection{Database Partitioning Strategies}
\textbf{Partitioning} divides a very large database table into smaller, physically independent segments while preserving a unified logical relational interface:
\begin{itemize}
    \item \textbf{Horizontal Partitioning (Sharding):} Divides records by \textbf{rows}. Common strategies include:
    \begin{enumerate}
        \item \textbf{Range Partitioning:} Tuples are segregated by contiguous continuous bounds (e.g., partitioning transaction ledgers by fiscal year: \texttt{Orders\_2024}, \texttt{Orders\_2025}).
        \item \textbf{List Partitioning:} Tuples are segregated by explicit discrete sets (e.g., partitioning by division: \texttt{'Dhaka'}, \texttt{'Chattogram'}).
        \item \textbf{Hash Partitioning:} Applies a uniform hash function to the partition key modulo $N$ ($h(\text{key}) \pmod N$), distributing tuples evenly across physical drives.
    \end{enumerate}
    \item \textbf{Vertical Partitioning:} Divides records by \textbf{columns}. Moves rarely accessed or sensitive attributes (e.g., customer medical history, profile photographs) into auxiliary tables linked by the primary key to maximize buffer pool density for high-frequency queries.
\end{itemize}

\subsection{Database Replication: Synchronous versus Asynchronous}
\begin{table}[htbp]
\centering
\small
\caption{Replication Topology Comparison}
\label{tab:replication_comparison}
\begin{tabularx}{\textwidth}{p{3.2cm} X X}
\toprule
\textbf{Replication Mode} & \textbf{Commit Acknowledgment Requirement} & \textbf{Engineering Trade-Offs} \\
\midrule
\textbf{Synchronous} & Primary node waits for confirmed write acknowledgment from replica(s) before returning \texttt{COMMIT} to client. & Guarantees zero data loss (RPO = 0) and immediate consistency; increases client write latency and risks downtime if network partitions. \\
\textbf{Asynchronous} & Primary node commits locally and returns success immediately; stream mutations to replicas in background. & Minimizes client write latency and survives replica outages; risks data loss (replication lag) if primary crashes before logs propagate. \\
\bottomrule
\end{tabularx}
\end{table}

\subsection{The CAP Theorem (Brewer's Theorem)}
Formulated by Eric Brewer and formally proven by Seth Gilbert and Nancy Lynch, the \textbf{CAP Theorem} governs distributed state systems:
\begin{itemize}
    \item \textbf{Consistency (C):} Every read receives the most recent committed write or an explicit error (Linearizability).
    \item \textbf{Availability (A):} Every non-failing node returns a non-error response for every request (without guaranteeing latest data).
    \item \textbf{Partition Tolerance (P):} The system continues operational service despite arbitrary network packet drops or network partitions separating cluster nodes.
\end{itemize}

\begin{quote}
\textbf{The Modern CAP Reality Formulation:}
Because physical network communication over Ethernet cables and routers can never be guaranteed to be 100\% reliable, \textbf{Partition Tolerance ($P$) is a mandatory property of any physical distributed system}. Therefore, when an inevitable network partition ($P$) occurs, a distributed database must choose between:
\begin{itemize}
    \item \textbf{CP Systems (Consistency over Availability):} Reject or block writes on partitioned sub-clusters to avoid conflicting data states (e.g., HBase, MongoDB primary writes, Spanner).
    \item \textbf{AP Systems (Availability over Consistency):} Allow partitioned nodes to accept writes locally, tolerating temporary replica drift and resolving conflicts later via eventual consistency (e.g., Apache Cassandra, DynamoDB).
\end{itemize}
\end{quote}

% =============================================================================
% SECTION: DATA WAREHOUSING, OLTP/OLAP, ETL & DIMENSIONAL MODELING
% =============================================================================
\section{Data Warehousing, OLTP versus OLAP, ETL/ELT, and Dimensional Modeling}
\label{sec:dbms_oltp_olap}

\subsection{OLTP versus OLAP: Operational Systems versus Analytics}
\begin{table}[htbp]
\centering
\small
\caption{Architectural Comparison: OLTP versus OLAP Systems}
\label{tab:oltp_vs_olap}
\begin{tabularx}{\textwidth}{p{3.2cm} X X}
\toprule
\textbf{Design Criterion} & \textbf{OLTP (Online Transaction Processing)} & \textbf{OLAP (Online Analytical Processing)} \\
\midrule
\textbf{Primary Purpose} & Execute day-to-day operational business transactions (e.g., banking fund transfers, retail orders). & Consolidate historical enterprise data to extract business intelligence and strategic executive trends. \\
\textbf{Workload Characteristics} & Massive volumes of short, low-latency \texttt{INSERT}, \texttt{UPDATE}, \texttt{DELETE} statements. & Periodic read-heavy batch queries computing complex aggregations (\texttt{SUM}, \texttt{AVG}) across millions of rows. \\
\textbf{Database Schema} & Highly normalized (3NF/BCNF) to eliminate update anomalies and optimize single-row writes. & Denormalized dimensional models (\textbf{Star Schema}, \textbf{Snowflake Schema}) to minimize join latency. \\
\textbf{Data Horizon} & Current, real-time snapshot; historical data purged or archived. & Multi-year historical time-series data retaining complete mutation histories. \\
\textbf{User Community} & End-user customers, frontline bank tellers, transaction microservices. & Data scientists, financial auditors, executive management. \\
\bottomrule
\end{tabularx}
\end{table}

\subsection{The Four Inman Characteristics of a Data Warehouse}
W. H. Inman defined a \textbf{Data Warehouse} as a central, integrated corporate analytical repository exhibiting four fundamental characteristics:
\begin{enumerate}
    \item \textbf{Subject-Oriented:} Structured around high-level enterprise operational entities (Sales, Claims, Balances) rather than application software functions.
    \item \textbf{Integrated:} Resolves naming conflicts, unit mismatches, and encoding variations from disparate transactional databases into a singular unified standard.
    \item \textbf{Time-Variant:} Every record is explicitly tied to an immutable historical timestamp, enabling trend analysis over multi-year windows.
    \item \textbf{Non-Volatile:} Once loaded, historical records are immutable; data is never updated in-place, only appended and read.
\end{enumerate}

\subsection{Data Integration Pipelines: ETL versus ELT}
\begin{itemize}
    \item \textbf{ETL (Extract, Transform, Load):} Source data is extracted from OLTP databases, transformed in a specialized staging server (cleaning, deduplication, type casting), and loaded into the relational data warehouse. Standard in legacy on-premise architectures.
    \item \textbf{ELT (Extract, Load, Transform):} Raw, untransformed data is dumped directly into modern scalable cloud warehouses/data lakes (e.g., BigQuery, Snowflake), with transformations executed in-place via massive parallel SQL compute engines.
\end{itemize}

\subsection{Dimensional Modeling: Star Schema versus Snowflake Schema}
\begin{figure}[htbp]
\centering
\small
\begin{tabularx}{\textwidth}{l p{3.2cm} X}
\toprule
\textbf{Dimensional Entity} & \textbf{Primary Contents} & \textbf{Banking Case Study Example} \\
\midrule
\textbf{Fact Table} & Numerical additive business measures and foreign keys pointing to dimension tables. & \texttt{Fact\_Transactions}(Date\_Key, Branch\_Key, Acc\_Key, Amount, Fee). \\
\textbf{Dimension Table} & Rich contextual textual attributes providing analytical filtering and slicing criteria. & \texttt{Dim\_Branch}(Branch\_Key, City, Division, Regional\_Manager). \\
\bottomrule
\end{tabularx}
\end{figure}

\begin{quote}
\textbf{Star Schema versus Snowflake Schema Structural Difference:}
\begin{itemize}
    \item \textbf{Star Schema:} Dimension tables are completely \textbf{denormalized}. Simple queries require fewer joins, maximizing OLAP scan speed at the expense of dimension redundancy.
    \item \textbf{Snowflake Schema:} Dimension tables are \textbf{normalized} into hierarchical sub-tables (e.g., splitting \texttt{Dim\_Product} into \texttt{Category} and \texttt{Brand}). Eliminates dimension redundancy but increases query complexity and join costs.
\end{itemize}
\end{quote}

% =============================================================================
% SECTION: DATABASE SECURITY, AUDITING & POINT-IN-TIME RECOVERY
% =============================================================================
\section{Database Security Administration, Role-Based Access, and Point-in-Time Recovery}
\label{sec:dbms_security_pitr}

\subsection{Role-Based Access Control (RBAC): GRANT and REVOKE}
In production database governance, administrative security mandates the Principle of Least Privilege:
\begin{lstlisting}[language=SQL]
-- Create functional role for reporting analysts
CREATE ROLE financial_auditor;
GRANT SELECT ON Accounts, Bank_Transactions TO financial_auditor;

-- Assign role to individual database user
GRANT financial_auditor TO user_tania;

-- Revoke permissions upon employee transfer
REVOKE financial_auditor FROM user_tania;
\end{lstlisting}

\subsection{Database Auditing Standards}
Auditing records immutable security logs tracking who accessed which relational objects, what commands were executed, the source IP address, and whether the attempt succeeded. Essential for regulatory compliance (e.g., PCI-DSS in banking, HIPAA in healthcare) to detect unauthorized privilege escalations and data exfiltration.

\subsection{Point-in-Time Recovery (PITR)}
\textbf{Point-in-Time Recovery (PITR)} is an advanced database restoration procedure that reconstructs the database state to an exact historic second:
\begin{equation}
\text{State at } T_{\text{target}} = \text{Full Base Backup} + \text{WAL Redo Log Replay up to } T_{\text{target}}
\end{equation}
If a developer accidentally executes an unconstrained \texttt{DELETE FROM Customers} at 14:32:15, the DBA restores the midnight full backup and replays the continuous Write-Ahead Logs up to 14:32:14, recovering all valid data without losing 14 hours of production transactions.



% =============================================================================
% SECTION 32: ENTERPRISE CORE BANKING DATABASE ARCHITECTURE
% =============================================================================
\section{Enterprise Core Banking Database Architecture and Relational Schema}
\label{sec:dbms_banking_schema}

\subsection{Architectural Requirements for Financial Ledger Systems}
Financial and banking database systems represent the highest tier of engineering rigor in enterprise data management. A banking relational database must satisfy non-negotiable invariants:
\begin{itemize}
    \item \textbf{Mathematical Balance Conservation:} Ledger debit and credit operations must balance with zero precision loss.
    \item \textbf{Strict ACID Transaction Isolation:} Prevents double-spending, phantom updates, and race conditions during concurrent ATM/online fund transfers.
    \item \textbf{Referential Auditing and Immutability:} Prohibits physical tuple deletions on historical transaction ledgers.
    \item \textbf{Multi-Entity Relational Topology:} Models complex one-to-many ($1:N$) and many-to-many ($M:N$) ownership structures (e.g., joint accounts).
\end{itemize}

\subsection{Core Banking Relational Schema Design}
The relational banking system comprises five interconnected core entities:
\begin{enumerate}
    \item \textbf{\texttt{Branches}:} Stores physical bank branch locations, clearing codes, and city jurisdictions.
    \item \textbf{\texttt{Customers}:} Captures personal identifying information, tax/national IDs, and contact vectors.
    \item \textbf{\texttt{Accounts}:} Represents ledger balance entities holding funds under specific deposit categories.
    \item \textbf{\texttt{Account\_Holders}:} A junction relation modeling the \textbf{Many-to-Many ($M:N$)} relationship between Customers and Accounts (supporting joint accounts and individual multi-account holdings).
    \item \textbf{\texttt{Bank\_Transactions}:} The central financial ledger recording all debit, credit, and transfer movements between accounts.
\end{enumerate}

\subsection{Production-Grade SQL DDL Schema Implementation}
The following standard SQL (PostgreSQL/ANSI compliant) DDL demonstrates domain constraints, foreign keys with restrict semantics, and check constraints:

\begin{lstlisting}[language=SQL]
-- 1. Branches Table
CREATE TABLE Branches (
    Branch_ID   INT PRIMARY KEY,
    Branch_Name VARCHAR(100) NOT NULL,
    Address     VARCHAR(255) NOT NULL,
    City        VARCHAR(100) NOT NULL,
    Phone       VARCHAR(20)
);

-- 2. Customers Table
CREATE TABLE Customers (
    Customer_ID   BIGINT PRIMARY KEY,
    Full_Name     VARCHAR(150) NOT NULL,
    Date_of_Birth DATE,
    Email         VARCHAR(150) UNIQUE,
    Phone         VARCHAR(20),
    Address       VARCHAR(255),
    Created_At    TIMESTAMP DEFAULT CURRENT_TIMESTAMP
);

-- 3. Accounts Table
CREATE TABLE Accounts (
    Account_ID     BIGINT PRIMARY KEY,
    Account_Number VARCHAR(30) NOT NULL UNIQUE,
    Branch_ID      INT NOT NULL,
    Account_Type   VARCHAR(20) NOT NULL CHECK (Account_Type IN ('Savings', 'Current', 'Fixed Deposit')),
    Balance        DECIMAL(18, 2) NOT NULL DEFAULT 0.00 CHECK (Balance >= 0.00),
    Status         VARCHAR(20) NOT NULL DEFAULT 'Active' CHECK (Status IN ('Active', 'Inactive', 'Closed', 'Frozen')),
    Created_At     TIMESTAMP DEFAULT CURRENT_TIMESTAMP,
    FOREIGN KEY (Branch_ID) REFERENCES Branches(Branch_ID) ON DELETE RESTRICT
);

-- 4. Account_Holders Junction Table (M:N Relationship)
CREATE TABLE Account_Holders (
    Customer_ID BIGINT NOT NULL,
    Account_ID  BIGINT NOT NULL,
    Holder_Role VARCHAR(30) DEFAULT 'Primary' CHECK (Holder_Role IN ('Primary', 'Joint', 'Beneficiary')),
    PRIMARY KEY (Customer_ID, Account_ID),
    FOREIGN KEY (Customer_ID) REFERENCES Customers(Customer_ID) ON DELETE RESTRICT,
    FOREIGN KEY (Account_ID)  REFERENCES Accounts(Account_ID)   ON DELETE RESTRICT
);

-- 5. Bank_Transactions Ledger Table
CREATE TABLE Bank_Transactions (
    Transaction_ID   BIGINT PRIMARY KEY,
    From_Account_ID  BIGINT,
    To_Account_ID    BIGINT,
    Amount           DECIMAL(18, 2) NOT NULL CHECK (Amount > 0.00),
    Transaction_Type VARCHAR(30) NOT NULL CHECK (Transaction_Type IN ('Deposit', 'Withdrawal', 'Transfer')),
    Status           VARCHAR(20) NOT NULL CHECK (Status IN ('Pending', 'Completed', 'Failed', 'Reversed')),
    Created_At       TIMESTAMP DEFAULT CURRENT_TIMESTAMP,
    FOREIGN KEY (From_Account_ID) REFERENCES Accounts(Account_ID) ON DELETE RESTRICT,
    FOREIGN KEY (To_Account_ID)   REFERENCES Accounts(Account_ID) ON DELETE RESTRICT,
    CHECK (From_Account_ID IS NOT NULL OR To_Account_ID IS NOT NULL),
    CHECK (From_Account_ID IS NULL OR To_Account_ID IS NULL OR From_Account_ID <> To_Account_ID)
);
\end{lstlisting}

% =============================================================================
% SECTION 33: HIGH-YIELD CORE BANKING SQL QUERIES
% =============================================================================
\section{High-Yield Core Banking SQL Analytical and Operational Queries}
\label{sec:dbms_banking_queries}

\subsection{Operational and Management Reporting Queries}
The following 12 production queries represent standard tasks tested in competitive government recruitment examinations for Assistant Programmer, Senior Officer (IT), and System Analyst posts:

\begin{enumerate}[label=\textbf{Query \arabic*:}]
    \item \textbf{Retrieve All Registered Customers:}
\begin{lstlisting}[language=SQL]
SELECT * FROM Customers;
\end{lstlisting}

    \item \textbf{Retrieve All Currently Active Bank Accounts:}
\begin{lstlisting}[language=SQL]
SELECT Account_ID, Account_Number, Account_Type, Balance 
FROM Accounts WHERE Status = 'Active';
\end{lstlisting}

    \item \textbf{Retrieve Comprehensive Customer Account Holdings via 3-Table Join:}
\begin{lstlisting}[language=SQL]
SELECT c.Customer_ID, c.Full_Name, a.Account_Number, a.Account_Type, a.Balance
FROM Customers c
INNER JOIN Account_Holders ah ON c.Customer_ID = ah.Customer_ID
INNER JOIN Accounts a         ON ah.Account_ID = a.Account_ID;
\end{lstlisting}

    \item \textbf{Identify High-Net-Worth Customers with Balance Exceeding BDT 100,000:}
\begin{lstlisting}[language=SQL]
SELECT DISTINCT c.Customer_ID, c.Full_Name
FROM Customers c
INNER JOIN Account_Holders ah ON c.Customer_ID = ah.Customer_ID
INNER JOIN Accounts a         ON ah.Account_ID = a.Account_ID
WHERE a.Balance > 100000.00;
\end{lstlisting}

    \item \textbf{Count Total Registered Accounts per Branch:}
\begin{lstlisting}[language=SQL]
SELECT Branch_ID, COUNT(*) AS Total_Accounts
FROM Accounts GROUP BY Branch_ID;
\end{lstlisting}

    \item \textbf{Calculate Aggregate Deposit Balance Managed per Branch:}
\begin{lstlisting}[language=SQL]
SELECT b.Branch_ID, b.Branch_Name, SUM(a.Balance) AS Total_Branch_Balance
FROM Branches b
INNER JOIN Accounts a ON b.Branch_ID = a.Branch_ID
GROUP BY b.Branch_ID, b.Branch_Name;
\end{lstlisting}

    \item \textbf{Find Branches Managing Total Deposits Exceeding BDT 10,000,000 (HAVING):}
\begin{lstlisting}[language=SQL]
SELECT b.Branch_ID, b.Branch_Name, SUM(a.Balance) AS Total_Branch_Balance
FROM Branches b
INNER JOIN Accounts a ON b.Branch_ID = a.Branch_ID
GROUP BY b.Branch_ID, b.Branch_Name
HAVING SUM(a.Balance) > 10000000.00;
\end{lstlisting}

    \item \textbf{Extract Complete Chronological Ledger History for Account 1001:}
\begin{lstlisting}[language=SQL]
SELECT Transaction_ID, Transaction_Type, From_Account_ID, To_Account_ID, Amount, Status, Created_At
FROM Bank_Transactions
WHERE From_Account_ID = 1001 OR To_Account_ID = 1001
ORDER BY Created_At DESC;
\end{lstlisting}

    \item \textbf{Calculate Grand Total of All Successfully Completed Fund Transfers:}
\begin{lstlisting}[language=SQL]
SELECT SUM(Amount) AS Total_Transferred_Capital
FROM Bank_Transactions
WHERE Transaction_Type = 'Transfer' AND Status = 'Completed';
\end{lstlisting}

    \item \textbf{Find Accounts with Balance Above the Global Enterprise Average:}
\begin{lstlisting}[language=SQL]
SELECT Account_ID, Account_Number, Balance
FROM Accounts
WHERE Balance > (SELECT AVG(Balance) FROM Accounts);
\end{lstlisting}

    \item \textbf{Identify Customers with No Associated Bank Accounts (NOT EXISTS):}
\begin{lstlisting}[language=SQL]
SELECT c.Customer_ID, c.Full_Name
FROM Customers c
WHERE NOT EXISTS (
    SELECT 1 FROM Account_Holders ah WHERE ah.Customer_ID = c.Customer_ID
);
\end{lstlisting}

    \item \textbf{Find the Highest-Balance Account within Each Individual Branch:}
\begin{lstlisting}[language=SQL]
SELECT a.Account_ID, a.Account_Number, a.Branch_ID, a.Balance
FROM Accounts a
WHERE a.Balance = (
    SELECT MAX(a2.Balance)
    FROM Accounts a2
    WHERE a2.Branch_ID = a.Branch_ID
);
\end{lstlisting}
\end{enumerate}

% =============================================================================
% SECTION 34: BANK TRANSFER TRANSACTION & ACID IMPLEMENTATION
% =============================================================================
\section{Banking Fund Transfer Case Study: ACID Transaction Semantics}
\label{sec:dbms_banking_acid}

\subsection{The Banking Transfer Paradigm}
Consider a fund transfer of \textbf{BDT 10,000} from source account $A$ (Account ID 1001, initial balance BDT 50,000) to destination account $B$ (Account ID 1002, initial balance BDT 30,000). The total financial capital across both accounts before transfer is:
\begin{equation}
\text{Total Capital} = 50,000 + 30,000 = 80,000 \text{ BDT}
\end{equation}

\subsection{SQL Multi-Statement Atomic Transaction Execution}
\begin{lstlisting}[language=SQL]
BEGIN TRANSACTION;

-- Step 1: Deduct transfer amount from source account with balance guard
UPDATE Accounts
SET Balance = Balance - 10000.00
WHERE Account_ID = 1001 AND Status = 'Active' AND Balance >= 10000.00;

-- Application checks: Exactly 1 row updated? If not, abort!

-- Step 2: Credit transfer amount to destination account
UPDATE Accounts
SET Balance = Balance + 10000.00
WHERE Account_ID = 1002 AND Status = 'Active';

-- Step 3: Insert immutable ledger audit transaction
INSERT INTO Bank_Transactions (
    Transaction_ID, From_Account_ID, To_Account_ID, Amount, Transaction_Type, Status
) VALUES (
    9001, 1001, 1002, 10000.00, 'Transfer', 'Completed'
);

-- If all steps succeed:
COMMIT;
-- If any step fails or violates integrity:
-- ROLLBACK;
\end{lstlisting}

\subsection{Post-Transfer Balance Invariant Verification}
\begin{equation}
\text{Source Balance} = 50,000 - 10,000 = 40,000 \text{ BDT}
\end{equation}
\begin{equation}
\text{Destination Balance} = 30,000 + 10,000 = 40,000 \text{ BDT}
\end{equation}
\begin{equation}
\text{Conserved Capital} = 40,000 + 40,000 = \mathbf{80,000 \text{ BDT}}
\end{equation}

\begin{table}[htbp]
\centering
\small
\caption{ACID Property Implementation in Banking Fund Transfer}
\label{tab:banking_acid_matrix}
\begin{tabularx}{\textwidth}{l X X}
\toprule
\textbf{Property} & \textbf{Banking Invariant Enforced} & \textbf{Underlying RDBMS Mechanism} \\
\midrule
\textbf{Atomicity} & The debit, credit, and ledger log either all commit together or all roll back. No money is deducted without being credited. & Undo logging in Write-Ahead Log (WAL); automatic rollback on failure. \\
\textbf{Consistency} & Account balance cannot drop below zero (\texttt{Balance >= 0.00}); global capital sum is preserved across ledger transitions. & Relational schema constraints, check predicates, and business invariants. \\
\textbf{Isolation} & Concurrent ATM withdrawals or online transfers cannot interleave to produce double-spending or lost updates. & Two-Phase Locking (Strict 2PL), Row-level Exclusive ($X$) locks, Snapshot Isolation. \\
\textbf{Durability} & Once \texttt{COMMIT} acknowledges success, the updated balances survive unexpected power loss or server operating system crash. & Synchronous WAL flush (\texttt{fsync}) to non-volatile disk before returning success. \\
\bottomrule
\end{tabularx}
\end{table}

% =============================================================================
% SECTION 35: BIG DATA & APACHE HADOOP HDFS ARCHITECTURE
% =============================================================================
\section{Big Data Frameworks and Distributed Storage: Apache Hadoop HDFS}
\label{sec:dbms_hadoop_hdfs}

\subsection{The Big Data Explosion and the 3 Vs}
\textbf{Big Data} refers to datasets whose scale, velocity of generation, and structural heterogeneity exceed the storage and processing limits of traditional monolithic relational database management systems.

\begin{figure}[htbp]
\centering
\small
\begin{tabularx}{\textwidth}{l p{2.8cm} X}
\toprule
\textbf{Dimension} & \textbf{Core Characteristic} & \textbf{Enterprise Example} \\
\midrule
\textbf{Volume} & Massive scale of data generation & Petabytes of global credit card transaction logs, satellite imagery. \\
\textbf{Velocity} & High arrival frequency & Thousands of financial stock market ticks per second; IoT sensor feeds. \\
\textbf{Variety} & Heterogeneous data formats & Structured SQL tables, semi-structured JSON/XML, unstructured text, video. \\
\textbf{Veracity} & Data trustworthiness & Filtering noise, fraudulent transactions, and sensor drift in telemetry. \\
\textbf{Value} & Business utility & Extracting predictive risk models and customer churn insights. \\
\bottomrule
\end{tabularx}
\end{figure}

\subsection{Apache Hadoop Architectural Philosophy}
\textbf{Apache Hadoop} is an open-source distributed storage and processing framework designed to scale across thousands of commodity hardware nodes. It is governed by a fundamental computing axiom:
\begin{quote}
\textbf{Data Locality Principle:} \textit{Moving computation to data is substantially cheaper than moving massive data to computation.} Instead of streaming terabytes of data over network switches to a central computation node, Hadoop routes compiled computation tasks (Map jobs) directly to the physical storage nodes hosting the data blocks.
\end{quote}

\subsection{The Four Core Modules of Apache Hadoop}
\begin{enumerate}
    \item \textbf{Hadoop Common:} Foundational Java utilities, configuration parsers, and RPC abstractions.
    \item \textbf{Hadoop Distributed File System (HDFS):} Scalable, fault-tolerant distributed filesystem.
    \item \textbf{YARN (Yet Another Resource Negotiator):} Central cluster resource scheduling and job coordination.
    \item \textbf{Hadoop MapReduce:} Parallel distributed computational programming framework.
\end{enumerate}

\subsection{HDFS Master-Worker Storage Architecture}
HDFS implements an asymmetrical master-worker topology:
\begin{itemize}
    \item \textbf{NameNode (Master):} Maintains the filesystem directory tree, file-to-block mappings, and DataNode block locations entirely in RAM for fast metadata lookup. Persists metadata snapshots (\texttt{FsImage}) and transaction mutation logs (\texttt{EditLog}) to local disk.
    \item \textbf{DataNodes (Workers):} Store physical raw block files on local disk drives. Execute read and write block requests from clients and send periodic \textbf{Block Reports} and \textbf{Heartbeats} (every 3 seconds) to the NameNode.
\end{itemize}

\subsection{HDFS Block Splitting and Sizing Mathematics}
Files in HDFS are split into large, contiguous chunks called \textbf{blocks} (default 128 MiB or 256 MiB). 
\begin{equation}
\text{Block Count} = \left\lceil \frac{\text{File Size}}{\text{Block Size}} \right\rceil
\end{equation}
\textbf{Why Are HDFS Blocks So Large?} Disk seek latency on commodity HDDs requires $\approx 10\text{ ms}$. At an average disk sequential transfer rate of $100\text{ MB/s}$, transferring 128 MB requires $\approx 1.28\text{ s}$. The seek time ($10\text{ ms}$) represents $< 1\%$ of the total transfer time, maximizing sequential I/O throughput.

\subsection{HDFS Block Replication and Fault Tolerance}
By default, HDFS maintains a \textbf{Replication Factor of 3}:
\begin{equation}
\text{Total Storage Required} = \text{File Size} \times \text{Replication Factor}
\end{equation}

\textbf{The Rack Awareness Placement Policy:}
\begin{itemize}
    \item \textbf{Replica 1:} Placed on the local DataNode within the client's own rack.
    \item \textbf{Replica 2:} Placed on a different DataNode within the same local rack.
    \item \textbf{Replica 3:} Placed on a DataNode located in an \textbf{entirely different remote rack}.
\end{itemize}
This policy strikes an optimal balance between write bandwidth (only one inter-rack network hop) and disaster resilience (surviving entire rack power/switch failures).

\subsection{Secondary NameNode versus Standby NameNode}
\begin{quote}
\textbf{Critical Examination Trap: Secondary NameNode Is Not a Failover Backup}
\begin{itemize}
    \item \textbf{Secondary NameNode:} Does \textbf{not} provide high availability or act as a warm standby. It is merely a \textbf{checkpointing helper} that periodically downloads the \texttt{FsImage} and \texttt{EditLog} from the NameNode, merges them into a clean snapshot, and uploads the compacted \texttt{FsImage} back to prevent the NameNode's \texttt{EditLog} from growing infinitely.
    \item \textbf{Standby NameNode:} The true high-availability (HA) standby in modern Hadoop. Synchronized via Quorum Journal Manager (QJM) and orchestrated by Apache ZooKeeper, it automatically takes over active cluster leadership if the primary NameNode fails.
\end{itemize}
\end{quote}

% =============================================================================
% SECTION 36: DISTRIBUTED ENGINES & HADOOP ECOSYSTEM
% =============================================================================
\section{Distributed Processing Engines and the Extended Hadoop Ecosystem}
\label{sec:dbms_hadoop_ecosystem}

\subsection{YARN: Yet Another Resource Negotiator}
YARN separates cluster resource management from computational programming abstractions, enabling multiple processing frameworks (MapReduce, Spark, Flink) to share the same physical cluster:
\begin{itemize}
    \item \textbf{ResourceManager (RM):} Cluster-wide authority that allocates memory and CPU resources across all running applications.
    \item \textbf{NodeManager (NM):} Per-node agent that launches and monitors resource-constrained \textbf{Containers}.
    \item \textbf{ApplicationMaster (AM):} Application-specific coordinator negotiated per job to manage task execution and handle worker failures.
\end{itemize}

\subsection{The MapReduce Computational Model}
MapReduce divides computations into three distinct, fault-tolerant phases:
\begin{enumerate}
    \item \textbf{Map Phase:} Reads input splits and transforms raw records into intermediate key-value pairs:
    \[
    \text{Map}: (k_1, v_1) \to \text{list}(k_2, v_2)
    \]
    \item \textbf{Shuffle and Sort Phase:} Collects, partitions across the network, and sorts intermediate records by key:
    \[
    \text{Shuffle}: \text{list}(k_2, v_2) \to (k_2, \text{list}(v_2))
    \]
    \item \textbf{Reduce Phase:} Aggregates all values associated with each unique intermediate key:
    \[
    \text{Reduce}: (k_2, \text{list}(v_2)) \to \text{list}(k_3, v_3)
    \]
\end{enumerate}

\subsection{Apache Spark versus Apache Hadoop MapReduce}
\begin{table}[htbp]
\centering
\small
\caption{Architectural Comparison: Hadoop MapReduce versus Apache Spark}
\label{tab:mapreduce_vs_spark}
\begin{tabularx}{\textwidth}{p{3.2cm} X X}
\toprule
\textbf{Architectural Dimension} & \textbf{Hadoop MapReduce} & \textbf{Apache Spark} \\
\midrule
\textbf{Processing Paradigm} & Rigid two-stage Map $\to$ Reduce model. & General Directed Acyclic Graph (DAG) computation engine. \\
\textbf{Intermediate Storage} & Persists intermediate outputs strictly to physical local disk. & Retains working datasets in RAM using \textbf{Resilient Distributed Datasets (RDDs)}. \\
\textbf{Iterative Workloads} & Extremely slow due to repeated disk I/O serialization bottlenecks. & Up to \textbf{100$\times$ faster} for iterative algorithms (machine learning, PageRank). \\
\textbf{Supported Workloads} & Batch processing only. & Unified platform: Batch (Spark Core), SQL (Spark SQL), Streaming, MLlib, GraphX. \\
\bottomrule
\end{tabularx}
\end{table}

\subsection{The Extended Apache Hadoop Ecosystem Suite}
\begin{figure}[htbp]
\centering
\small
\begin{tabularx}{\textwidth}{l p{3.2cm} X}
\toprule
\textbf{Technology} & \textbf{Category} & \textbf{Primary Engineering Responsibility} \\
\midrule
\textbf{Apache Hive} & Data Warehousing & Provides SQL-like abstraction (\textbf{HiveQL}) for batch analytics over HDFS datasets. \\
\textbf{Apache Pig} & Data Flow Scripting & Procedural ETL data-flow scripting platform using \textbf{Pig Latin}. \\
\textbf{Apache HBase} & NoSQL Database & Distributed, column-family NoSQL database providing low-latency $O(1)$ random read/write access on HDFS. \\
\textbf{Apache Sqoop} & Bulk Data Transfer & Bidirectional bulk data transfer tool connecting RDBMSs (MySQL, Oracle) with HDFS. \\
\textbf{Apache Flume} & Streaming Log Ingestion & Distributed service for collecting and ingesting high-volume streaming server logs into HDFS. \\
\textbf{Apache Oozie} & Workflow Scheduling & Server-based workflow scheduling engine for executing multi-stage Hadoop job DAGs. \\
\textbf{Apache ZooKeeper} & Distributed Coordination & High-performance coordination service for configuration management, leader election, and distributed locking. \\
\textbf{Apache Kafka} & Event Streaming Platform & High-throughput distributed publish-subscribe messaging log for real-time streaming architectures. \\
\bottomrule
\end{tabularx}
\end{figure}

\begin{quote}
\textbf{Definitive Technical Mnemonic: HDFS vs HBase vs Hive vs Hadoop}
\begin{itemize}
    \item \textbf{Hadoop:} The umbrella distributed ecosystem and framework.
    \item \textbf{HDFS:} The distributed storage layer (large sequential files).
    \item \textbf{HBase:} The NoSQL database layer (low-latency random point lookups).
    \item \textbf{Hive:} The data warehouse query layer (SQL analytical batch processing).
\end{itemize}
\end{quote}



\newpage
\section{Comprehensive Worked Problem Bank (Part 1): Relational Foundations, Normalization, and Transactions}

This section contains thirty-three meticulously solved examination problems spanning the entire syllabus of database management, relational algebra, constraints, transaction processing, concurrency anomalies, and normalization.

\subsection{Section A: Basic DBMS Concepts}

\begin{enumerate}[label=\textbf{Problem \arabic*.},leftmargin=1.5cm]
\item \textbf{Determine Degree and Cardinality:}
Consider the relation \texttt{Employees}:
\begin{table}[htbp]
\centering
\begin{tabularx}{0.8\textwidth}{c l c c}
\toprule
\textbf{Employee\_ID} & \textbf{Employee\_Name} & \textbf{Department} & \textbf{Salary} \\
\midrule
101 & Rahim & IT & 50000 \\
102 & Karim & HR & 45000 \\
103 & Nabila & Finance & 60000 \\
104 & Tania & IT & 55000 \\
105 & Hasan & HR & 48000 \\
\bottomrule
\end{tabularx}
\end{table}
Determine the Degree and Cardinality of this relation.
\begin{itemize}
\item \textbf{Solution:}
\begin{itemize}
\item \textbf{Degree:} Total number of attributes (columns) $= 4$ (\texttt{Employee\_ID, Employee\_Name, Department, Salary}).
\item \textbf{Cardinality:} Total number of tuples (rows) $= 5$ (Five employee records).
\end{itemize}
\[ \boxed{\text{\textbf{Degree}} = \mathbf{4}, \quad \text{\textbf{Cardinality}} = \mathbf{5}} \]
\end{itemize}

\item \textbf{Identify Primary Key and Foreign Key:}
Consider two relations: \texttt{Department(Department\_ID, Department\_Name)} and \texttt{Student(Student\_ID, Name, Department\_ID)}. Identify the primary and foreign keys.
\begin{itemize}
\item \textbf{Solution:}
\begin{itemize}
\item In \texttt{Department}: \texttt{Department\_ID} is the \textbf{Primary Key}.
\item In \texttt{Student}: \texttt{Student\_ID} is the \textbf{Primary Key}.
\item In \texttt{Student}: \texttt{Department\_ID} references \texttt{Department(Department\_ID)}, making it the \textbf{Foreign Key}.
\end{itemize}
\[ \boxed{\text{PK}_{\text{Dept}} = \texttt{Department\_ID}, \quad \text{PK}_{\text{Student}} = \texttt{Student\_ID}, \quad \text{FK}_{\text{Student}} = \texttt{Department\_ID}} \]
\end{itemize}

\item \textbf{Determine Candidate Keys and Super Keys:}
A relation \texttt{Employee(Employee\_ID, Email, Name, Department)} has unique \texttt{Employee\_ID} and unique \texttt{Email}. Identify candidate keys and give examples of super keys.
\begin{itemize}
\item \textbf{Solution:}
Both attributes independently guarantee unique row identification and are minimal:
\[ \boxed{\text{\textbf{Candidate Keys}} = \{\{\texttt{Employee\_ID}\}, \{\texttt{Email}\}\}} \]
\textbf{Super Keys} include any superset containing at least one candidate key: $\{\texttt{Employee\_ID}\}$, $\{\texttt{Email}\}$, $\{\texttt{Employee\_ID, Name}\}$, $\{\texttt{Email, Department}\}$, $\{\texttt{Employee\_ID, Email}\}$.
\end{itemize}

\item \textbf{Identify Relationship Cardinality Ratio:}
A university allows a student to enroll in multiple courses, and each course contains multiple enrolled students. Identify the relationship type and explain its relational implementation.
\begin{itemize}
\item \textbf{Solution:}
\[ \boxed{\text{\textbf{Many-to-Many (M:N) Relationship}}} \]
Implemented via an intermediate junction table \texttt{Enrollment(Student\_ID, Course\_ID, Grade)} with composite primary key $(\texttt{Student\_ID, Course\_ID})$ and foreign keys referencing \texttt{Students} and \texttt{Courses}.
\end{itemize}
\end{enumerate}

\subsection{Section B: Database Constraints and Integrity}

\begin{enumerate}[label=\textbf{Problem \arabic*.,resume},leftmargin=1.5cm]
\item \textbf{Write SQL DDL with Complete Constraints:}
Create a table named \texttt{Employees} where: (1) \texttt{Employee\_ID} is the primary key; (2) \texttt{Employee\_Name} is mandatory; (3) \texttt{Email} is unique; (4) \texttt{Salary} must be strictly positive; (5) \texttt{Status} defaults to \texttt{'Active'}.
\begin{itemize}
\item \textbf{Solution:}
\begin{lstlisting}[language=SQL]
CREATE TABLE Employees (
    Employee_ID INT PRIMARY KEY,
    Employee_Name VARCHAR(100) NOT NULL,
    Email VARCHAR(150) UNIQUE,
    Salary DECIMAL(12,2) CHECK (Salary > 0),
    Status VARCHAR(20) DEFAULT 'Active'
);
\end{lstlisting}
\end{itemize}

\item \textbf{Identify Referential Integrity Violation:}
A student record is inserted with $\texttt{Department\_ID} = 99$. However, \texttt{Departments} contains only departments 10 and 20. Which integrity rule is violated?
\begin{itemize}
\item \textbf{Solution:}
\[ \boxed{\text{\textbf{Referential Integrity Violation}}} \]
The child foreign key references a nonexistent parent primary key.
\end{itemize}

\item \textbf{Identify Entity Integrity Violation:}
A \texttt{Students} table with primary key \texttt{Student\_ID} contains two rows with $\texttt{Student\_ID} = 101$. What integrity rule is violated?
\begin{itemize}
\item \textbf{Solution:}
\[ \boxed{\text{\textbf{Entity Integrity Violation}}} \]
Entity integrity mandates that primary keys must be strictly unique and non-null.
\end{itemize}

\item \textbf{Identify Domain Integrity Violation:}
An account balance column defined to accept non-negative numbers records a value of $-5,000$. Which integrity rule is violated?
\begin{itemize}
\item \textbf{Solution:}
\[ \boxed{\text{\textbf{Domain Integrity Violation}}} \]
Enforced via \texttt{CHECK (Balance >= 0)}.
\end{itemize}
\end{enumerate}

\subsection{Section C: ACID Properties and Transaction Processing}

\begin{enumerate}[label=\textbf{Problem \arabic*.,resume},leftmargin=1.5cm]
\item \textbf{Bank Transfer Balance Calculation:}
Account A holds BDT 20,000 and Account B holds BDT 12,000. Account A transfers BDT 5,000 to Account B. Calculate the final individual and total balances.
\begin{itemize}
\item \textbf{Solution:}
\[ A_{\text{final}} = 20,000 - 5,000 = \text{BDT } 15,000 \]
\[ B_{\text{final}} = 12,000 + 5,000 = \text{BDT } 17,000 \]
\[ \text{Total Balance} = 15,000 + 17,000 = \text{BDT } 32,000 \quad (\text{Conserved}) \]
\[ \boxed{A = \text{\textbf{BDT 15,000}}, \quad B = \text{\textbf{BDT 17,000}}, \quad \text{\textbf{Total}} = \text{\textbf{BDT 32,000}}} \]
\end{itemize}

\item \textbf{Identify the ACID Property (Atomicity):}
A bank server deducts BDT 3,000 from Account A, but crashes before depositing the money into Account B. Which ACID property prevents the incomplete transaction from remaining committed?
\begin{itemize}
\item \textbf{Solution:}
\[ \boxed{\text{\textbf{Atomicity (All-or-Nothing Guarantee)}}} \]
\end{itemize}

\item \textbf{Identify the ACID Property (Consistency):}
A bank rule mandates that balances cannot drop below zero. A customer with BDT 5,000 attempts to withdraw BDT 10,000, and the database rejects the transaction. Which ACID property is demonstrated?
\begin{itemize}
\item \textbf{Solution:}
\[ \boxed{\text{\textbf{Consistency}}} \]
\end{itemize}

\item \textbf{Identify the ACID Property (Isolation):}
Two users attempt to purchase the final available airline seat simultaneously. Concurrency controls ensure that only one transaction succeeds while the other is rejected. Which ACID property is demonstrated?
\begin{itemize}
\item \textbf{Solution:}
\[ \boxed{\text{\textbf{Isolation}}} \]
\end{itemize}

\item \textbf{Identify the ACID Property (Durability):}
A database server confirms a payment transaction commit. A catastrophic power outage occurs 2 seconds later. Upon reboot, the payment record is present and intact. Which ACID property is demonstrated?
\begin{itemize}
\item \textbf{Solution:}
\[ \boxed{\text{\textbf{Durability}}} \]
\end{itemize}

\item \textbf{Write a Robust Transaction in SQL:}
Write an explicit SQL transaction to transfer BDT 1,000 from Account 101 to Account 102, ensuring that Account 101 has sufficient balance.
\begin{itemize}
\item \textbf{Solution:}
\begin{lstlisting}[language=SQL]
BEGIN;

UPDATE Accounts
SET Balance = Balance - 1000
WHERE Account_ID = 101 AND Balance >= 1000;

-- Application checks affected row count == 1
UPDATE Accounts
SET Balance = Balance + 1000
WHERE Account_ID = 102;

COMMIT;
\end{lstlisting}
\end{itemize}

\item \textbf{Identify the Concurrency Anomaly (Non-Repeatable Read):}
Transaction T1 reads balance $= 10,000$. Transaction T2 updates the balance to $12,000$ and commits. T1 reads the same row again and observes $12,000$. Identify the anomaly.
\begin{itemize}
\item \textbf{Solution:}
\[ \boxed{\text{\textbf{Non-Repeatable Read (Fuzzy Read)}}} \]
\end{itemize}

\item \textbf{Identify a Dirty Read:}
Transaction T1 updates an employee salary from $50,000$ to $60,000$ without committing. Transaction T2 reads the salary of $60,000$. T1 subsequently aborts and rolls back. What occurred?
\begin{itemize}
\item \textbf{Solution:}
\[ \boxed{\text{\textbf{Dirty Read Anomaly (Reading Uncommitted Data)}}} \]
\end{itemize}

\item \textbf{Identify a Lost Update and Calculate True Balance:}
Account balance is BDT 5,000. T1 reads 5,000 and adds 500. Concurrently, T2 reads 5,000 and subtracts 300. T1 writes 5,500. T2 subsequently writes 4,700 based on its stale read. Calculate the correct balance and identify the anomaly.
\begin{itemize}
\item \textbf{Solution:}
True balance $= 5,000 + 500 - 300 = \mathbf{5,200}$.
T1's deposit was completely overwritten.
\[ \boxed{\text{\textbf{Correct Balance}} = \text{\textbf{BDT 5,200}}, \quad \text{\textbf{Anomaly}} = \text{\textbf{Lost Update}}} \]
\end{itemize}

\item \textbf{Select the Appropriate Isolation Level:}
A financial ledger requires the strongest standard isolation guarantee to eliminate all dirty reads, non-repeatable reads, and phantom reads. Which isolation level must be configured?
\begin{itemize}
\item \textbf{Solution:}
\[ \boxed{\text{\textbf{SERIALIZABLE}}} \]
\end{itemize}
\end{enumerate}

\subsection{Section D: Normalization and Functional Dependencies}

\begin{enumerate}[label=\textbf{Problem \arabic*.,resume},leftmargin=1.5cm]
\item \textbf{Identify 1NF Violation and Remediate:}
A table has row: $\langle 101, \text{'Rahim'}, \text{'01711111111, 01822222222'} \rangle$. Does it satisfy 1NF? How is it resolved?
\begin{itemize}
\item \textbf{Solution:}
\textbf{No, it violates 1NF} because \texttt{Phone\_Numbers} contains non-atomic multi-valued data.
\textbf{Remediation:} Split into two rows: $\langle 101, \text{'Rahim'}, \text{'01711111111'} \rangle$ and $\langle 101, \text{'Rahim'}, \text{'01822222222'} \rangle$.
\end{itemize}

\item \textbf{Identify Partial Functional Dependency:}
In \texttt{Enrollment(Student\_ID, Course\_ID, Student\_Name, Course\_Name, Grade)} with primary key $(\texttt{Student\_ID, Course\_ID})$:
\[ \texttt{Student\_ID} \rightarrow \texttt{Student\_Name}, \quad \texttt{Course\_ID} \rightarrow \texttt{Course\_Name}, \quad \{\texttt{Student\_ID, Course\_ID}\} \rightarrow \texttt{Grade} \]
Identify all partial dependencies.
\begin{itemize}
\item \textbf{Solution:}
Both \texttt{Student\_Name} and \texttt{Course\_Name} depend on proper subsets of the composite primary key.
\[ \boxed{\texttt{Student\_ID} \rightarrow \texttt{Student\_Name} \quad \text{and} \quad \texttt{Course\_ID} \rightarrow \texttt{Course\_Name} \quad \text{\textbf{(Partial Dependencies)}}} \]
This table violates \textbf{2NF}.
\end{itemize}

\item \textbf{Identify Transitive Dependency:}
In \texttt{Students(Student\_ID, Student\_Name, Department\_ID, Department\_Name)}:
\[ \texttt{Student\_ID} \rightarrow \texttt{Department\_ID} \quad \text{and} \quad \texttt{Department\_ID} \rightarrow \texttt{Department\_Name} \]
Identify the transitive dependency.
\begin{itemize}
\item \textbf{Solution:}
\texttt{Department\_Name} depends on non-key attribute \texttt{Department\_ID}, which depends on candidate key \texttt{Student\_ID}.
\[ \boxed{\texttt{Student\_ID} \rightarrow \texttt{Department\_Name} \quad \text{\textbf{(Transitive Dependency)}}} \]
This table violates \textbf{3NF}.
\end{itemize}

\item \textbf{Decompose Relation into 2NF:}
Given \texttt{Order\_Details(Order\_ID, Product\_ID, Order\_Date, Product\_Name, Quantity)} with dependencies:
\[ \{\texttt{Order\_ID, Product\_ID}\} \rightarrow \texttt{Quantity}, \quad \texttt{Order\_ID} \rightarrow \texttt{Order\_Date}, \quad \texttt{Product\_ID} \rightarrow \texttt{Product\_Name} \]
Convert into 2NF relations.
\begin{itemize}
\item \textbf{Solution:}
Decompose to isolate partial dependencies:
\begin{enumerate}
\item \texttt{Orders(}\underline{\texttt{Order\_ID}}, \texttt{Order\_Date)}
\item \texttt{Products(}\underline{\texttt{Product\_ID}}, \texttt{Product\_Name)}
\item \texttt{Order\_Items(}\underline{\texttt{Order\_ID}}, \underline{\texttt{Product\_ID}}, \texttt{Quantity)}
\end{enumerate}
\end{itemize}

\item \textbf{Decompose Relation into 3NF:}
Given \texttt{Employees(Employee\_ID, Employee\_Name, Department\_ID, Department\_Name)} where \texttt{Employee\_ID} is primary key and $\texttt{Department\_ID} \rightarrow \texttt{Department\_Name}$. Convert into 3NF.
\begin{itemize}
\item \textbf{Solution:}
Decompose to remove transitive dependency:
\begin{enumerate}
\item \textbf{\texttt{Employees:}} $(\underline{\texttt{Employee\_ID}}, \texttt{Employee\_Name}, \texttt{Department\_ID})$
\item \textbf{\texttt{Departments:}} $(\underline{\texttt{Department\_ID}}, \texttt{Department\_Name})$
\end{enumerate}
\end{itemize}

\item \textbf{Identify Normal Form Flaw:}
A relation satisfies 2NF but contains a transitive dependency where non-prime attribute $C$ depends on non-prime attribute $B$. Which normal form is violated?
\begin{itemize}
\item \textbf{Solution:}
\[ \boxed{\text{\textbf{Third Normal Form (3NF) is violated}}} \]
\end{itemize}

\item \textbf{Identify BCNF Violation:}
A relation contains non-trivial dependency $\texttt{Instructor} \rightarrow \texttt{Course}$, but \texttt{Instructor} is not a super key. Does the relation satisfy BCNF?
\begin{itemize}
\item \textbf{Solution:}
\textbf{No, BCNF is violated.} BCNF mandates that every non-trivial determinant must be a super key.
\[ \boxed{\text{\textbf{Violates BCNF}}} \]
\end{itemize}

\item \textbf{Identify 4NF Violation:}
A student has multiple independent skills and multiple independent hobbies. A table stores all cross-product combinations $(\texttt{Student\_ID, Skill, Hobby})$. Which normal form addresses this?
\begin{itemize}
\item \textbf{Solution:}
\[ \boxed{\text{\textbf{Fourth Normal Form (4NF) - Addresses Multivalued Dependencies}}} \]
\end{itemize}

\item \textbf{Normal Form Matching Directory:}
Match each relational defect to the normal form that eliminates it:
\begin{itemize}
\item Multi-valued non-atomic cells $\to$ \textbf{1NF}.
\item Partial dependency on composite key $\to$ \textbf{2NF}.
\item Transitive dependency of non-prime attributes $\to$ \textbf{3NF}.
\item Non-trivial determinant is not a super key $\to$ \textbf{BCNF}.
\item Independent multivalued dependencies ($X \twoheadrightarrow Y$) $\to$ \textbf{4NF}.
\item Join dependencies requiring 3-way lossless decomposition $\to$ \textbf{5NF (PJNF)}.
\end{itemize}
\end{enumerate}

\subsection{Section E: Practical SQL Queries and Relational Operations}

\begin{enumerate}[label=\textbf{Problem \arabic*.,resume},leftmargin=1.5cm]
\item \textbf{Retrieve All Records:} Write a query to retrieve all attributes and rows from \texttt{Students}.
\begin{lstlisting}[language=SQL]
SELECT * FROM Students;
\end{lstlisting}

\item \textbf{Filter Records by Exact Match:} Retrieve all students enrolled in the 'CSE' department.
\begin{lstlisting}[language=SQL]
SELECT * FROM Students WHERE Department = 'CSE';
\end{lstlisting}

\item \textbf{Filter by Numerical Range:} Retrieve student ID, name, and CGPA for students with $\text{CGPA} > 3.50$.
\begin{lstlisting}[language=SQL]
SELECT Student_ID, Name, CGPA FROM Students WHERE CGPA > 3.50;
\end{lstlisting}

\item \textbf{Sort Query Results:} Retrieve all students sorted in descending order of CGPA (highest to lowest).
\begin{lstlisting}[language=SQL]
SELECT Student_ID, Name, CGPA FROM Students ORDER BY CGPA DESC;
\end{lstlisting}

\item \textbf{Aggregate Group By Query:} Count the total number of students enrolled in each department.
\begin{lstlisting}[language=SQL]
SELECT Department, COUNT(*) AS Total_Students
FROM Students
GROUP BY Department;
\end{lstlisting}

\item \textbf{Inner Join Query:} Retrieve student ID, student name, and department name by joining \texttt{Students} and \texttt{Departments} on \texttt{Department\_ID}.
\begin{lstlisting}[language=SQL]
SELECT s.Student_ID, s.Name, d.Department_Name
FROM Students s
INNER JOIN Departments d ON s.Department_ID = d.Department_ID;
\end{lstlisting}
\end{enumerate}

\newpage



% =============================================================================
% SECTION 38: WORKED PROBLEM BANK (PART 2): SQL, JDBC, MONGODB & BIG DATA
% =============================================================================
\section{Comprehensive Worked Problem Bank (Part 2): SQL Queries, JDBC, MongoDB, and Big Data}
\label{sec:dbms_problems_part2}

This section delivers twenty-eight worked examination problems covering SQL DQL queries, aggregate calculations, subqueries, JDBC programming, MongoDB operations, HDFS block mathematics, and MapReduce trace simulations.

\subsection{Section A: SQL Data Query Problems}

Consider the relational schema below representing an organizational enterprise:
\begin{center}
\textbf{Employees}(\underline{Employee\_ID}, Name, Department\_ID, Salary) \\
\textbf{Departments}(\underline{Department\_ID}, Department\_Name)
\end{center}

\begin{table}[htbp]
\centering
\small
\caption{Sample Operational Dataset}
\begin{minipage}{0.55\textwidth}
\centering
\textbf{Employees}
\begin{tabularx}{\textwidth}{c l c r}
\toprule
\textbf{ID} & \textbf{Name} & \textbf{Dept\_ID} & \textbf{Salary} \\
\midrule
101 & Rahim  & 10 & 50,000 \\
102 & Karim  & 20 & 45,000 \\
103 & Nabila & 10 & 60,000 \\
104 & Tania  & 30 & 55,000 \\
105 & Hasan  & 20 & 48,000 \\
106 & Sara   & 10 & 70,000 \\
\bottomrule
\end{tabularx}
\end{minipage}
\hfill
\begin{minipage}{0.40\textwidth}
\centering
\textbf{Departments}
\begin{tabularx}{\textwidth}{c l}
\toprule
\textbf{Dept\_ID} & \textbf{Dept\_Name} \\
\midrule
10 & IT \\
20 & HR \\
30 & Finance \\
40 & Marketing \\
\bottomrule
\end{tabularx}
\end{minipage}
\end{table}

\begin{enumerate}[label=\textbf{Problem \arabic*.},leftmargin=1.5cm]
\item \textbf{Retrieve All Employees in the IT Department:}
\begin{lstlisting}[language=SQL]
SELECT * FROM Employees WHERE Department_ID = 10;
\end{lstlisting}

\item \textbf{Retrieve Employees with Salary Exceeding BDT 50,000:}
\begin{lstlisting}[language=SQL]
SELECT Name, Salary FROM Employees WHERE Salary > 50000;
\end{lstlisting}

\item \textbf{Find the Maximum Company Salary:}
\begin{lstlisting}[language=SQL]
SELECT MAX(Salary) AS Highest_Salary FROM Employees;
\end{lstlisting}
\textbf{Result:} $\boxed{70,000\text{ BDT}}$ (Sara).

\item \textbf{Calculate Average Company Salary:}
\begin{lstlisting}[language=SQL]
SELECT AVG(Salary) AS Average_Salary FROM Employees;
\end{lstlisting}
\textbf{Mathematical Calculation:}
\[
\text{Total Salary} = 50,000 + 45,000 + 60,000 + 55,000 + 48,000 + 70,000 = 328,000
\]
\[
\text{Average Salary} = \frac{328,000}{6} \approx \mathbf{54,666.67\text{ BDT}}
\]

\item \textbf{Count Employees in Each Department:}
\begin{lstlisting}[language=SQL]
SELECT Department_ID, COUNT(*) AS Total_Employees
FROM Employees GROUP BY Department_ID;
\end{lstlisting}
\textbf{Result:} Dept 10: 3; Dept 20: 2; Dept 30: 1.

\item \textbf{Find Departments with Headcount Greater Than One:}
\begin{lstlisting}[language=SQL]
SELECT Department_ID, COUNT(*) AS Total_Employees
FROM Employees GROUP BY Department_ID HAVING COUNT(*) > 1;
\end{lstlisting}
\textbf{Result:} Dept 10 (3 employees) and Dept 20 (2 employees).

\item \textbf{Display Employee Names Alongside Department Names (INNER JOIN):}
\begin{lstlisting}[language=SQL]
SELECT e.Name, d.Department_Name
FROM Employees e
INNER JOIN Departments d ON e.Department_ID = d.Department_ID;
\end{lstlisting}

\item \textbf{Find Departments Without Any Assigned Employees (LEFT JOIN):}
\begin{lstlisting}[language=SQL]
SELECT d.Department_ID, d.Department_Name
FROM Departments d
LEFT JOIN Employees e ON d.Department_ID = e.Department_ID
WHERE e.Employee_ID IS NULL;
\end{lstlisting}
\textbf{Result:} Department 40 (Marketing) has no associated employees.

\item \textbf{Find Employees Earning Above Company Average (Scalar Subquery):}
\begin{lstlisting}[language=SQL]
SELECT Name, Salary FROM Employees
WHERE Salary > (SELECT AVG(Salary) FROM Employees);
\end{lstlisting}
\textbf{Result:} Nabila ($60,000$), Tania ($55,000$), and Sara ($70,000$) earn above $54,666.67$.

\item \textbf{Find the Second-Highest Distinct Salary:}
\begin{lstlisting}[language=SQL]
SELECT MAX(Salary) AS Second_Highest_Salary
FROM Employees
WHERE Salary < (SELECT MAX(Salary) FROM Employees);
\end{lstlisting}
\textbf{Analysis:} The maximum salary is $70,000$. The maximum salary strictly strictly less than $70,000$ is $\boxed{60,000\text{ BDT}}$ (Nabila). This query gracefully handles ties because \texttt{MAX} evaluates the distinct scalar ceiling.
\end{enumerate}

\subsection{Section B: SQL Conceptual and Analytical Problems}

\begin{enumerate}[resume*]
\item \textbf{WHERE versus HAVING Selection:}
\textit{Question:} Which clause must be used to filter departments whose average salary exceeds BDT 50,000? Explain why. \\
\textbf{Solution:} \textbf{\texttt{HAVING}} must be used:
\begin{lstlisting}[language=SQL]
SELECT Department_ID, AVG(Salary) AS Avg_Salary
FROM Employees GROUP BY Department_ID HAVING AVG(Salary) > 50000;
\end{lstlisting}
\textbf{Reason:} The \texttt{WHERE} clause filters individual candidate tuples \textbf{prior} to grouping; it cannot evaluate aggregate functions. The \texttt{HAVING} clause executes \textbf{after} \texttt{GROUP BY} aggregation and can filter on aggregated results.

\item \textbf{Behavioral Analysis of COUNT(*) versus COUNT(column):}
\textit{Question:} Suppose the \texttt{Students} table has 5 rows, but the \texttt{Email} column contains 1 \texttt{NULL} value. What are the results of \texttt{SELECT COUNT(*)} and \texttt{SELECT COUNT(Email)}? \\
\textbf{Solution:}
\[
\text{COUNT}(*) = 5, \quad \text{COUNT}(\text{Email}) = 4
\]
\texttt{COUNT(*)} counts total rows regardless of nulls; \texttt{COUNT(column)} counts only non-NULL entries.

\item \textbf{UNION versus UNION ALL Evaluation:}
\textit{Question:} Relation $A$ contains $\{\text{CSE}, \text{EEE}, \text{BBA}\}$. Relation $B$ contains $\{\text{CSE}, \text{EEE}, \text{LAW}\}$. What are the output cardinalities of \texttt{UNION} and \texttt{UNION ALL}? \\
\textbf{Solution:}
\begin{itemize}
    \item \texttt{SELECT Dept FROM A UNION SELECT Dept FROM B;} produces \textbf{4 rows} ($\{\text{CSE}, \text{EEE}, \text{BBA}, \text{LAW}\}$) because duplicate tuples are eliminated.
    \item \texttt{SELECT Dept FROM A UNION ALL SELECT Dept FROM B;} produces \textbf{6 rows} ($\{\text{CSE}, \text{EEE}, \text{BBA}, \text{CSE}, \text{EEE}, \text{LAW}\}$) because duplicates are preserved.
\end{itemize}

\item \textbf{Selecting the Correct Join Type for Total Customer Retention:}
\textit{Question:} A banking portal query must display all registered customers, including those who have not yet opened an active bank account. Which join type is required? \\
\textbf{Solution:} A \textbf{\texttt{LEFT OUTER JOIN}}, placing \texttt{Customers} as the left table:
\begin{lstlisting}[language=SQL]
SELECT c.Customer_ID, c.Full_Name, a.Account_Number
FROM Customers c
LEFT JOIN Account_Holders ah ON c.Customer_ID = ah.Customer_ID
LEFT JOIN Accounts a         ON ah.Account_ID = a.Account_ID;
\end{lstlisting}
Customers without accounts will retain their personal rows with \texttt{NULL} account columns.
\end{enumerate}

\subsection{Section C: Java Database Connectivity (JDBC) Problems}

\begin{enumerate}[resume*]
\item \textbf{JDBC Component Functional Mapping:}
\begin{table}[htbp]
\centering
\small
\begin{tabularx}{0.85\textwidth}{l X}
\toprule
\textbf{Application Database Task} & \textbf{Standard JDBC API Component} \\
\midrule
Obtain active database connection & \texttt{DriverManager.getConnection()} or \texttt{DataSource} \\
Execute parameterized dynamic SQL & \texttt{PreparedStatement} \\
Execute database stored procedures & \texttt{CallableStatement} \\
Navigate returned query records & \texttt{ResultSet} (via \texttt{rs.next()}) \\
Persist atomic transaction changes & \texttt{Connection.commit()} \\
Abort transaction and rollback & \texttt{Connection.rollback()} \\
\bottomrule
\end{tabularx}
\end{table}

\item \textbf{Identify JDBC Driver Architecture:}
\textit{Question:} An enterprise Java microservice connects directly to an Oracle database instance over TCP/IP without requiring any native C/C++ client libraries or middleware proxy servers. Which JDBC driver type is deployed? \\
\textbf{Solution:} $\boxed{\textbf{Type 4 JDBC Driver (Thin / Native-Protocol Driver)}}$.

\item \textbf{Remediating SQL Injection in Java:}
\textit{Question:} A legacy login servlet executes:
\begin{lstlisting}[language=Java]
String sql = "SELECT * FROM Users WHERE User = '" + user + "' AND Pass = '" + pass + "'";
\end{lstlisting}
Provide the secure refactoring. \\
\textbf{Solution:} Refactor to use \texttt{PreparedStatement} with parameterized placeholders and salted cryptographic password verification (e.g., Argon2 or BCrypt):
\begin{lstlisting}[language=Java]
String sql = "SELECT Password_Hash, Salt FROM Users WHERE Username = ?";
try (PreparedStatement ps = conn.prepareStatement(sql)) {
    ps.setString(1, user);
    try (ResultSet rs = ps.executeQuery()) {
        if (rs.next()) {
            String hash = rs.getString("Password_Hash");
            // Verify password using cryptographically secure hashing function
        }
    }
}
\end{lstlisting}
\end{enumerate}

\subsection{Section D: MongoDB Document Database Problems}

\begin{enumerate}[resume*]
\item \textbf{Document Insertion:}
Insert a document into collection \texttt{students} for student Rahim (ID 101, CSE, CGPA 3.75).
\begin{lstlisting}[language=Java]
db.students.insertOne({ _id: 101, name: "Rahim", department: "CSE", cgpa: 3.75 });
\end{lstlisting}

\item \textbf{Filter by Equality:}
Retrieve all students belonging to the CSE department.
\begin{lstlisting}[language=Java]
db.students.find({ department: "CSE" });
\end{lstlisting}

\item \textbf{Range Filter Operator:}
Retrieve all students having a CGPA strictly greater than 3.50.
\begin{lstlisting}[language=Java]
db.students.find({ cgpa: { $gt: 3.50 } });
\end{lstlisting}

\item \textbf{Atomic In-Place Update:}
Update Rahim's CGPA to 3.90 in collection \texttt{students}.
\begin{lstlisting}[language=Java]
db.students.updateOne({ _id: 101 }, { $set: { cgpa: 3.90 } });
\end{lstlisting}

\item \textbf{Aggregation Pipeline Grouping:}
Calculate the total number of students enrolled per department.
\begin{lstlisting}[language=Java]
db.students.aggregate([
  { $group: { _id: "$department", totalStudents: { $sum: 1 } } }
]);
\end{lstlisting}

\item \textbf{Identify MongoDB Scaling Mechanism:}
\textit{Question:} An e-commerce product collection expands beyond the storage limit of a single database server. Which architectural mechanism partitions the collection across multiple machines? \\
\textbf{Solution:} $\boxed{\textbf{Horizontal Sharding}}$ (using a Shard Key and \texttt{mongos} query routers).

\item \textbf{Identify High-Availability Clustering Feature:}
\textit{Question:} Which MongoDB feature maintains multiple mirrored copies of data across distinct servers and executes automated primary election upon failure? \\
\textbf{Solution:} $\boxed{\textbf{Replica Set / Replication}}$.
\end{enumerate}

\subsection{Section E: Apache Hadoop and Distributed Computing Problems}

\begin{enumerate}[resume*]
\item \textbf{Hadoop Component Responsibility Mapping:}
\begin{table}[htbp]
\centering
\small
\begin{tabularx}{0.85\textwidth}{l X}
\toprule
\textbf{Component} & \textbf{Core Architectural Role} \\
\midrule
\textbf{HDFS} & Scalable, distributed, fault-tolerant filesystem storage. \\
\textbf{NameNode} & Centralized manager of HDFS metadata and block mapping. \\
\textbf{DataNode} & Worker daemon storing physical data blocks on local disk. \\
\textbf{YARN} & Cluster-wide resource management and application scheduling. \\
\textbf{ResourceManager} & Master scheduler allocating memory/CPU across the cluster. \\
\textbf{NodeManager} & Worker agent managing local task containers and heartbeats. \\
\textbf{MapReduce} & Distributed programming framework (Map $\to$ Shuffle $\to$ Reduce). \\
\textbf{Apache Hive} & SQL-like analytical data warehousing interface (HiveQL). \\
\textbf{Apache HBase} & Distributed column-family NoSQL database for random read/write. \\
\bottomrule
\end{tabularx}
\end{table}

\item \textbf{HDFS Block Count Calculation:}
\textit{Question:} A dataset file of size \textbf{500 MB} is ingested into HDFS configured with a default block size of \textbf{128 MB}. Calculate the number of blocks created and specify the size of each block. \\
\textbf{Solution:}
\[
\text{Number of Blocks} = \left\lceil \frac{\text{File Size}}{\text{Block Size}} \right\rceil = \left\lceil \frac{500}{128} \right\rceil = \lceil 3.90625 \rceil = \mathbf{4\text{ blocks}}
\]
\textbf{Block Size Distribution:}
\begin{itemize}
    \item Block 1 = 128 MB
    \item Block 2 = 128 MB
    \item Block 3 = 128 MB
    \item Block 4 = $500 - (3 \times 128) = 500 - 384 = \mathbf{116\text{ MB}}$
\end{itemize}
\textit{Note: HDFS DataNodes allocate only the actual bytes consumed by a partial block (116 MB), avoiding disk wastage.}

\item \textbf{HDFS Storage Calculation with Replication:}
\textit{Question:} A data file of size \textbf{600 MB} is stored in an HDFS cluster configured with the standard replication factor of 3. What is the total physical disk storage consumed across the cluster? \\
\textbf{Solution:}
\[
\text{Physical Disk Storage} = \text{File Size} \times \text{Replication Factor} = 600\text{ MB} \times 3 = \mathbf{1,800\text{ MB}} \approx \mathbf{1.8\text{ GB}}
\]

\item \textbf{MapReduce Word Count Computational Trace:}
\textit{Question:} Trace the execution of a distributed MapReduce Word Count job given the input string:
\begin{center}
\texttt{"database sql database sql hadoop database hadoop sql"}
\end{center}
\textbf{Solution:}
\begin{itemize}
    \item \textbf{Map Phase Output (Key-Value Emission):}
    \[
    (\text{database}, 1), (\text{sql}, 1), (\text{database}, 1), (\text{sql}, 1), (\text{hadoop}, 1), (\text{database}, 1), (\text{hadoop}, 1), (\text{sql}, 1)
    \]
    \item \textbf{Shuffle and Sort Phase (Grouping by Key):}
    \begin{align*}
    \text{database} &\to [1, 1, 1] \\
    \text{hadoop}   &\to [1, 1] \\
    \text{sql}      &\to [1, 1, 1]
    \end{align*}
    \item \textbf{Reduce Phase Output (Key Aggregation):}
    \begin{center}
    \begin{tabular}{l c}
    \toprule
    \textbf{Word (Key)} & \textbf{Final Count (Value)} \\
    \midrule
    \texttt{database} & 3 \\
    \texttt{hadoop}   & 2 \\
    \texttt{sql}      & 3 \\
    \bottomrule
    \end{tabular}
    \end{center}
\end{itemize}
\end{enumerate}

% =============================================================================
% SECTION 39: MASTER COMPARATIVE ANALYSIS: 22 HIGH-PRIORITY DISTINCTIONS
% =============================================================================
\section{Master Technical Comparison Tables: 22 High-Priority Distinctions}
\label{sec:dbms_master_comparisons}

Technical recruitment interviews for IT positions frequently hinge on articulating crisp, semantically accurate distinctions between closely related concepts. The following 22 comparative matrices synthesize the core curriculum:

\begin{table}[htbp]
\centering
\scriptsize
\caption{The 22 High-Priority Recruitment Examination Distinctions}
\label{tab:22_distinctions}
\begin{tabularx}{\textwidth}{l p{3.2cm} X}
\toprule
\textbf{\#} & \textbf{Comparison Concept} & \textbf{Core Engineering Distinction} \\
\midrule
1 & \textbf{DBMS vs. RDBMS} & General data storage system vs. mathematically formalized relational engine enforcing Codd's rules, primary/foreign keys, and relational integrity. \\
2 & \textbf{Primary Key vs. Foreign Key} & Uniquely identifies a tuple within its own relation vs. references a primary/candidate key in another relation to enforce referential integrity. \\
3 & \textbf{Schema vs. Instance} & Static structural blueprint/metadata (intension) vs. dynamic state of stored tuples at an exact snapshot in time (extension). \\
4 & \textbf{Degree vs. Cardinality} & Total number of attributes (columns, $n$) vs. total number of records (tuples/rows, $m$) in a relation. \\
5 & \textbf{Candidate Key vs. Super Key} & Minimal super key with no redundant attributes vs. any attribute set that uniquely identifies a tuple. \\
6 & \textbf{CHAR vs. VARCHAR} & Fixed-length blank-padded character storage vs. variable-length storage with length prefix bytes. \\
7 & \textbf{DECIMAL vs. FLOAT} & Exact base-10 packed decimal arithmetic (mandatory for money) vs. approximate IEEE 754 binary floating point. \\
8 & \textbf{WHERE vs. HAVING} & Pre-aggregation row filter (no aggregates allowed) vs. post-aggregation group filter (aggregates standard). \\
9 & \textbf{GROUP BY vs. ORDER BY} & Collapses rows into summary equivalence groups vs. sorts final result tuples lexicographically/numerically. \\
10 & \textbf{INNER JOIN vs. LEFT JOIN} & Returns matching tuples from both tables vs. preserves all left tuples and pads unmatched right attributes with NULL. \\
11 & \textbf{JOIN vs. UNION} & Combines tables horizontally by column association vs. combines query results vertically by row concatenation. \\
12 & \textbf{UNION vs. UNION ALL} & Concatenates row sets with duplicate elimination (sort/hash pass) vs. pure concatenation preserving duplicates. \\
13 & \textbf{COUNT(*) vs. COUNT(col)} & Counts total candidate rows including NULLs vs. counts only rows where the specified attribute is non-NULL. \\
14 & \textbf{Correlated vs. Uncorrelated} & Nested query re-evaluated per outer candidate row vs. nested query executed once and cached globally. \\
15 & \textbf{RDBMS vs. MongoDB} & Rigid tabular schemas, ACID joins, SQL vs. flexible BSON document collections, sharded horizontal scaling, MQL. \\
16 & \textbf{JDBC vs. JDBC Driver} & Abstract Java SE API specification (\texttt{java.sql}) vs. vendor-supplied wire protocol implementation library. \\
17 & \textbf{Statement vs. PreparedStatement} & Static SQL compiled at runtime (SQL injection hazard) vs. precompiled parameterized query template (safe). \\
18 & \textbf{Replication vs. Sharding} & Duplicating identical data copies across nodes (high availability) vs. partitioning disjoint data across nodes (scalability). \\
19 & \textbf{HDFS vs. HBase} & Distributed sequential file storage system vs. distributed column-family NoSQL database for random $O(1)$ read/write. \\
20 & \textbf{Hive vs. HBase} & SQL-like batch analytics and data warehousing (HiveQL) vs. low-latency real-time key-value NoSQL database. \\
21 & \textbf{HDFS vs. YARN} & Physical distributed storage layer vs. cluster computational resource management and scheduling layer. \\
22 & \textbf{MapReduce vs. Apache Spark} & Rigid disk-oriented two-stage batch engine vs. in-memory DAG computation engine up to 100$\times$ faster for iterative workloads. \\
\bottomrule
\end{tabularx}
\end{table}

% =============================================================================
% SECTION 40: MASTER QUICK-REVISION CHEAT SHEET & GOLDEN FORMULAS
% =============================================================================
\section{Master Quick-Revision Cheat Sheet and Golden Formulas}
\label{sec:dbms_quick_revision}

\subsection{Consolidated High-Yield Revision Fast Facts}
\begin{itemize}
    \item \textbf{SQL Processing Sequence:} $\text{FROM} \to \text{WHERE} \to \text{GROUP BY} \to \text{HAVING} \to \text{SELECT} \to \text{DISTINCT} \to \text{ORDER BY} \to \text{LIMIT}$.
    \item \textbf{NULL Arithmetic:} Any standard arithmetic expression with \texttt{NULL} yields \texttt{NULL} ($x + \text{NULL} = \text{NULL}$). Any boolean comparison with \texttt{NULL} yields \texttt{UNKNOWN}.
    \item \textbf{Aggregate NULL Rule:} Except \texttt{COUNT(*)}, all aggregates ignore \texttt{NULL}. If all values are \texttt{NULL}, \texttt{SUM/AVG/MIN/MAX} return \texttt{NULL}, but \texttt{COUNT} returns \texttt{0}.
    \item \textbf{Banking Invariant:} Monetary amounts must use \texttt{DECIMAL(p, s)}, ledgers must enforce \texttt{ON DELETE RESTRICT}, and funds transfers require atomic multi-statement transactions.
    \item \textbf{JDBC Rule:} Modern production uses \textbf{Type 4 thin drivers} and \textbf{\texttt{PreparedStatement}} to defeat SQL injection.
    \item \textbf{MongoDB Rule:} BSON limit is 16 MiB. Single-document operations are atomic. Replication provides availability; Sharding provides horizontal capacity.
    \item \textbf{Hadoop Rule:} NameNode manages metadata; DataNodes store blocks. Default block size is 128 MB; default replication factor is 3. Secondary NameNode is a checkpointer, not a standby.
\end{itemize}

\subsection{The Eight Golden Mathematical Formulas for DBMS Examinations}
\begin{enumerate}
    \item \textbf{Upper Bound on Super Keys:} For a relation with $n$ attributes where candidate key $K$ consists of $k$ attributes:
    \begin{equation}
    \text{Total Super Keys containing } K = 2^{n - k}
    \end{equation}
    \item \textbf{Cartesian Product Cardinality and Degree:}
    \begin{equation}
    |\mathcal{R} \times \mathcal{S}| = |\mathcal{R}| \times |\mathcal{S}|, \quad \text{deg}(\mathcal{R} \times \mathcal{S}) = \text{deg}(\mathcal{R}) + \text{deg}(\mathcal{S})
    \end{equation}
    \item \textbf{Binary Decomposition Lossless Join Test:} A decomposition of $R$ into $(R_1, R_2)$ is lossless if and only if:
    \begin{equation}
    (R_1 \cap R_2 \to R_1) \in F^+ \quad \text{or} \quad (R_1 \cap R_2 \to R_2) \in F^+
    \end{equation}
    \item \textbf{Normal Form Hierarchy Formal Inclusion:}
    \begin{equation}
    5\text{NF} \subset 4\text{NF} \subset \text{BCNF} \subset 3\text{NF} \subset 2\text{NF} \subset 1\text{NF}
    \end{equation}
    \item \textbf{3NF Condition:} For every non-trivial FD $X \to A$, either $X$ is a super key, or $A$ is a prime attribute.
    \item \textbf{BCNF Condition:} For every non-trivial FD $X \to A$, $X$ must strictly be a super key.
    \item \textbf{HDFS Block Allocation Formula:}
    \begin{equation}
    N_{\text{blocks}} = \left\lceil \frac{\text{File Size}}{\text{Block Size}} \right\rceil
    \end{equation}
    \item \textbf{HDFS Physical Storage Requirement Formula:}
    \begin{equation}
    \text{Physical Disk Space} = \text{File Size} \times \text{Replication Factor}
    \end{equation}
\end{enumerate}

% =============================================================================
% SECTION 41: CATEGORIZED WRITTEN & VIVA EXAMINATION QUESTION BANK
% =============================================================================
\section{High-Yield Written Examination and Viva Voce Question Bank}
\label{sec:dbms_viva_questions}

This bank categorizes forty-eight comprehensive technical questions frequently posed in recruitment written tests and viva voce boards for government IT posts:

\subsection{Module A: DBMS and SQL Fundamentals}
\begin{enumerate}
    \item Define DBMS and RDBMS. Detail five fundamental structural differences between them.
    \item Differentiate between Database Schema (intension) and Database Instance (extension).
    \item Define Degree and Cardinality of a relation with a concrete mathematical example.
    \item Compare \texttt{CHAR(n)} and \texttt{VARCHAR(n)} in terms of physical disk storage, padding, and retrieval speed.
    \item Why is \texttt{DECIMAL(p, s)} strictly mandated over \texttt{FLOAT} or \texttt{DOUBLE} for financial and banking ledgers?
    \item Formally define Super Key, Candidate Key, Primary Key, and Alternate Key. State their set-theoretic subset relation.
    \item Explain the concept of Foreign Key. How does it enforce referential integrity between tables?
    \item Explain One-to-One, One-to-Many, and Many-to-Many entity relationships. How is an $M:N$ relationship physically implemented in relational SQL?
    \item Describe the ANSI/SPARC Three-Schema Architecture (External, Conceptual, Internal).
    \item Differentiate between Physical Data Independence and Logical Data Independence with practical schema change examples.
\end{enumerate}

\subsection{Module B: SQL Queries and Advanced DQL}
\begin{enumerate}[resume]
    \item Detail the standard syntax and functional purpose of all primary clauses of an SQL \texttt{SELECT} statement.
    \item Provide a comprehensive 5-point comparative analysis between the \texttt{WHERE} and \texttt{HAVING} clauses.
    \item Formally explain \texttt{INNER JOIN}, \texttt{LEFT OUTER JOIN}, \texttt{RIGHT OUTER JOIN}, \texttt{FULL OUTER JOIN}, and \texttt{CROSS JOIN}.
    \item Differentiate between \texttt{UNION} and \texttt{UNION ALL} in terms of tuple deduplication and execution overhead.
    \item Explain standard ANSI SQL aggregate functions (\texttt{COUNT}, \texttt{SUM}, \texttt{AVG}, \texttt{MIN}, \texttt{MAX}) and their behavior when encountering \texttt{NULL} values.
    \item Write an SQL query to retrieve the second-highest distinct salary from an \texttt{Employees} table without using vendor-specific limit offsets.
    \item Define Correlated Subquery and explain how its execution model differs from an uncorrelated subquery.
    \item Compare the set predicates \texttt{IN} and \texttt{EXISTS}. Why is \texttt{NOT EXISTS} preferred over \texttt{NOT IN} when handling columns with potential \texttt{NULL} values?
    \item Explain Database Indexing. Contrast Clustered Indexes with Non-Clustered Indexes and state why a table can have only one clustered index.
    \item Detail the referential actions: \texttt{ON DELETE CASCADE}, \texttt{ON DELETE SET NULL}, and \texttt{ON DELETE RESTRICT}. Why is \texttt{CASCADE} prohibited on financial ledgers?
\end{enumerate}

\subsection{Module C: Java Database Connectivity (JDBC) and MongoDB}
\begin{enumerate}[resume]
    \item What is JDBC? Illustrate the 4-tier architectural communication path from Java code to the physical RDBMS.
    \item Explain the four types of JDBC drivers. Why is the Type 4 Thin Driver considered the modern enterprise industry standard?
    \item Compare the \texttt{Statement}, \texttt{PreparedStatement}, and \texttt{CallableStatement} interfaces.
    \item How does \texttt{PreparedStatement} prevent SQL injection attacks? Explain using abstract syntax tree compilation mechanics.
    \item List the complete sequence of programmatic steps required to connect, query, and close database resources in Java.
    \item What is MongoDB? Explain its document-oriented data model and define BSON.
    \item Provide a 7-point comparative matrix between Relational Database Systems (RDBMS) and MongoDB.
    \item Write MongoDB commands to perform the four fundamental CRUD operations on a collection.
    \item Explain the MongoDB Aggregation Pipeline framework and detail the roles of \texttt{\$match}, \texttt{\$group}, \texttt{\$sort}, and \texttt{\$project}.
    \item Compare Embedded Subdocuments (denormalization) with Document References (normalization) in MongoDB schema design.
    \item Explain High Availability in MongoDB via Replica Sets. Describe the roles of Primary, Secondary, and the automatic failover election mechanism.
    \item Explain Horizontal Scalability in MongoDB via Sharding. Detail the functions of Shards, Config Servers, and \texttt{mongos} query routers.
\end{enumerate}

\subsection{Module D: Banking Database Architecture and Transactions}
\begin{enumerate}[resume]
    \item Design a relational schema for a core banking application including \texttt{Branches}, \texttt{Customers}, \texttt{Accounts}, \texttt{Account\_Holders}, and \texttt{Bank\_Transactions}.
    \item Explain how Primary Keys, Foreign Keys, and Check Constraints maintain financial integrity in a banking schema.
    \item Write an atomic SQL transaction script demonstrating a funds transfer of BDT 10,000 between two bank accounts.
    \item Map the four ACID properties (Atomicity, Consistency, Isolation, Durability) to specific mechanisms in a core banking funds transfer.
    \item Write SQL queries for a banking system to calculate total deposits per branch and identify customers who hold no bank accounts.
\end{enumerate}

\subsection{Module E: Big Data and Apache Hadoop Ecosystem}
\begin{enumerate}[resume]
    \item Define Big Data and detail the 3 Vs (Volume, Velocity, Variety) along with Veracity and Value.
    \item What is Apache Hadoop? Name and define its four core modules.
    \item Explain the master-worker architecture of HDFS. Contrast the responsibilities of the NameNode and DataNodes.
    \item Why does HDFS use large default block sizes (128 MB or 256 MB)? Explain using disk seek versus sequential transfer latency mathematics.
    \item Explain the HDFS Block Replication Mechanism and describe the 3-replica Rack Awareness Policy.
    \item What is the function of the Secondary NameNode? Explain why it is an architectural misconception to consider it a failover backup.
    \item Describe the YARN architecture. Detail the responsibilities of the ResourceManager, NodeManager, ApplicationMaster, and Containers.
    \item Explain the MapReduce distributed programming paradigm using the canonical Word Count algorithm as an operational case study.
    \item Compare Hadoop MapReduce with Apache Spark. Why is Spark significantly faster for iterative machine learning algorithms?
    \item Provide a clear architectural distinction between HDFS, Apache HBase, and Apache Hive.
    \item What are Apache Pig and Apache Oozie, and what roles do they perform in big data pipelines?
    \item Explain the fundamental architectural difference between Data Replication and Data Sharding in distributed data systems.
    \item Discuss the major engineering advantages and architectural limitations of the Apache Hadoop framework.
\end{enumerate}



% =============================================================================
% SECTION: PRACTICAL SQL EXAMINATION QUERY BANK (25 SOLVED QUERIES)
% =============================================================================
\section{Practical SQL Examination Query Bank: 25 Production Queries and Solutions}
\label{sec:dbms_practical_sql_25}

This practical query bank uses the standard four-table enterprise schema tested across government technical recruitment examinations:
\begin{itemize}
    \item \textbf{Employees}(\underline{Employee\_ID}, Employee\_Name, Department\_ID, Salary, Manager\_ID, Join\_Date)
    \item \textbf{Departments}(\underline{Department\_ID}, Department\_Name)
    \item \textbf{Projects}(\underline{Project\_ID}, Project\_Name, Department\_ID, Budget)
    \item \textbf{Employee\_Projects}(\underline{Employee\_ID}, \underline{Project\_ID}, Assigned\_Date)
\end{itemize}

\begin{enumerate}[label=\textbf{Q\arabic*.},leftmargin=1.2cm]
\item \textbf{Retrieve all employees in ascending order of salary:}
\begin{lstlisting}[language=SQL]
SELECT * FROM Employees ORDER BY Salary ASC;
\end{lstlisting}

\item \textbf{Find the three highest-paid employees:}
\begin{lstlisting}[language=SQL]
SELECT Employee_ID, Employee_Name, Salary 
FROM Employees 
ORDER BY Salary DESC, Employee_ID 
LIMIT 3;
\end{lstlisting}

\item \textbf{Find the second-highest distinct salary without vendor-specific offsets:}
\begin{lstlisting}[language=SQL]
SELECT MAX(Salary) AS Second_Highest_Salary 
FROM Employees 
WHERE Salary < (SELECT MAX(Salary) FROM Employees);
\end{lstlisting}

\item \textbf{Find employees earning more than the global company average salary:}
\begin{lstlisting}[language=SQL]
SELECT Employee_ID, Employee_Name, Salary 
FROM Employees 
WHERE Salary > (SELECT AVG(Salary) FROM Employees);
\end{lstlisting}

\item \textbf{Calculate the average salary for each department:}
\begin{lstlisting}[language=SQL]
SELECT Department_ID, AVG(Salary) AS Average_Salary 
FROM Employees 
GROUP BY Department_ID;
\end{lstlisting}

\item \textbf{Find departments having more than five employees:}
\begin{lstlisting}[language=SQL]
SELECT Department_ID, COUNT(*) AS Employee_Count 
FROM Employees 
GROUP BY Department_ID 
HAVING COUNT(*) > 5;
\end{lstlisting}

\item \textbf{Find departments without any assigned employees (LEFT JOIN):}
\begin{lstlisting}[language=SQL]
SELECT d.Department_ID, d.Department_Name 
FROM Departments d 
LEFT JOIN Employees e ON d.Department_ID = e.Department_ID 
WHERE e.Employee_ID IS NULL;
\end{lstlisting}

\item \textbf{Retrieve employee names alongside their department names (INNER JOIN):}
\begin{lstlisting}[language=SQL]
SELECT e.Employee_ID, e.Employee_Name, d.Department_Name 
FROM Employees e 
INNER JOIN Departments d ON e.Department_ID = d.Department_ID;
\end{lstlisting}

\item \textbf{Retrieve every department including those with zero employees:}
\begin{lstlisting}[language=SQL]
SELECT d.Department_ID, d.Department_Name, e.Employee_Name 
FROM Departments d 
LEFT JOIN Employees e ON d.Department_ID = e.Department_ID;
\end{lstlisting}

\item \textbf{Find employees with no assigned projects (NOT EXISTS):}
\begin{lstlisting}[language=SQL]
SELECT e.Employee_ID, e.Employee_Name 
FROM Employees e 
WHERE NOT EXISTS (
    SELECT 1 FROM Employee_Projects ep WHERE ep.Employee_ID = e.Employee_ID
);
\end{lstlisting}

\item \textbf{Find employees earning more than the average salary of their own department:}
\begin{lstlisting}[language=SQL]
SELECT e.Employee_ID, e.Employee_Name, e.Department_ID, e.Salary 
FROM Employees e 
WHERE e.Salary > (
    SELECT AVG(e2.Salary) 
    FROM Employees e2 
    WHERE e2.Department_ID = e.Department_ID
);
\end{lstlisting}

\item \textbf{Find the highest-paid employee(s) in each department using Window Functions:}
\begin{lstlisting}[language=SQL]
WITH Ranked_Salaries AS (
    SELECT Employee_ID, Employee_Name, Department_ID, Salary,
           RANK() OVER (PARTITION BY Department_ID ORDER BY Salary DESC) AS Rank_Num
    FROM Employees
)
SELECT * FROM Ranked_Salaries WHERE Rank_Num = 1;
\end{lstlisting}

\item \textbf{Assign unique sequential numbers to employees based on salary:}
\begin{lstlisting}[language=SQL]
SELECT Employee_ID, Employee_Name, Salary,
       ROW_NUMBER() OVER (ORDER BY Salary DESC, Employee_ID) AS Row_Num 
FROM Employees;
\end{lstlisting}

\item \textbf{Find employees and their direct manager names (SELF JOIN):}
\begin{lstlisting}[language=SQL]
SELECT e.Employee_Name AS Employee, m.Employee_Name AS Manager 
FROM Employees e 
LEFT JOIN Employees m ON e.Manager_ID = m.Employee_ID;
\end{lstlisting}

\item \textbf{Find projects with budgets exceeding the average project budget:}
\begin{lstlisting}[language=SQL]
SELECT Project_ID, Project_Name, Budget 
FROM Projects 
WHERE Budget > (SELECT AVG(Budget) FROM Projects);
\end{lstlisting}

\item \textbf{Calculate total project budget allocated per department:}
\begin{lstlisting}[language=SQL]
SELECT Department_ID, SUM(Budget) AS Total_Department_Budget 
FROM Projects 
GROUP BY Department_ID;
\end{lstlisting}

\item \textbf{Count the number of employees assigned to each project (including zero):}
\begin{lstlisting}[language=SQL]
SELECT p.Project_ID, p.Project_Name, COUNT(ep.Employee_ID) AS Total_Assigned_Staff 
FROM Projects p 
LEFT JOIN Employee_Projects ep ON p.Project_ID = ep.Project_ID 
GROUP BY p.Project_ID, p.Project_Name;
\end{lstlisting}

\item \textbf{Find employees who joined during the calendar year 2025 (SARGable query):}
\begin{lstlisting}[language=SQL]
SELECT * FROM Employees 
WHERE Join_Date >= DATE '2025-01-01' AND Join_Date < DATE '2026-01-01';
\end{lstlisting}

\item \textbf{Increase the salary of all employees in Department 10 by 10\%:}
\begin{lstlisting}[language=SQL]
UPDATE Employees 
SET Salary = Salary * 1.10 
WHERE Department_ID = 10;
\end{lstlisting}

\item \textbf{Delete all employees assigned to temporary test Department 99:}
\begin{lstlisting}[language=SQL]
DELETE FROM Employees 
WHERE Department_ID = 99;
\end{lstlisting}

\item \textbf{Create a secure database view displaying employee names and department names:}
\begin{lstlisting}[language=SQL]
CREATE VIEW Employee_Department_View AS 
SELECT e.Employee_ID, e.Employee_Name, d.Department_Name 
FROM Employees e 
INNER JOIN Departments d ON e.Department_ID = d.Department_ID;
\end{lstlisting}

\item \textbf{Create a secondary B-tree index to accelerate department-based employee filtering:}
\begin{lstlisting}[language=SQL]
CREATE INDEX idx_emp_dept ON Employees(Department_ID);
\end{lstlisting}

\item \textbf{Find employees whose salary exceeds their own direct manager's salary:}
\begin{lstlisting}[language=SQL]
SELECT e.Employee_Name, e.Salary, m.Employee_Name AS Manager_Name, m.Salary AS Manager_Salary 
FROM Employees e 
INNER JOIN Employees m ON e.Manager_ID = m.Employee_ID 
WHERE e.Salary > m.Salary;
\end{lstlisting}

\item \textbf{Find the single department having the highest average salary:}
\begin{lstlisting}[language=SQL]
SELECT Department_ID, AVG(Salary) AS Avg_Salary 
FROM Employees 
GROUP BY Department_ID 
ORDER BY Avg_Salary DESC 
FETCH FIRST 1 ROW ONLY;
\end{lstlisting}

\item \textbf{Find duplicate non-null email addresses in an employee database:}
\begin{lstlisting}[language=SQL]
SELECT Email, COUNT(*) AS Duplicate_Occurrences 
FROM Employees 
WHERE Email IS NOT NULL 
GROUP BY Email 
HAVING COUNT(*) > 1;
\end{lstlisting}
\end{enumerate}

% =============================================================================
% SECTION: NUMERICAL, SIZING & ARCHITECTURAL CALCULATIONS (PROBLEMS 1-10)
% =============================================================================
\section{Numerical, Sizing, and Architectural Calculations: Solved Problems}
\label{sec:dbms_numerical_problems}

\begin{enumerate}[label=\textbf{Problem \arabic*.},leftmargin=1.5cm]
\item \textbf{CROSS JOIN Cardinality Calculation:}
\textit{Question:} Relation $A$ contains 5 records; Relation $B$ contains 4 records. What is the cardinality of their Cartesian product $A \times B$? \\
\textbf{Solution:}
\[
|A \times B| = |A| \times |B| = 5 \times 4 = \mathbf{20\text{ rows}}
\]

\item \textbf{Group Formation Upper Bound:}
\textit{Question:} An \texttt{Employees} table has 100 rows containing exactly 5 distinct \texttt{Department\_ID} keys. What is the maximum number of output tuples produced by \texttt{GROUP BY Department\_ID}? \\
\textbf{Solution:} The grouping operator partitions the input multiset into distinct equivalence classes. Because there are 5 unique keys, the query produces exactly $\mathbf{5\text{ groups}}$ (tuples).

\item \textbf{Arithmetic Mean Calculation with NULL Values:}
\textit{Question:} A numeric attribute contains the values: $\{30000, 40000, \text{NULL}, 50000, 60000\}$. Calculate the exact value of \texttt{AVG(Salary)}. \\
\textbf{Solution:} By ANSI SQL rules, all aggregate functions except \texttt{COUNT(*)} ignore nulls:
\[
\text{AVG} = \frac{\sum \text{non-null values}}{\text{COUNT(Salary)}} = \frac{30000 + 40000 + 50000 + 60000}{4} = \frac{180000}{4} = \mathbf{45,000}
\]

\item \textbf{Evaluating COUNT(*) versus COUNT(column) with NULLs:}
\textit{Question:} For the dataset $\{30000, 40000, \text{NULL}, 50000, 60000\}$, evaluate \texttt{COUNT(*)} and \texttt{COUNT(Salary)}. \\
\textbf{Solution:}
\[
\text{COUNT}(*) = \mathbf{5\text{ (total rows)}}, \quad \text{COUNT}(\text{Salary}) = \mathbf{4\text{ (non-null entries)}}
\]

\item \textbf{Secondary Index Physical Storage Estimation:}
\textit{Question:} A table has 1,000,000 records. A secondary B-tree index entry occupies an average of 40 bytes per entry. Estimate the raw storage required for the index leaf entries. \\
\textbf{Solution:}
\[
\text{Raw Storage} = 1,000,000 \times 40\text{ bytes} = 40,000,000\text{ bytes} = \mathbf{40\text{ MB}}
\]

\item \textbf{Query Latency and Throughput Computation:}
\textit{Question:} A critical OLTP query executes 200 times and consumes a total cumulative server time of 10 seconds. Calculate the average query execution latency. \\
\textbf{Solution:}
\[
\text{Average Latency} = \frac{10\text{ seconds}}{200} = 0.05\text{ seconds} = \mathbf{50\text{ milliseconds}}
\]

\item \textbf{Analytic Window Ranking Evaluation under Ties:}
\textit{Question:} Given descending salary values $\{80000, 70000, 70000, 50000\}$, determine the exact values computed by \texttt{ROW\_NUMBER()}, \texttt{RANK()}, and \texttt{DENSE\_RANK()}. \\
\textbf{Solution:}
\begin{center}
\begin{tabular}{l c c c}
\toprule
\textbf{Salary} & \textbf{ROW\_NUMBER()} & \textbf{RANK()} & \textbf{DENSE\_RANK()} \\
\midrule
80,000 & 1 & 1 & 1 \\
70,000 & 2 & 2 & 2 \\
70,000 & 3 & 2 & 2 \\
50,000 & 4 & 4 (gap) & 3 (no gap) \\
\bottomrule
\end{tabular}
\end{center}

\item \textbf{Referential Cascade Cascading Destruction Impact:}
\textit{Question:} A department managing 20 employees is deleted. The foreign key on \texttt{Employees(Department\_ID)} specifies \texttt{ON DELETE CASCADE}. What is the operational effect? \\
\textbf{Solution:} The RDBMS automatically deletes all 20 referencing employee records from the \texttt{Employees} table.

\item \textbf{LEFT OUTER JOIN Cardinality Determinism:}
\textit{Question:} Relation $A$ contains 10 rows; Relation $B$ contains 8 rows. Exactly 6 rows match on the join key (with unique 1:1 join keys). How many rows are returned by $A \text{ } {}_{{}^{\text{L}}}\!\bowtie B$? \\
\textbf{Solution:} A LEFT JOIN preserves every row from the left relation ($A$). The 6 matching rows are paired with $B$ attributes; the 4 non-matching rows are padded with NULL. Total rows = $\mathbf{10\text{ rows}}$.

\item \textbf{Distributed Cluster Replication Physical Sizing:}
\textit{Question:} A distributed database stores 2 TB of logical table data with a replication factor of 3. What is the total raw disk storage consumed for replicas? \\
\textbf{Solution:}
\[
\text{Total Physical Storage} = 2\text{ TB} \times 3 = \mathbf{6\text{ TB}}
\]
\end{enumerate}

% =============================================================================
% SECTION: 25 MODEL SHORT-ANSWER EXAMINATION QUESTIONS & MODEL ANSWERS
% =============================================================================
\section{Theoretical Examination Bank: 25 Core Questions and Model Answers}
\label{sec:dbms_model_answers_25}

\subsection{Module A: Database Theory and Relational Algebra}
\begin{enumerate}[label=\textbf{TA\arabic*.},leftmargin=1.2cm]
\item \textbf{Differentiate DBMS and RDBMS.} \\
\textbf{Answer:} A DBMS is general software for managing stored data without mandatory relational rules. An RDBMS is based on E. F. Codd's relational model, storing data in normalized tabular relations, enforcing primary/foreign keys, and ensuring mathematical integrity.

\item \textbf{Define Relational Algebra.} \\
\textbf{Answer:} Relational algebra is a formal procedural query language where operators ($\sigma, \pi, \cup, -, \times, \bowtie$) take relations as input and produce a resultant relation.

\item \textbf{Differentiate Relational Algebra and Relational Calculus.} \\
\textbf{Answer:} Relational algebra is \textit{procedural} (prescribes \textit{how} to retrieve data); relational calculus is \textit{declarative} (prescribes \textit{what} data satisfies given logical predicates).

\item \textbf{Define a Database View.} \\
\textbf{Answer:} A virtual table defined by an SQL query stored in the system catalog. It contains no physical data of its own and simplifies queries while enforcing column/row access control.

\item \textbf{What is a Materialized View?} \\
\textbf{Answer:} A materialized view physically stores the query result on disk pages. It provides fast reads for heavy OLAP aggregations but requires refreshing when base data changes.

\item \textbf{What is a Common Table Expression (CTE)?} \\
\textbf{Answer:} A named temporary result set defined using the \texttt{WITH} clause within a single SQL statement, enhancing query readability and enabling recursive traversals.

\item \textbf{What is an SQL Window Function?} \\
\textbf{Answer:} An analytic function that computes metrics across related rows (\texttt{OVER}) while preserving the individual row identities of the input relation.

\item \textbf{Differentiate RANK() and DENSE\_RANK().} \\
\textbf{Answer:} Both assign equal rank to tied values; \texttt{RANK()} skips subsequent integers based on duplicate counts, whereas \texttt{DENSE\_RANK()} never skips integers.

\item \textbf{What is a Stored Procedure?} \\
\textbf{Answer:} A precompiled, named collection of procedural SQL statements stored within the database server, executed explicitly via \texttt{CALL} to perform business operations.

\item \textbf{What is a Database Trigger?} \\
\textbf{Answer:} An active database object that executes automatically in response to specified DML events (\texttt{INSERT}, \texttt{UPDATE}, \texttt{DELETE}) for auditing or validation.

\item \textbf{Differentiate Stored Procedure and Trigger.} \\
\textbf{Answer:} A stored procedure is invoked explicitly by an application or user; a trigger fires implicitly and automatically when an underlying table event occurs.

\item \textbf{What is a Database Cursor?} \\
\textbf{Answer:} A procedural pointer allowing applications to fetch and process query result tuples row-by-row rather than as a set.

\item \textbf{Explain Three-Valued Logic in SQL.} \\
\textbf{Answer:} SQL logic accommodates \texttt{TRUE}, \texttt{FALSE}, and \texttt{UNKNOWN}. Evaluating comparisons against \texttt{NULL} yields \texttt{UNKNOWN}; \texttt{WHERE} clauses filter for strictly \texttt{TRUE}.

\item \textbf{Define Database Index.} \\
\textbf{Answer:} An auxiliary search structure (typically B+ Tree) that maps search keys to physical storage addresses, accelerating query retrieval from $O(N)$ to $O(\log N)$.

\item \textbf{What is Query Optimization?} \\
\textbf{Answer:} The process by which the Cost-Based Optimizer analyzes alternative execution plans and selects the lowest-cost strategy based on I/O, CPU, and data statistics.
\end{enumerate}

\subsection{Module B: Advanced Database Management and Distributed Systems}
\begin{enumerate}[label=\textbf{TB\arabic*.},leftmargin=1.2cm]
\item \textbf{What is Multi-Version Concurrency Control (MVCC)?} \\
\textbf{Answer:} A concurrency technique that maintains multiple physical versions of rows, allowing readers to view consistent snapshots without acquiring shared locks or blocking writers.

\item \textbf{Compare Pessimistic and Optimistic Concurrency Control.} \\
\textbf{Answer:} Pessimistic locking acquires locks upfront assuming collisions are frequent; optimistic concurrency control proceeds without locks and checks for version collisions at commit.

\item \textbf{Explain Two-Phase Commit (2PC).} \\
\textbf{Answer:} A distributed consensus protocol ensuring atomic commits across multiple database nodes via a \textit{Prepare Phase} (voting) and a \textit{Commit/Abort Phase}.

\item \textbf{What is Database Partitioning?} \\
\textbf{Answer:} Dividing a large table into smaller physical partitions (horizontal by rows, vertical by columns) to optimize I/O performance and facilitate data archiving.

\item \textbf{Explain the CAP Theorem.} \\
\textbf{Answer:} States that during an unavoidable network partition ($P$), a distributed data system must choose between guaranteeing full Consistency ($C$) or full Availability ($A$).

\item \textbf{Differentiate OLTP and OLAP.} \\
\textbf{Answer:} OLTP handles high-volume, real-time, low-latency operational transactions on normalized schemas; OLAP executes complex read-heavy historical analysis on dimensional schemas.

\item \textbf{Define Data Warehouse.} \\
\textbf{Answer:} A centralized, subject-oriented, integrated, time-variant, and non-volatile database repository designed for organizational business intelligence.

\item \textbf{Differentiate ETL and ELT.} \\
\textbf{Answer:} ETL transforms raw data on an external processing server before loading into the warehouse; ELT loads raw data directly and transforms it in-place using warehouse compute power.

\item \textbf{What is a Star Schema?} \\
\textbf{Answer:} A dimensional data warehouse structure consisting of a central numerical fact table directly referencing surrounding denormalized dimension tables.

\item \textbf{What is Point-in-Time Recovery (PITR)?} \\
\textbf{Answer:} A disaster recovery method restoring a database to a specific historical second by applying continuous Write-Ahead Log redo records over a base physical backup.
\end{enumerate}

% =============================================================================
% SECTION: 21 COMMON DBMS EXAMINATION TRAPS & MISCONCEPTIONS
% =============================================================================
\section{Common DBMS Examination Traps and Diagnostic Misconceptions}
\label{sec:dbms_exam_traps}

Technical examination questions often contain subtle traps designed to test deep conceptual understanding. Review these 21 definitive diagnostic principles:

\begin{table}[htbp]
\centering
\scriptsize
\caption{The 21 Definitive Examination Traps and Corrections}
\label{tab:21_traps}
\begin{tabularx}{\textwidth}{l p{4.5cm} X}
\toprule
\textbf{\#} & \textbf{Common Examination Misconception} & \textbf{Rigorous Technical Reality} \\
\midrule
1 & A view is a stored duplicate of table data. & \textbf{False.} A standard view is purely virtual; it stores only the SQL text definition in the catalog. \\
2 & A materialized view is a virtual object. & \textbf{False.} A materialized view physically materializes and stores its result tuples on disk. \\
3 & \texttt{GROUP BY} retains every individual row. & \textbf{False.} \texttt{GROUP BY} collapses all rows sharing identical grouping keys into a single summary tuple. \\
4 & Window functions collapse rows like \texttt{GROUP BY}. & \textbf{False.} Window functions calculate metrics across a partition while preserving all individual input rows. \\
5 & \texttt{COUNT(col)} counts \texttt{NULL} entries. & \textbf{False.} \texttt{COUNT(col)} counts only rows where the specified attribute is strictly non-null. \\
6 & \texttt{COUNT(*)} ignores \texttt{NULL} rows. & \textbf{False.} \texttt{COUNT(*)} counts every candidate tuple in the relation regardless of null contents. \\
7 & \texttt{NULL} is equivalent to numeric zero ($0$). & \textbf{False.} \texttt{NULL} represents the complete absence of a value; zero is a concrete integer. \\
8 & SQL boolean expressions evaluate only to T or F. & \textbf{False.} SQL uses Three-Valued Logic ($\{\text{TRUE}, \text{FALSE}, \text{UNKNOWN}\}$). \\
9 & Triggers are executed explicitly via \texttt{CALL}. & \textbf{False.} Triggers fire implicitly and automatically when an underlying DML/DDL event occurs. \\
10 & JDBC is a standalone relational database engine. & \textbf{False.} JDBC is a Java SE API specification for connecting to external database engines. \\
11 & \texttt{PreparedStatement} prevents SQL injection. & \textbf{True.} Precompiled AST parameter binding guarantees input data is never executed as syntax. \\
12 & Creating an index always accelerates queries. & \textbf{False.} Indexes incur write overhead, consume storage, and the CBO may reject them for low selectivity. \\
13 & MVCC is an ANSI SQL transaction isolation level. & \textbf{False.} MVCC is an internal concurrency implementation mechanism, not an isolation standard. \\
14 & 2PC and 2PL are identical protocols. & \textbf{False.} 2PC is distributed commit consensus; 2PL is local concurrency lock management. \\
15 & Partitioning is another name for normalization. & \textbf{False.} Normalization eliminates logical update anomalies; partitioning divides physical storage. \\
16 & Database replication eliminates the need for backups. & \textbf{False.} Replication immediately copies accidental drops or corrupted writes to all replicas. \\
17 & CAP allows systems to choose any 2 of 3 freely. & \textbf{False.} Network partitions are unavoidable; systems must choose between C and A during partitions. \\
18 & OLTP is designed for historical trend reporting. & \textbf{False.} OLTP executes high-frequency, low-latency operational updates; OLAP handles reporting. \\
19 & ETL loads data before performing transformations. & \textbf{False.} ETL transforms data before loading; ELT loads raw data before transforming in-place. \\
20 & Secondary NameNode is a high-availability backup. & \textbf{False.} Secondary NameNode only checkpoints logs; Standby NameNode provides live HA failover. \\
21 & \texttt{NOT IN} is always safe with subqueries. & \textbf{False.} If the subquery returns even a single \texttt{NULL}, \texttt{NOT IN} evaluates to UNKNOWN and returns 0 rows! \\
\bottomrule
\end{tabularx}
\end{table}

% =============================================================================
% SECTION: FINAL COMPREHENSIVE 48-TOPIC MASTER REVISION CHECKLIST
% =============================================================================
\section{Final Comprehensive 48-Topic Master Examination Revision Checklist}
\label{sec:dbms_master_checklist}

This comprehensive audit checklist summarizes the forty-eight essential knowledge modules of the complete Database Management Systems and SQL curriculum. Candidates should verify full conceptual mastery across all six modules before attending competitive written tests and viva voce boards:

\begin{figure}[htbp]
\centering
\small
\begin{tabularx}{\textwidth}{l X}
\toprule
\textbf{Curricular Module} & \textbf{Required Core Competencies and Knowledge Areas} \\
\midrule
\textbf{1. Relational Foundations} & $\square$ Database vs. DBMS vs. RDBMS \quad $\square$ Schemas vs. Instances vs. Metadata \newline
$\square$ Degree and Cardinality \quad $\square$ Key Hierarchy ($\text{PK} \subseteq \text{AK} \subset \text{CK} \subseteq \text{SK}$) \newline
$\square$ Cardinality Ratios (1:1, 1:N, M:N) \quad $\square$ ANSI/SPARC 3-Schema Architecture \newline
$\square$ Logical and Physical Data Independence \quad $\square$ Relational Algebra ($\sigma, \pi, \cup, -, \times, \bowtie$) \newline
$\square$ Tuple and Domain Relational Calculus \quad $\square$ Codd's Relational Completeness \\
\midrule
\textbf{2. SQL Query Engine} & $\square$ SQL Command Taxonomy (DDL, DML, DQL, DCL, TCL) \newline
$\square$ Data Types: \texttt{DECIMAL} vs. \texttt{FLOAT}, \texttt{CHAR} vs. \texttt{VARCHAR} \newline
$\square$ Logical Query Processing Sequence (FROM $\to$ WHERE $\to$ GROUP BY $\to$ HAVING $\to$ SELECT) \newline
$\square$ Subqueries: Scalar, Correlated, EXISTS vs. IN, ANY/ALL \newline
$\square$ Aggregate Functions and NULL Semantics \quad $\square$ WHERE vs. HAVING Paradigm \newline
$\square$ Relational JOINs (INNER, OUTER, CROSS, SELF) \quad $\square$ UNION vs. UNION ALL \newline
$\square$ Views, Updatability, and WITH CHECK OPTION \quad $\square$ Materialized Views \newline
$\square$ Common Table Expressions (CTEs) and Recursive Graph Traversal \newline
$\square$ Window Functions: \texttt{ROW\_NUMBER}, \texttt{RANK}, \texttt{DENSE\_RANK}, \texttt{LAG}/\texttt{LEAD}, Frames \newline
$\square$ Procedural SQL: Procedures, Functions, Triggers, and Cursors \newline
$\square$ Three-Valued Logic (3VL), UNKNOWN, \texttt{COALESCE}, and \texttt{NULLIF} \\
\midrule
\textbf{3. Integrity \& Transactions} & $\square$ Entity, Referential, Domain, and User-Defined Integrity \newline
$\square$ Foreign Key Cascade Actions: CASCADE, SET NULL, RESTRICT \newline
$\square$ ACID Properties (Atomicity, Consistency, Isolation, Durability) \newline
$\square$ Transaction States (Active, Committed, Failed, Aborted) \newline
$\square$ Concurrency Anomalies: Lost Update, Dirty Read, Non-Repeatable Read, Phantom \newline
$\square$ ANSI SQL Isolation Levels (Read Uncommitted to Serializable) \newline
$\square$ Shared and Exclusive Locks \quad $\square$ Strict Two-Phase Locking (Strict 2PL) \newline
$\square$ Deadlocks, Coffman Conditions, Wait-For Graph, and Prevention \newline
$\square$ Multi-Version Concurrency Control (MVCC) \quad $\square$ Optimistic vs. Pessimistic Locking \newline
$\square$ Write-Ahead Logging (WAL) Protocol \quad $\square$ Checkpoints, Redo and Undo Recovery \\
\midrule
\textbf{4. Normalization \& Design} & $\square$ Insertion, Update, and Deletion Anomalies \newline
$\square$ Functional Dependencies (Full, Partial, Transitive) \quad $\square$ Armstrong's Axioms \newline
$\square$ Normal Forms: 1NF, 2NF, 3NF, BCNF, 4NF, 5NF \newline
$\square$ Binary Lossless Join Decomposition Condition \quad $\square$ Dependency Preservation \newline
$\square$ Controlled Denormalization Trade-offs \\
\midrule
\textbf{5. Systems \& Big Data} & $\square$ Database Storage: 8 KB Pages, Slotted Pages, Buffer Pool, Dirty Pages \newline
$\square$ Index Architectures: B+ Trees, Hash Indexes, Clustered vs. Secondary, Covering \newline
$\square$ Query Optimization, Cost-Based Optimizer (CBO), EXPLAIN ANALYZE, SARGability \newline
$\square$ Database Partitioning: Horizontal (Range, List, Hash) vs. Vertical \newline
$\square$ Distributed Databases, Synchronous vs. Asynchronous Replication \newline
$\square$ CAP Theorem (Consistency, Availability, Partition Tolerance) \newline
$\square$ Two-Phase Commit (2PC) Protocol \quad $\square$ JDBC Driver Types 1--4 and Security \newline
$\square$ Apache MongoDB: BSON, CRUD, Aggregation Pipelines, Replica Sets, Sharding \newline
$\square$ Apache Hadoop: Big Data 3 Vs, HDFS Architecture, NameNode, DataNodes, Sizing \newline
$\square$ YARN Resource Management, MapReduce Execution Engine, Apache Spark \\
\midrule
\textbf{6. Data Engineering \& Admin} & $\square$ OLTP vs. OLAP Systems \quad $\square$ Data Warehouse Inman Characteristics \newline
$\square$ Data Marts \quad $\square$ ETL vs. ELT Pipelines \newline
$\square$ Dimensional Modeling: Star Schema vs. Snowflake Schema, Fact vs. Dimension Tables \newline
$\square$ Role-Based Access Control (RBAC), GRANT and REVOKE Security \newline
$\square$ Database Auditing \quad $\square$ Point-in-Time Recovery (PITR) Log Replay \\
\bottomrule
\end{tabularx}
\end{figure}
